
%%%%%%%%%%%%%%%%%%%%%%%%%%%%%%%% SCRIPT FOR E-RARA AND MDZ FILES     %%%%%%%%%%%%%%%%%%%%%%%%%%%%%%%%%%%%%%%%%%%%%%%%
%%%%%%%%%%%%%%%%%%%%%%%%%%% fini le 30.04.2025 par F. GOY            %%%%%%%%%%%%%%%%%%%%%%%%%%%%%%%%%%%%%%%%%%%%%%%%
% !TeX TS-program = lualatex
\documentclass{article}
\usepackage[T1]{fontenc}
\usepackage{microtype}
\usepackage[pdfusetitle,hidelinks]{hyperref}

\usepackage{polyglossia}
\setmainlanguage{english}
\setotherlanguages{latin,greek}
\usepackage[series={},nocritical,noend,noeledsec,nofamiliar,noledgroup]{reledmac}
\usepackage{reledpar}

\usepackage{fontspec}
\setmainfont{TeX Gyre Termes}

\usepackage{sectsty}
\usepackage{xcolor}

\usepackage{fancyhdr}
\pagestyle{fancy}
\fancyhf{}
\fancyhead[LE,RO]{\nouppercase{\leftmark}}  
\cfoot{\thepage}
\renewcommand{\headrulewidth}{0.4pt}

% Redefine \section to remove numbering
\usepackage{titlesec}
\titleformat{\section}[block]{\normalfont\scshape\color{gray}}{}{0pt}{} % no number in heading
\titleformat{\subsection}[hang]{\normalfont}{}{0pt}{} % also remove subsection number
\titleformat{\subsubsection}[hang]{\normalfont\footnotesize\color{black}}{}{0pt}{}

% Modify how section marks are stored to exclude numbers
\makeatletter
\renewcommand{\sectionmark}[1]{%
	\markboth{#1}{}} % Only store the section title, without number
\renewcommand{\subsectionmark}[1]{%
	\markright{#1}} % Only store the subsection title, without number
\renewcommand{\numberline}[1]{} % Hide the section number in TOC
\makeatother

\begin{document}

\date{}
        \title{Commentarii in utranque Pauli epistolam ad Timotheum: [Calvin Jean], [1548]}
\maketitle
\tableofcontents
\clearpage
\begin{pages} 
\beginnumbering
        IOANNIS CALVINI COMMENTA. RII, IN VTRANQVE Pauli Epistolam ad Timotheum. GENEVAE, PER IOANNEA GERARDVRS 1548 ILLVSTRISS.ET VERE CHRISTIANO PRINCIPI, D. Etuardo Sommerseti Duci. Comiti Hertphordiae, et c. Angliae et Hiberniae protectori, tutorique. Regio: Ioannes Caluinus, S. \pstart \huge\textbf{C}\normalsize V M praeclari isti sermones, qui tuin de aliis tuis virtutibus plane Heroicis, tum vero de singulari tua pietate habentur (Illustriss. Princeps) tantum vbique tui amorem in bonorum omnium animis, quibus alioqui de facie ignotus es, excitent: necesse est incredibili te charitatis ac reuerent iae affectu prosequantur, quicunque in Regno Angliae recte sapiunt: quibus non modo palam cernere, quae alii audientes mirantur bona, datum est, sed fructum inde percipere, quantus a praestantissimo gubernatore prouenire, cum ad totum populi corpus, tum ad vnuinquodque membrum solet. Neque enim quae constanti faina celebrantur tuae laudes, vanitatis suspectae esse possunt, quasi ab adulatoribus profectae, cum in rebus tuis gestis illustre extet earum testimonium. In priuato et mediocris fortunae pupillo, difficilis est administratio tutelae. Tu vero non Regis modo, sed amplissimi regni tutelam tibi iniunctam ea dexteritate ac prudentia sustines, vt successum omnes mirentur. Ac ne inter leges modo et in tranquillo reip.statu niteret tua vir  \pend\pstart tus, eam in bello quoque spectandam proposuit Deus, quod pari foelicitate ac fortitudine hactenus auspiciis tuis gestum est. Neque tamen interea tot tantae que difficultates, quas reputare cuiuis promptum est, tu expertus es, impedimento tibi fuerunt, quominus instaurandae religionis curam imprimis susciperes. Cogitat io sane non magis principe viro digna, quam publico Regni bono vtilis. Nam haec demum stabilis est regnorumn foelicitas, fidelis que custodia, si is in quo fundata sunt, et per quem seruantur, Filius Dei illis praesideat. Itaque melius rem Anglicanam stabilire non poteras, quam profligatis idolis, purum ibi Dei cultum erigendo. Nam id fieri non potest, quin restituatur genuina pietatis doctrina, quae sacrilega Antichristi Romani tyrannide, nimis diu oppressa iacuit. Quod est, Christum in suo tribunali collocare. Atque hoc eximium alioqui facinus eo maiori laude propterea dignum est, quo pauciores hodie inter eos qui gubernacula tenent, reperiuntur, qui princ ipatus sui insignia spirituali Christi sceptro submittant. Praeclare igitur serenissimo Regi consultum est, cuius pueritiae talis ex sua cognatione obtigit moderator. Nam tametsi indolem eius imprimis generosam omnes praedicant: tamen et formandis ad virilem constantiam eius moribus, ex constituendae Anglicanae Ecclesiae, dum aetas ipsum non patitur has obire vices, ta\pend\pstart lis artifex erat plusquam necessarius. Nec dubito, quin ipse iam nunc agnoscat, diuinitus te sibi datum esse, ex cuius manu res ex voto compositas pau lopost reciperet. Me vero nec tanta locorum distantia, nec conditionis meae humilitas impedire potuit, quominus praestantissimum hunc propagandae Christi gloriae affectum tibi gratularer. Et sane, cum me vnum ex iis esse voluerit Deus, quorum opera, laboreque sinceriorem Euangelii doctrinam hodie mundo restituit: cur non te quoque sin gula ri Dei beneficio, eiusdem doctrinae patronum ac vindicem ordinatum, qualicunque locorum interuallo sim abs te dissitus, qua possum obseruant ia prosequar? Et quoniam non aliud huius rei testimonium reddere licebat, saltem arrae loco, hos in duas Pauli Epistolas meos commentarios tibi offerendos esse putaui. Nec vero fortuitum munus, quod offerrem, arripui: sed iudicio elegi, quod videbatur optime conuenire. Hic Timotheum suum admonet PauIus, quali doctrinae genere aedificanda sit Ecclesia Dei, quibus et vitiis et hostibus resistend um, quam multa deuoranda taedia:hortatur ne vllis difficultatibus caedat: vt omnia discrimina animi fortitudine superet: vt improborum proteruiam authoritate coerceat, ne quid hominum gratiae ambit iose largiatur. Denique in his duabus Epistolis, quasi in viua tabula depictum habemus verum Ecclesiae regi\pend\pstart men. Tu, cum instaurandae Ecclesiae Anglicanae, quam scelerata Papatus impietas, vt alias fere omnes Christiani orbis, misere corruperat, strenuain operam impendas sub Regis tui auspiciis: et multos ad cam rem Timotheos in opere habeas: sancta certe vestra studia melius componere non potestis, quam ad hoc Pauli quasi praescriptum exigendo. Nam neque hic quidquam habetur, quod non sit temporibus nostris valde aptum: et nihil ferme rursum in Ecclesiae architectura necessarium est, quod non hinc etiam haurire liceat. Laborautem meus mediocre saltem adiumetum, vt spero, afferet. Quod expe rient ia tamen cognosci malo, quam verbis iactare. Ti bi porro (Nobiliss. Princeps) si fuerit probatus: erit cur mihi plurimùm gratuler. Pro singulari quidem tua humanitate quin officium hoc meum boni consulturus sis, non dubito. Dominus, in cuius manu sunt fines terrae, statuin Regni Anglicani diu incolumem florentemque tueatur: Regem serenissimum tum Spiritu principali ornet, tum bonis omnibus cumulet, tibi que in praeclaro isto cursu det foeliciter pergere: vt per te Nomen eius magis ac magis propagetur.  \pend\pstart Geneuae 8. calendas Augusti.  \pend\pstart 548.  \pend
\phantomsection
\addcontentsline{toc}{subsection}{\textit{IN PRIOREM PAVLI EPISTO}}
\subsection*{\textit{IN PRIOREM PAVLI EPISTO}}
\phantomsection
\addcontentsline{toc}{subsection}{\textit{lam ad Timotheum, argumentum. }}
\subsection*{\textit{lam ad Timotheum, argumentum. }}\pstart \huge\textbf{H}\normalsize ANC Epistolam aliorum magis, quam Timothei cau sa scriptam fuisse, iudico: et mihi assentientur, qui diligenter omnia expenderint. Non equidem nego, quin eius quoque docendi et monendi rationem Paulus habuerit. Sed multa hic contineri dico, quae superuacuum fuisset scribere, si cum solo Timotheo habuisset negocium. Iuuenis erat, nondum ea authoritate instructus, quae sufficeret ad cohibendos homines proteruos, qui contra insurgerent. Apparet autem ex Pauli verbis, fuisse tunc nonnullos, qui cum ostentatione pollerent, non libenter cuiuis erant cessuri: qui tanta deinde ambitione flagrabant, vt nul lum turbandae Ecclesiae finem facturi essent, nisi quispiam maior Timotheo se interponeret. Apparet similiter, multa Ephesi fuisse consticuenda, quae Pauli calculum et nomen desiderarent. Ergo cum Timotheum multis de rebus monere haberet in animo, sub eius quoque nomine alios, eadem opera, monere voluit. Ac primo quidem capite, aduersus quosdam ambitiosos armat, qui vanis quaestionibus agitan dis se venditabant. Eos porro Iudaeos fuisse, colligi potest, qui cum legis zelum simularent, negleda aedificatione, tantùm curiosas disputationes ingerebant. Haec autem profanatio est legis Dei minime ferenda, ex ea nihil excerpere quod prosit: sed tantùm garriendi materiam captare, eiûsque praetextu ad onerandam friuolis nugis Ecclesiam, abuti. Cuiusmodi scilicet corruptela pluribus quam oportuit, s-culis regnauit in Papatu. Quid enim aliud fuit scholastica theologis quam immensumu inanium et nihili speculationum chaost Et hodie extant plerique, qui vt acumen suum ostentent, in tractando Dei verbo, non secus ac in profana philosophia, ludere fibi permittunt. In eiusmodi vitio coarguendo authorem se Timotheo profitetur Paulus: et quid ex lege praecipue discendum sit, admonet: vt appareat legis esse corruptores, qui ea secus vtuntur. Porro, ne contemptui sit sua authoritas, postquam indignitatem suam confessus est. simul qualis esse coeperit per Dei gratiam, honorifice praedicat. Tandem caput concludit graui denunciatione, qua et Timotheum in sana doctrina, et bona conscientia confirmat: et alios terret, proposito Hymenei et Alexandri exemplo. Secundo capite, iubet omnes in precibus publicis commendari Deo et nominatim principes ac magistratus vbi obiter etiam attingit, quid vtilitatis percipiat mundus, ex ciuili gubernatione. Cur autem pro omnibus concipi debeant preces, ratio additur, quia Deus Euangelium omnibus offerendo et Christum mediatorem, omnes se velle saluos fieri demonstret. Quod  \pend\pstart etiam confirmat a suo Apostolatu, qui Gentibus proprie destinatus erat. Inde omnes, cuiuscunque sint loci aut regionis, ad Deum orandum inuitat: et hac occasione, qualem modestiam et subiectionem in sacro coetu praeseferre debeant mulieres, disserit. Capite tertio, vbi indicauit, quam praeclarum sit munus episcopatus, verum Episcopum depingit, ac dotes enumerat, quae in eo requiruntur. Postea quales esse debeant Diaconi et vxores tam horum quam illorum, praescribit. Ac quo plus diligentiae ac religionis in omnibus obseruandis adhibeat Timotheus, reducit illi in memoriam, quid sit in Ecclesiae administratione, quae Domus Dei est, ac columna veritatis, versari. Tandem totius coelestis doctrinae praecipuum caput et quasi cardinem de Filio Dei in carne manifestato, commemorat: prae quo reliqua omnia leuia ducenda sint. Quod ad sequentia pertinet. Nam initio capitis quarti tam falsas doctrinas de coniugii et ciborum prohibitione, quam ineptas fabulas cum hoc capite non congruentes, seuere damnat. Et deinde subiicit, eos sibi et piis omnibus qui caput hoc tenent, solos esse aduersarios, qui ferre nequeunt, vt spes collocetur in Deum viuum. In fine capitis rursus noua exhortatione Timotheum confirmat. Quinto capite, vbi modestiam in reprehensionibus et senitatem illi commendauit, de viduis disputat, quae tunc in Ecclesiae ministeriuni cooptabantur. Ac vetat omnes promiscue recipi: sed que tota vita probatae ad annum sexagesimum peruenerint, et nullum habeant onus domesticum. Inde ad presbyteros transit, et qualiter tractandi sint, tum in victu, tum in diseiplina, admonet. Hanc doctrinam graui contestatione quasi sancit: et iterum prohibet, ne quos temere in ordinem presbyterii recipiat. Hortatur vt leuandae valetudinis causa vinum bibat aquae loco. In fine capitis hortatur, vt in peccatis occultis iudicium suspendat. Sexto capite, de seruorum officio praecipit: et hac occasione in falsos doctores inuehitur, qui de ociosis speculationibus litigando. quaestui magis quam aedificationi intenti sunt: et quam exitiosa sit pestis auaritia, docet. Inde iterum ad contestationem priori similem reuoluitur, ne frustra Timotheo haec praeceperit. Sub finem. postquam obiter de diuitibus meminit, iterum prohibet Timorheo, ne se inanibus doctrinis impediat. Vulgatae Graecorum inscriptioni, quae habet hanc Epistolam fuisse Laodicea missam, non assentior. Nam cum Paulus ex vinculis scribens Colossensibus testetur, Laodicenses nunquam sibi fuisse visos: coguntur, qui in illa sunt opinione, quam repudio, duplicem facere Laodiceam in Asia minori: quum alteram non celebrent scriptores. Deinde Paulus Timotheum Ephesi reliquerat, vt declarant cius verba, cum iret in Macedoniam. Hanc Epistolam, aut scripsit ex itinere, priusquam illuc  \pend\pstart appulisset, aut iam profectione confecta, reuersus. Atqui longius a Macedonia distare Laodiceam constat, quam Ephesum ipsam. Neque rursus verisimile est, ipsum in reditu, praeterita Epheso, Laodiceani se contulisse. Praesertim cum illuc eum multae causae vocarent. Itaque magis existimo aliunde eum scripsisse. Quanquam hoc ipsum non est tanti, vt velim cum dissentientibus contendere. Fruatur quisque suo iudicio. Tantùm indico, quid magis iit, meo iudicio, probabile.  \pend
\phantomsection
\addcontentsline{toc}{subsection}{\textit{IOANNIS CALVINI }}
\subsection*{\textit{IOANNIS CALVINI }}
\phantomsection
\addcontentsline{toc}{subsection}{\textit{COMMENTARII IN PRIOrem Pauli Epistolam ad Timotheum. }}
\subsection*{\textit{COMMENTARII IN PRIOrem Pauli Epistolam ad Timotheum. }}
\section{CAPVT I.}
\phantomsection
\addcontentsline{toc}{subsection}{\textit{\huge\textbf{P}\normalsize AVLVS APOSTOlus Iesu Christi secundùm ordinationem Dei saluatoris nostri, et Domini Iesu Christi spei nostrae: Timotheo germano filio infide, gratia, misericordia, pax a Deo patre nostro, et Chri sto Iesu Domino nostro. Quemadmodum rogauite, vt maneres Ephesi, cum proficiscerer in Macedo niam, volo denuncies quibusdam, ne aliter doceant, neque attendant fabulis et genealogiis nunquam finiendis: quae quaestiones praebent magis quam aedificationem Dei, quae in fide consistit. }}
\subsection*{\textit{\huge\textbf{P}\normalsize AVLVS APOSTOlus Iesu Christi secundùm ordinationem Dei saluatoris nostri, et Domini Iesu Christi spei nostrae: Timotheo germano filio infide, gratia, misericordia, pax a Deo patre nostro, et Chri sto Iesu Domino nostro. Quemadmodum rogauite, vt maneres Ephesi, cum proficiscerer in Macedo niam, volo denuncies quibusdam, ne aliter doceant, neque attendant fabulis et genealogiis nunquam finiendis: quae quaestiones praebent magis quam aedificationem Dei, quae in fide consistit. }}\pstart Paulus Apostolus.) Si soli Timotheo scripsisset: non opus erat adscribere hoc epitheton, et ita asserere, vt facit. Contentus enim proculdubio fuisset Timotheus solo nomine. Sciebat enim esse Chri sti Apostolum: nec indigebat vlla probatione, quo fidem illi deferret. Alios ergo potius respicit, qui non tam erant ad ipsum audiendum voluntarii: vel non tam facile eius dictis acquiescebant. Talium hominum causa, se Apostolum Christi esse testatur, ne paruipendere au sint, quod scribit. Et apostolatum suum confirmat ordinatione, vel decreto Dei: quia nemo Apostolum creare se ipse potest: sed quem Deus ordinauit, ille verus est Apostolus et dignus honore. Neque Patrem solum apostolatus sui facit authorem, sed illi Christum adiungit. Et certe in Ecclesiae gubernatione nihil agit Pater, nisi per Fi lium, et ita communiter cum Filio. Deum saluatorem vocat: quem titulum Filio tribuere saepius solet. Patri tamen eo ipso conuenit, quiais est qui nobis dedit Filium. Merito igitur illi salutis nostrae gloria adscribitur. Vnde enim salui sumus? Nempe quia sic nos dilexit Pateryt per Filium nos redimere voluerit, ac saluare. Quod au\pend
\section*{I. TIMOTH. }
\marginpar{[ p.2 ]}\pstart tem Christum appellat spem nostram, proprie illi conuenit. Tunc enim incipimus bene sperare, cum in Christum respicimus: quandoquidem in eo solo residet tota salutis nostrae materia.  \pend\pstart Timotheo germano filio.) Non est hoc paruae laudis elogium, quo significat Paulus, se pro vero ac non degenere filio agnoscere Timotheum, et talem ab aliis vult agnosci. Imo perinde Timotheum commendat, acsi esset alter ipse. Sed qui conueniet hoc cum sententia Christi? Nolite vobis patrem vocare in terra. Item cum illo Apostoli testimonio: Et si multos habetis patres secundùm carnemiivnus tamen est Pater spirituum. Respondeo, Paulum ita uibi vsurpare nomen patris, vt Deo nullam honoris sui particulam abroget, aut minuat. Vulgare dictum est: non pugnare quod subordinatur Tale est paternum in Paulo nomen respectu Dei. Vnicus in fide Pater omnium est Deus, quia omnes verbo suo et Spiritus sui virtute regenerat: quia solus est qui fidem confert. Sed quibus ad eam rem dignatur vti ministris, eos in honoris sui communicationem etiam admitcit, nihil tauien sibi derogando. Erat ergo Deus spiritualis Timothei pater: et quidem solus, proprie loquendo. Sed Paulus, qui minister fuerat Dei in gignendo Timotheo, quasi subalterno iure titulum sibi vindicat.  \pend\pstart Gracia, misericordia.) Secundum hoc praeter suum morem interposuit, forte singulari erga Timotheum amore impulsus, Porro non seruat exactum ordinem. Quod enim posterius est, priore ioco posuit. Gratiam scilicet, quae ex misericordia manat. Nam ideo nos in gratiam recipit Deus, et amore nos prosequitur, quia misericors est. Verum non est insolitum, causam subiungi effectui, explicationis causa. De vocabulis gratiae et pacis, alibi dictum est.  \pend\pstart Quemadmodum.) Vel oratio est defectiua: vel particula 'νae abundat. Vtroque modo sensus erit clarus. Primo reducit Timotheo in memoriam, qua de causa Ephesi manere rogatus fuerit. Non enim frustra eum a se dimiserat: sed potius vt vices illic suas sustineret. Hinc quasi per modum illationis hortatur, vt se opponat falsis doctoribus, qui puram doctrinam adulterabant. In eo quod Timotheo Ephesi mandauerat suas partes, obseruanda est pia eius solicitudo. Sie enim satagebat in colligendis pluribus Ecclesiis, vt priores non relin queret destitutas pastore. Et sane, vt ille inquit: non minor est virtus, quam quaerere, parta tueri, Denunciandi verbum potestatem significat. Vult enim Paulus eum instruere potestate ad alios coercendos. Ne aliter doceant.) Verbum Graecum, quo vsus est Paulus, est compositum. Ideóque transferri potest, vel secus docere et nouo modo: vel tradere diuersam doctrinam. Quod Erasmus vertit, sectari, non placet: quia de auditoribus posset accipi. Paulus autem eos intel  \pend
\section*{CAPVT I. }
\marginpar{[ p.3 ]}\pstart ligit, qui ambitionis eausa, nouam doctrinam ingerebant. Si legamus, aliter docere: latius hoc patebit. Nam prohibebit hoc verbo, ne patiatur Timotheus nouas docendi formas induci, quae cum vera et genuina, quam tradiderat, minus congruant. Sicut enim vna est Dei veritas, ita simplex eius tradendae ratio, minime fucata scilicet, et quae Spiritus maie statem potius quam humanae eloquentiae pompam resipiat. Ab hac quisquis recedit, deformat ipsam doctrinam et vitiat.Er go aliter docere, referre oportebit ad formam. Si legamus, aliud docere, ad materiam referetur. Notandum tamen est, non vocari aliam do ctrinam tantùm, quae palam cum pura Euangelii doctrina pugnet: sed quaecunque purum Euangelium, vel corrumpit nouis et aduentitiis figmentis, vel profanis speculationibus obscurat. Omnia enim hominum figmenta totidem sunt Euangelii corruptelae. Et qui scripturis ludunt profano more, vt Christianismum in artem ostentationis conuertant, Euangelio tenebras inducunt. Totum itaque hoc docendi genus alienum est a verbo Dei, et ab illa doctrinae puritate, in qua Ephesios manere iubet Paulus.  \pend\pstart Neque attendant fabulis.) Fabulas vocat, meo iudicio, non tam mendacia conficta, quam nugas aut ineptias, quae nihil habent solidi. Fieri enim potest, vt non sit aliquid falsum quod tamen erit fa bulosum. Qua significatione Tranquillus fabularem historiam posuit: et Liuius, fabulari, pro inepte, et absque ratione garrire. Nec dubium est, quin μύδος, quo verbo hic vsus est Paulus, Graecis significet βλοιαρίαν. hoc est, nugas. Porro cum speciem vnam addidit exempli causa, omnem dubitationem sustulit. Disputationes enim de genealogiis numerat inter fabulas: non quia commentitium sit, quidquid de illis dici potest: sed quia inane et infructuosum. Sic ergo resoluere licebit hunc locum: Non attendant fabulis: quales aut cuius generis sunt genealogiae. Atque illa est fabularis historia, cuius meminit Suetonius, quae etiam inter Grammaticos a sanis hominibus merito semper derisa fuit. Neque enim poterat non ridieula videri haec curiositas, vtili rerum cognitione posthabita, aetatem consumere in quaerenda Achillis, vel Aiacis prosapia, ingenium suum fatigare in numerandis Priami filiis. Si hoc non fertur in puerili scientia, vbi oblectationi est locus: quantominus in eoelesti sapientia erit tolerabile? Genealogias vocat infinitas, propterea quod inanis curiositas nullum habet modum: sed ex labyrintho subinde in labyrinthum reuoluitur.  \pend\pstart Quae quaestiones.) AEstimat a fructu doctrinam. Quaecunque enim non aedificat, repudianda est:etiam si nihil aliud habeat vitii. Quaecunque vero ad concertationes solum excitandas valet, duplici nomine damnanda est. Tales porro sunt omnes argutiae, quibus  \pend
\section*{I. TIMOTH. }
\marginpar{[ p.4 ]}\pstart ingenium suum ambitiosi homines exercent. Meminerimus ergo, ad hanc regulam exigendas esse omnes doctrinas: vt ea demum probetur, quae ad aedificationem facit: quae vero disceptationum materiam praebent absque fructu, respuantur tanquam Ecclesia Dei indignae. Hoc examen si fuisset aliquot supra seculis obseruatum, etiam si multis erroribus contaminata fuisset religio, saltem non adeo inualuisset diaboliea ista ars litigandi, quae Scholasticae theologiae nomen obtinuit. Quid enim illic continetur praeter rixas, aut ociosas speculationes, vnde nullus profectus redit? Proinde quoquisquein ea doctior, eo miserior iudicandus est. »cio quibus eam coloribus excusationum praetexant., ed nunquam efficient, vt l'aulus ipsas dam nando sit mentitus.  \pend\pstart AEdificationem Dei.) AEdificant enim eiusmodi argutiae in superbia aedificant in vanitate: sed non in Deo. AEdificationem Dei nominat, vel quae Deo probatur, vel quae secundùm Deum est. Eam porrô docet in fide consistere. Quo dicto non excludit charitatem, nec timorem Dei, nec poenitentiam. Nam haec omnia quid aliud sunt, quam fructus fidei? Quia ergo in sola fide sciebat fundatum esse totum Dei cultum, ideo satis habuit fidem nominare, vnde reliqua dependent.  \pend
\phantomsection
\addcontentsline{toc}{subsection}{\textit{Porro finis praecepti est charitas, ex puro corde, et conscientia bona, et fide non simulata. A quibus postquam nonnulli aberrarunt, deflexerunt ad vaniloquium, volentes esse Legis Doctores, non intelligentes quae loquuntur, neque de quibus affirmant. Scimus autemn quod Lex bona sit: si quis ea legitime vtatur, sciens illud, quod iusto non sit Lex posita, sed iniustis et inobsequentibus, impiis, et peccatoribus, irreuerentibus et profanis, parricidis et matricidis, homicidis, scortatoribus, masculorum concubitoribus, plagiariis, mendacibus, periuris, et si quid aliud est quod sanae doctrinae aduersetur, secundùm Euangelium gloriae beati Dei quod concreditum est mihi. }}
\subsection*{\textit{Porro finis praecepti est charitas, ex puro corde, et conscientia bona, et fide non simulata. A quibus postquam nonnulli aberrarunt, deflexerunt ad vaniloquium, volentes esse Legis Doctores, non intelligentes quae loquuntur, neque de quibus affirmant. Scimus autemn quod Lex bona sit: si quis ea legitime vtatur, sciens illud, quod iusto non sit Lex posita, sed iniustis et inobsequentibus, impiis, et peccatoribus, irreuerentibus et profanis, parricidis et matricidis, homicidis, scortatoribus, masculorum concubitoribus, plagiariis, mendacibus, periuris, et si quid aliud est quod sanae doctrinae aduersetur, secundùm Euangelium gloriae beati Dei quod concreditum est mihi. }}\pstart Quoniam nebulones illi, quibuscum negocium habebat Timotheus, legis praetextu se iactabantideo anticipat Paulus, ac docet  \pend
\section*{CAPVT I. }
\marginpar{[ p.5 ]}\pstart Legem non modd illis nihil suffragari, sed potius esse aduersam, cum Euangelio autem, quod ipse docuerat, oprime conuenire. Neque enim absimilis erat illorum defensio ei quam hodie nobis quaestionarii opponunt: nos scilicet nihil moliri aliud, quam vt aboleatur sacra theologia quasi eam soli in sinu suo fouerent. Ita illi de Lege loquebantur vt inuidiam Paulo conflarent. Quid autem ille? Quo fumos illos discutiat, vltro per occupationem suam doctrinam optime cum lege congruere demonstrat et eos omnes perperam lege abuti, qui in alium finem ea vtuntur. Quemadmodum hodie posita verae theologiae definitione, palam apparet, nos cupere ipsam restitutam, quae misere lacerata et deformata fuit ab istis nugatoribus, qui theologorum titulo frustra inflati, nihil praeter dilutas ac frigidas nugas tenent. Praeceptum hic pro lege accepit, partem pro toto.  \pend\pstart Charitas ex puro corde.) Ad hunc scopum si lex dirigenda est, vt instituamur in charitate, quae ex fide et bona conscientia procedat: sequitur exaduerso, peruersos esse legis interpretes, qui ad curiosas quaestiones eius doctrinam conuertunt. Caeterdm non magni interest, charitatem hoc loco ad vtranque legis tabulam, an ad secun dam tantùm referas. Iubemur Deum ex toto corde diligere, et proximos, vt nosmetipsos. Sed cum de charitate fit mentio in Scriptura, saepius ad secundum membrum restringitur. Hic quidem non dubitarem Dei et proximi charitatem intelligere, si solum charitatis nomen positum esset a Paulo. Sed cum addat fidem et bonam conscientiam et cor purum: non erit haec interpretatio aliena ab eius mente, et circunstantiae loci apte quadrabit. Haec legis summa est: vt Deum colamus syncera fide, et pura conscientia deinde vt nos mutuo amemus. Quisquis hinc deflectit, ille deprauat legem Dei, in alienum finem detorquendo. Sed hic exoritur dubitatio, quôd charitatem fidei praeferre videtur Paulus. Respondeo, eos qui ita sentiunt, nimis pueriliter ratiocinari. Neque enim si charitas priore loco nominatur, ideo priorem habet honoris gradum: quandoquidem simul docet Pau lus, ex fide eam manare. Causa autem effectum suum profecto ordine antecedit. Ergo si rite expendatur totus contextus, perinde valet, quod Paulus dicit, aesi dixisset: Hoc consilio data nobis est lex, vt in fide nos erudiret, quae mater est bonae conscientiae et charitatis. Ita exordium a fide non a charitate fieri debet. Cor purum et bona conscientia parum inter se differunt. Vtrunque ex fide manat. Nam de puro corde habetur Actorum 15. quod Deus fide corda purificet. De bona conscientia testatur Petrus, quod fundata sit in resurrectione Christi. Hinc porrô docemur, nusquam esse veram charitatem, vbi non est timor Dei, et conscientiae integritas. Nec leuiter praetereundum est, quod singulis addit epitheta. Nihil enim magis vulgare, sicu  \pend
\section*{I. TIMOTH. }
\marginpar{[ p.6 ]}\pstart ti nihil magis facile est, quam fidem, bonam conscientiam iactare. Se d quotusquisque est, qui factis ab omni se hypocrisi remotum esse demonstret? Praesertim notandum est epitheton fidei:quo significat fallacem esse eius professionem, vbi non apparet bona conscientia, vbi charitas se non exerit. Iam cum fide constet salus hominum, fide autem et bona conscientia et charitate constet perfectus Dei cultus:nihil mirum, si in his Paulus constituat legis summam.  \pend\pstart A quibus postquam.) Perstat in metaphora scopi, vel finis. Nam verbum ἀστοxᾶν, cuius participium hic habetur, significat, a scopo deflectere, vel aberrare. Est autem insignis locus: quia vaniloquii condemnat doctrinas omnes, quibus hic vnus finis non est propositus. Et simul omnium ingenia et cogitationes euanescere admonet, vbi alio respiciunt. Fieri quidem poterit, vt multos sui admiratione percellant inutiles argutiae. Stat tamen Pauli sententia, quidquid in pietate non aedificat, esse ματαιολογέαν. Quare summopere cauendum, ne aliud quaeramus in sacrosancto Dei verbo, quam solidam aedificationem, ne ipse alioqui in nobis vindicet eius abusum, tam seuero exemplo.  \pend\pstart Volentes esse.) Non taxat eos, qui aperte impugnant legis doctrinam, sed qui potius se eius titulo venditant. Tales negat quidquam intelligere, cum inutiliter in quaestionibus curiosis ingenium fatigant. Et simul eorum arrogantiam perstringit, cuni addit, de quibus affirmant. Nulli enim audaciores reperientur ad res incognitas temere asserendas, quam eiusmodi fabularum magistri. Videmus hodie quo fastu et supercilio magistrales suas determinationes crepent Sorbonicae scholae. Quibus autem de rebus? Nempe quae penitus abditae sunt humanis ingeniis, nullo Scripturae verbo, nulla vnquam reuelatione patefactae. Maiore confidentia purgatorium suum affirmant, quam mortuorum resurrectionem. Quod de sanctorum intercessione commenti sunt nisi habeatur pro certo oraculo, roram religionem conuulsam esse clamitant. Quid de coelestibus Hierarchiis, de relationibus, et similibus commentis commemorem immensos labyrinthos Genus ipsum est infinitum. In his omnibus pronunciat Apostolus impleri quod veteri prouerbio celebratum est, Audacem esse inscitiam.  \pend\pstart Scimus quod lex bona.) Iterum calumniam paeoccupat, qua ipsum grauabant. Nam quoties eorum ostentationi resistebat. hoc quasi clypeo se tuebantur. Quid ergo? Vis ne sepultani esse legem, et deletam ex hominum memoria? Fatetur ergo Paulus, vt hac calumnia eos depellat, bonam esse legem: sed vsum eius legitimum requiri. Vbi argumentatur a coniugatis. Nomen enim legitimi ab ipsa lege deductum est. Verùm vltra progreditur. Nam et legem pul\pend
\section*{CAPVT I. }
\marginpar{[ p.7 ]}\pstart chre cum sua doctrina conuenire docet, et eandem in ipsos retor\pend\pstart Taala Quod iusto non sit lex.) Non instituit Apostolus disputare de toto legis officio:sed homines ipsos respicit. Fere enim contingit, vt qui summi legis zelotae videri affectant, tota vita se prodant sum mos esse contemptores. Cuius rei praeclarum hodie specimen relucet in assertoribus iustitiae operum, et liberi arbitrii patronis. Perpetud in ore habent perfectam vitae sanctitatem, merita, satiffactiones. Tota auteni vita clamitat, eos plusquam impios et profanos esse, pro uocare quibuscunque modis iram Dei, et iudicium eius secure con temnere. Magnifice extollunt liberam boni et mali electionem. Factis autem produnt, se esse Sathanae mancipia, et quidem arctissimis seruituris compedibus districta. Cum tales essent Pauli aduersaril, vt stolidam eorum insolentiam compesceret, admonuit, legem esse qua si Dei gladium, ad ipsos iugulandos exertum. Sibi vero suique similibus non esse causam, cur legem horrori vel odio haberet, quae iu stis, hoc est, piïs et voluntariis Dei cultoribus non sit aduersa. Nec me latet, doctos quosdam viros subtiliorem ex his verbis sensum elicere: quasi penitus theologice de legis natura disputaret Paulus. Colligunc enim legem nihil ad filios Dei, qui spiritu sunt regenerati, per tinere: quia iustis non sit lata. Sed loci circunstantia simplicius me accipere cogit hanc sententiam. Sumit enim vulgare illud principium:ex malis moribus natas esse bonas leges: ac legem Dei asserit latam esse ad coercendam impiorum nequitiam:quia qui sponte boni sunt, legis imperio non indigent. Nunc quaeritur, an vnus quispiam mortalium eximatur ab hoc numeros? Respondeo, Paulum hic nominare iustos, non qui sint omnibus numeris absoluti: qualis nemo reperierur sed qui praecipuo cordis affectu ad bonum aspirant: vt illis pium desiderium sit quasi lex voluntaria, absque alieno impulsu aut freno. Ergo aduersariorum impudentiam coarguere voluit. qui legis nomine se armabant aduersus pios viros, ac veram legis regulam tota vita exprimentes eum ipsi maxime opus haberent lege, et tamen eam non magnopere curarent. Quod melius opposito membro exprimitur. Quod si quis non admittat, tacite vel oblique exprobrari aduersariis, quae hic flagitia recenset Paulus, simplex tamen manebit calumniae depulsio: quod si recto et non fucato legis studio ducerentur, instruere se potius illius armis debuissent ad gerendum cum sceleribus et maleficiis bellum, quam eius colorem suae ambitioni et ineptae garrulitati praetexere.  \pend\pstart Iniustis et inobsequentibus.) Iniustos vocat exleges, a qui bus parum differunt inobsequentes, quos secundo loco posuit. Peccatores accipit pro sceleratis, aut qui vitae sunt turpis ac flagitiosae.  \pend
\section*{I. TIMOTH. }
\marginpar{[ p.8 ]}\pstart Pro irreuerentibus et profanis non absurde quis transferret profanos et impuros. Sed ego in vocibus leuis momenti nolui esse tam morosus. Plagium apud veteres alieni serui subreptio aut solicitatio vocabatur, aut falsa liberi hominis venditio. Si quis plura requi rat, petat a Iureconsultis, ad legem flauiam. Caeterùm genera quaedam attigit hic Paulus, quae breuiter omnes transgressiones compre hendunt. Radix est contumacia et rebellio, quam primis duobus vocabulis notauit. Per impios et peccatores videntur prioris et secundae tabulae transgressores designari. His subiicit profanos et im puros, vel turpi vita inquinatos. Cum autem tribus praecipue modis laedantur proximi, violentia, fraudibus et libidine: tres istos modos ordine taxauit: vt promptum est videre.  \pend\pstart Si quid est aliud.) Hoc membro contendit, aded̀ non pugnare suum Euangelium cum lege, vt sit optima eius confirmatio. Testatur enim se comprobare sua praedicatione eandem sententiam, quam Dominus lege sua tulit aduersus ea omnia, quae sanae doctrinae contraria sunt. Vnde sequitur, eos legis animam minime tene re, qui ab Euangelio recedunt: sed vmbram duntaxat sectari. Euangelium gloriae, pro glorioso. In quo tamen acriter eos perstringit, qui extenuando Euangelio imminebant:in quoDeus suam gloriam exerit. Nominatim addit sibi fuissecreditum:vt omnes intelligant non aliud esse Dei Euangelium, quam quod praedicat: ideoque fabulas omnes quas prius reprehendit, alienas esse tam a lege quam ab Euangelio Dei.  \pend
\phantomsection
\addcontentsline{toc}{subsection}{\textit{Et gratiam habeo, qui me potentem reddidit Christo Iesu Domino nostro, quod fidelem me iudicauit ponendo inministerium, qui prius eram blasphemus et persecutor, et violentus. Sed et misericordiam adeptus sum, quod ignorans feci in incredulitate. }}
\subsection*{\textit{Et gratiam habeo, qui me potentem reddidit Christo Iesu Domino nostro, quod fidelem me iudicauit ponendo inministerium, qui prius eram blasphemus et persecutor, et violentus. Sed et misericordiam adeptus sum, quod ignorans feci in incredulitate. }}\pstart Gratiam habeo.) Magna est Apostolatus dignitas, quam sibi Paulus vindicauit: ipse autem superioris vitae respectu nequaquam, tanto honore dignus censeri poterat:ideo ne arrogantiae damnetur, necessa rio ad personae suae mentionem descendit. Et suae quidem indignitatis liberam confessionem edit. Caeterùm nihilominus se gratia Dei esse Apostolum. Imo quod videbatur authoritatem illi abrogare, in suum commodum tranffert: quod eo magis in se reluceat Dei gratia. Christo gratias dum agit, a seipso remouet inuidiam:  \pend
\section*{CAPVT I. }
\marginpar{[ p.9 ]}\pstart et ansam huie quaestioni praecidit, dignus sit necne functione, tam ho norifica. Multi quidem eadem verborum forma humilitatem praeseferunt: sed procul absunt a Pauli sinceritate. Qua de re autem agit gratias? quod in ministerio fuerit collocatus. Inde enim colligit, se fidelem iudicatum. Neque enim Christus ambitiose, quosuis assumit: sed tantùm eligit idoneos. Ergo quoscunque honore dignatus est, hos tanquam dignos amplectamur. Atqui videtur hoc modo Paulus, fidelitatem qua prius pollebat, causam facere suae voca tionis. Si ita esset, ficta essetatque absurda gratiarum actio. Neque enim tam Deo acceptum referret suum Apostolatum, quam proprio merito. Nego igitur hunc esse sensum: quod ideo in ordinem Apostolicum fuerit cooptatus: quia praecognita esset Deo eius fides. Nihil enim praeuidere boni in illo potuit Christus, nisi quod Pater in eum contulerat. Quare manet illud semper verum: Non vos me elegistis: sed ego elegi vos. Potius argumentum suae fidei inde ducere voluit, quia Apostolus creatus erat a Christo. Nam quos ita creat, quasi suo decreto pronunciat habendos esse fideles. In summa, iudicium hoc non refertur ad praescientiam: sed potius testimonium significat quod hominibus redditur. Acsi dixisset: Gratiam habeo Christo, qui me in ministerium vocando, palam fecit, meam sibi fidem probari. Interponit et aliud Christi beneficium, quod ipsum confirmet. Quo verbo, intelligo continuam gratiae subministrationem. Neque enim satis fuisset, semel esse fidelem declaratum, nisi eum perpetuo auxilio confirmasset Christus. Vtrunque ergo se ex Christi gratia fatetur habere, quod semel euectus est, et quod permaneat in statu.  \pend\pstart Blasphemus.) Blasphemus aduersus Deum, persecutor et violentus aduersus Ecclesiam. Videmus quam ingenue agnoscat, quod illi probro dari poterat, quam non extenuet sua peccata. Nam persecutoris nomine non contentus, furorent suum ac saeuitiam magis exprimere voluit.  \pend\pstart Quia ignorans.) Veniam, inquit, meae incredulitatis obtinui, quia ex ignorantia manabat. Nam persecutio et violentia nihil quam incredulitatis fructus erant. Atqui videtur innuere, nullum esse veniae locum, nisi vbi suppetit ignorantiae excusatio. Quid ergoran si quis sciens peccauerit, nunquam ignoscet Deus? Respondeo, notandum esse incredulitatis nomen, quod dictum Pauli ad priorem legis tabulam restringit. Secundae tabulae transgressiones, etiam si voluntariae sunt, condonantur. Qui sciens et volens priorem tabulam violat, quia recta se Deo opponit, peccat in spiritum sanctum. Neque enim infirmitate labitur, sed aduersus Deum impie ruendo, certum reprobationis suae signum edit. Et hinc colligi potest definitio peccati in spiritum sanctum primum, quod sit directa  \pend
\section*{I. TIMOTH. }
\marginpar{[ p.10 ]}\pstart in Deum contumacia in transgressione prioris tabulae. deinde quod sit malitiosa veritatis reiectio. Vbi enim non respuitur destinata malitia Dei veritas, non resistitur Spiritui sancto. Denique incredulitas loco generis ponitur: malitiosum autem propositum, quod opponitur ignorantiae, est quasi differentia. Perperam itaque faciunt, qui peccatum in Spiritum sanctum constituunt in trangressione secundae tabulae. Perperam et illi, qui inconsideratum et caecum impetum damnant tanti criminis. Nam tunc demum peccatur in spi ritum sanctum, cum voluntarium bellum suscipiunt aduersus Deum homines mortales, vt lucem spiritus sibi oblatam extinguant. Horrenda est enim eiusmodi impieras, et immanis audacia. Nec dubium est, quin tacita denunciatione terrere voluerit omnes, qui semel illu minati fuerunt, ne aduersus cognitam veritatem impingant: quia exitialis sit casus. Nam si propter ignorantiam Paulo suas blasphemias ignouit Deus, qui scienter dataque opera blasphemant, nullam veniam sperare debent. Sed videtur id frustra dici. Nusquam enim sine ignorantia porest esse incredulitas, quae semper caeca est. Respondeo, ex infidelibus alios ita esse caecos, vt falsa recti imaginatione fallantur:alios autem sic excaecatos, vt malicia tamen praeualeat. Paulus non omnino vacuus erar prauo affectu. Sed inconsideratus zelus abripiebat hominem, vt rectum eise putaret quod agebat. Non erat ergo aduersarius Christi proposito, sed errore et inscientia. Phari saei, qui mala conscientia Christum calumniabantur, non omnino crant immunes ab errore et ignorantia sed ambitio, et impium sanae doctrinae odium, imo furiosa aduersus Deum contumacia illos inci tabat, vt malitiose et consulto, non inscienter, Christo se opponerent.  \pend
\phantomsection
\addcontentsline{toc}{subsection}{\textit{Exuberauit autem supramodum gratia Domini nostri cumfide, et dilectione quae est in Christo lesu. Fidelis serino, et dignus omnino qui accipiatur, quod Christus Iesus venit in mundum vt peccatores saluos faceret, quorum primus sum ego. Verùm ideo misericordiam sum adeptus, vt in me primo ostenderet Iesus Christus omnem clementiam, in exemplar iis qui creditur i essent in ipso in vitam aeternam. Regi autem seculorum immortali, inuisibili, soli sapienti Deo, honor et gloria in secula seculorum, Amen. }}
\subsection*{\textit{Exuberauit autem supramodum gratia Domini nostri cumfide, et dilectione quae est in Christo lesu. Fidelis serino, et dignus omnino qui accipiatur, quod Christus Iesus venit in mundum vt peccatores saluos faceret, quorum primus sum ego. Verùm ideo misericordiam sum adeptus, vt in me primo ostenderet Iesus Christus omnem clementiam, in exemplar iis qui creditur i essent in ipso in vitam aeternam. Regi autem seculorum immortali, inuisibili, soli sapienti Deo, honor et gloria in secula seculorum, Amen. }}
\section*{CAPVT I. }
\marginpar{[ p.11 ]}\pstart Iterdm Dei erga se gratiam amplificat, non modo inui dix tol lendae ac gratitudinis suae testandae causa: sed etiam vt eam opponat improborum calumniis, qui toti in hoc erant, vt eius Apostolatum deprimerent. Cum autem dicit exuperasse, et quidem supramodum, innuit deletam esse superiorum memoriam, adeôque absorptam, ne quid illi obesset apud bonos talem prius fuisse.  \pend\pstart Cuni fide et dilectione.) Vtrunque potest ad Deum referri, hoc sensu, quod Deus veracem se praestiterit, et dilectionis in Christo suae specimen praebuerit, cum dignatus ipsum fuerat sua gratia. MaIo tamen simplicius accipere: vt fides et dilectio signa sint ac testimo nia eius gratiae, cuius meminerat: ne quid frustra et temere iactare pu taretur. Ac fides quidem incredulitati opponitur dilectio in Christo, saeuitiae quam exercuerat aduersus fideles acsi diceret, se ita fuisse a Deo mutatum, vt alius sit ac nouus homo.  \pend\pstart Fidelis sermo.) Postquam suum ministerium ab infamia et iniquis obtrectationibus asseruit, eo non contentus, retorquet in suum commodum, quod illi poterat ab aduersariis obiici probri loco. Vtile enim Ecelesiae fuisse docet, se prius talem extitisse, quam ad Apostolatum vocaretur, quia hoc modo Christus peccatores omnes, tanquam dato pignore, ad certam spem veniae obtinendae vocauerit. Nam cum velut ex immani truculentaque belua versus fuerat in pastorem, insigne gratiae suae specimen ediderat Christus, vnde fiduciam omnes conciperent, nullis, quamlibet grauibus et atrocibus peccatis aditum sibi ad salutem praecludi. Verùm primô generalem sententiam profert: Christum venisse, vt saluos faceret peccatores: et eam praefatione ornat, qua vti solet in rebus maxime seriis. Vt certe in religionis doctrina hoc praecipuum est caput: ad Christum venire, vt in nobis perditi, salutem ab eo recuperemus. Praefatio ergo haec nobis sit instar buccinae sonantis ad publicandum gratiae Christi praeconium: quo plus fidei apud nos habeat. Sit nobis quasi sigillum ad impriniendam cordibus nostris de remissione peccatorum fiducian, quae difficulter alioqui in animos hominum penetrat. Quae etiam causa fuit Paulo, cur attentionem excitaret his verbis: Fidelis sermo. Nam vt millies in Christo salutem nobis offerat Deus Pater, et Christus ipse de officio suo concionetur: non tamen propterea desinimus trepidare, vel saltem disputare nobiscum, an ita sit. Proinde quoties nobis in mentem veniet vlla de peccatorum remissione dubitatio, hoc veluti clypeo fortiter eam repellere discamus: Veritatem esse indubiam, etdignam quae absquae controuersia recipiatur.  \pend\pstart Peccatores.) Subest in hoc verbo emphasis. Nam qui offieium Christi esse fatentur, saluare, cogitationem hanc difficilius ad\pend
\section*{I. TIMOTH. }
\marginpar{[ p.12 ]}\pstart mittunt: quod eiusmodi salus ad peccatores pertineat. Semper enim abripitur sensus noster ad respectum dignitatis. Simul atque indigni tas apparet, concidit fiducia. Proinde quo quisque peccatis suis magis grauatur, eo maiore animo ad Christum confugiat, hac scilicet doctrina fretus, venisse vt salutem non iustis, sed peccatoribus afferat.  \pend\pstart Quorum primus.) Caue existimes modestiae causa Apostolum mentitum esse. Veram enim non minus quam humilem con fessionem edere voluit, atque ex intimo cordis sensu depromptam. Sed quaeret hic quispiam, curnam se inter peccatores praecipuum nu merat, qui tantum sanae doctrinae ignorantia lapsus est, alioqui tota vita inculpabilis apud homines? Atqui his verbis admonemur, quam gra ue sit apud Deum et atrox crimen infidelitas: praesertim vbi accedit obstinatio et saeuiendi rabies. Facile quidem est apud homines, totum id quod de se confessus est Paulus, extenuare colore zeli inconsiderati. Sed pluris obedientiam fidei facit Deus, quam vt pro leui delicto incredulitatem cum peruicacia coniunctam ducat. Quare diligenter notandus est hic locus, hominem coram mundo non modo innoxium, sed eximiis virtutibus praestantem vitaeque laudatissimae, quia Euangelii doctrinae fuerit aduersarius, propter incredulitatis suae contumaciam, vnum ex grauissimis peccatoribus censeri. Nam inde colligere promptum est, quid valeant coram Deo omnes hypocritarum pompae, dum contumaciter Christo resistunt.  \pend\pstart Vt primo ostenderet.) Cum dicit, primo, alludit ad id quod nuper dixerat, se primum esse inter peccatores. Significat autem, statim ab initio Deum proposuisse tale exemplar, quod tanquam ex illustri, excelsôque theatro conspici posset, ne quis diffideret paratam sibi fore veniam, modo fide accederet ad Christum. Et certe praeuenitur nostra omnium diffidentia, dum eius, quam quaerimus, gratiae typum in paulo sic videmus expressum.  \pend\pstart Regi autem seculorum.) Prae ardore erumpit in istam exclamationem:quia deerant verba, quibus suam gratitudinem exprimeret. Nam epiphonemata praecipue locum habent, vbi abrumpere orationem cogimur, quia rei magnitudo superat. Quid autem Pauli conuersione admirabilius? Quanquam nos simul omnes admonet suo exemplo, nunquam de gratia diuinae vocationis esse cogitandum, quin tandem efferamur admiratione. Adde quod hoc tam magnificum encomium totius superioris vitae memoriam absorbet. Qualis enim abyssus est Dei gloria? Epitheta quae hic tribuit Deo, tametsi perpetua sunt, tamen proprie conueniunt loci circunstantiae. Vocat regem aeternum vel seculorum, in quem nulla cadit mutatio:inuisibilem, quia lucem habitat inaccessam, vt postea dicet.  \pend
\section*{CAPVT I. }
\marginpar{[ p.13 ]}\pstart Solum sapientem:quia infatuet ac vanitatis damnet omnem hominum sapientiam. Summa, redit ad eam clausulam, qua vtitur ad Ro. vndecimo. O profunditatem diuitiarum et c. Quam abscondira sunt Dei consilia? Quam inserutabiles viae eius? Dubium tamen est, velit ne Deo gloriam omnem soli vindicare, an vocet ipsum solum sapientem:an vero solum Deum esse dicat. Secundus sensus magis mihi probarur. Praesentis enim causae maxime interfuit, arcano Dei conuilio subiici quanuis hominum intelligentiam.  \pend
\phantomsection
\addcontentsline{toc}{subsection}{\textit{Hoc praeceptum commendo tibi, fili Timothee, secundùm praecedentes super te Prophetias: vt milites in illis bonam militiam, habens fidem et bonam conscientiam, a qua auersi quidam circa fidem naufragium fecerunt: ex quibus sunt Hymenaeus et Alexander, quos tradidi Sathanae, vt discant non maledicere. }}
\subsection*{\textit{Hoc praeceptum commendo tibi, fili Timothee, secundùm praecedentes super te Prophetias: vt milites in illis bonam militiam, habens fidem et bonam conscientiam, a qua auersi quidam circa fidem naufragium fecerunt: ex quibus sunt Hymenaeus et Alexander, quos tradidi Sathanae, vt discant non maledicere. }}\pstart Hoc praeceptum.) Quidquid de persona sua interposuit, fuit quasi digressio ab instituto. Nam cum Timotheum instruere authoritate vellet, summa ipsum praeditum esse oportuit. Mature itaque opinionem refutauit, quae sibi obstare poterat. Nunc autem, postquam probauit, nihilo minoris apud pios aestimandum esse suum Apostolatum, quia aliquando Christi regnum oppugnasset. quasi sublato obstaculo, ad exhortationis suae contextum redit. Ergo praeceptum est, euius initio meminit. Filium vocando, non suam modo beneuolentiam apud eum testatur, sed aliis etiam commendat ipsum hoc nomine. Porro vt magis eum animet, reducit in memoriam, quale a spiritu Dei testimonium habuerit. Nam erat hic non paruus stimulus, quod diuinitus ministerium eius fuerat approbatum et quod prius Dei oraculis vocatus fuerat, quam electus hominum suffragiis. Turpe est non respondere hominum expectationi. Quanto itaque turpius erit, irritum, quoad in te sit, facere Dei iudicium?o ed primum sciendum est, de quibus loquatur propheriis. Putant aliqui reuelatione admonitum fuisse Pauluni, vt Timocheo munus iniungeret. Quod verum esse fateor: sed addo reuelationes ab aliis proditas. Neque enim Paulus numero plurali frustra vsus est. Proinde ex eius verbis colligi mus, plures fuisse editas Prophetias de Timotheo, quae ipsum Ecclesiae commendarent. Nam cum iuuenis adhuc esset, poterat contemptui esse illius aetas: expositus etiam calumniis fuisset Paulus, quod adulescen tulos ante tempus promoueret ad senile officium. Praetera Deus  \pend
\section*{I. TIMOTH. }
\marginpar{[ p.14 ]}\pstart illum ad res magnas et arduas destinauerat. Non enim fuit vnns ex vulgo: sed Apostolis proximus, qui Paulo absente saepe eius perso nam substinebat. Quare opus habuit singulari testimonio vnde pateret, non hominum temeritate tantùm illi deferri: sed electum ab ipso Deo esse. Neque enim tunc ordinarium fuit, aut multis commune, ornari Prophetarum encomiis. Sed quia in Timotheo peculiares erant causae, noluit eum Deus ab hominibus admitti, nisi sua voce iam probatum: noluit eum ad obeundam functionem accedere, nisi Propheticis oraculis accitum. Sicuti etiam Paulo et Barnabae contigit, dum Gentibus ordinati sunt doctores. Erat enim res noua et insolens: nec alioqui potuissent famam temeritatis effugere. Obiiciet nunc quispiam:si Deus per suos Proplietas ante pronuncia uerat, qualis futurus esset minister Timotheus: quorsum attinebat eum moneri, vt se talem praestaret? An potuit Dei oraculis fidem abrogare? Respondeo, non aliter potuisse euenire, quam Deus promi serat. Sed interim non debuit Timotheus desidiae se ac torpori dare, sed viuum agileque organum Dei prouidentiae, obsequendo, se exhibere. Quare Paulus non abs re dum calcar illi addere vult, Prophetias commemorat, quibus se pro illo Deus quasi sponsorem Ecclesiae constituerat. Nam inde admonebatur, ad quid vocatus foret. Ideóque addit: milites in illis. quo significat, Timotheum, tali Dei approbatione fretum, animosius debere pugnare. Quid enim plus alacritatis addere nobis vel debet, vel potest, quam cum scimus nos diuinitus ordinatos ad agendum quod agimus? Haec arma sunt nostras haec praesidia, quibus muniti nunquam deficimus. Militiae nomine, subindicat certandum esse. Atque in vniuersum piis id conuenit: proprie tamen Christianis doctoribus, qui sunt velut antesignani aut duces. Perinde ergo est, acsi diceret: Tu si non potes absque certamine munus tuum exequi, memineris te Dei oraculis ad victoriae spem certissimam esse instructum. Atque hac recordatione te excites. Bona enim est militia, hoc est gloriosa et salutaris, quam sub Dei auspiciis militamus.  \pend\pstart Habens fidem et bonam conscientiam. Fidei nomen generaliter accipio, pro sana doctrina. Qua significatione etcapite tertio mysterium fidei ponit. Et certe duo haec praecipue in doctore requiruntur, vt puram Euangelii veritatem retineat deinde eam administret bona conscientia, et sincero zelo. Vbi vtrunque istorum aderit, sponte etiam reliqua sequentur. A qua auersi quidant.) Ostendit quam necessarium sit, bonam conscientiam cum fide coniungi: quia econuerso malae conscientiae poena sit a recta fide defectio. Qui non pura animi integritate seruiunt Domino, sed prauis affectibus indulgent, etiam si prius sanam intelligentiam habuerint, ab ea prorsus excidunt. Locus dili\pend
\section*{CAPVT I. }
\marginpar{[ p.15 ]}\pstart genter obseruandus. Seimus incomparabilem esse sanae doctrinae chesaurum. Quare nihil magis timendum, quam vt nobis auferatur. Atqui vnicam eius custodiendi rationem Paulus hic praescribit, si ob seratus fuerit bona conscientia. Idque quotidie experimur. Vnde enim fit, vt abiecto Euangelio tam multi ad impias sectas proruant, vel errorum portentis implicentur? Nempe quia hoc excaecationis gene re hypocrisin vlciscitur Deus. Quemadmodum sincerus Dei timor ad perseuerantiam confirmat. Duo hinc discenda sunt. Primûm admonentur Doctores et Euangelii ministri, atque in eorum persona totae Ecclesiae, quantopere fucatam et fallacem verae doctrinae professionem horrere debeant, quam adeo seuere puniri audiunt. Deinde tollitur scandalum, quod plerosque non leuiter turbat, dum aliquos vident, qui prius Christo et Euangelio nomen dederant, non modô relabi ad pristinas superstitiones, sed quod longe deterius est, prodigiosis erroribus fascinatos teneri. Nam talibus exemplis Euangelii maiestas palam vindicatur: ac palam ostendit Deus se eius profanationem nullo modo ferre posse. Atque id omnium seculorum experien tia docuit. Quidquid errorum ab exordio Christianae Ecclesiae extitit, ab hoc fonte manauit: quia in aliis ambitio, in aliis auaritia sincerum Dei timorem obruebat. Mala igitur conscientia omnium haereseon mater. Et hodie permultos videmus, quia rectam fidem non sincere amplexi fuerant, pecudum instar ad Epicureorum deliria rapi, vt ipsoruni hypocrisis detegatur. Quinetiam cum passim grassetur Dei contemptus, et flagitiosi perditique omnium fere ordinum mores demonstrent, nihil aut quam minimum esse integritatis in mundo: valde timendum est, ne breui lux, quae accensa fuerat, extinguatur: et paucissimis Deus puram Euangelii intelligentiam relinquat. Metaphora a naufragio sumpta aptissime quadrat. Nam innuit, vt salua fides ad portum vsque perueniat, nauigationis nostrae eursus bona conscientia regendum esse. Alids naufragii esse periculum: hoc est, ne fides mala conscientia, tanquam gurgite aliquo mergatur.  \pend\pstart Hymenaeus et Alexander.) Priorem iterum nominabit in secunda Epistola vbi etiam naufragii, quod fecerat, species notatur. Dicebat enim resurrectionem esse factam. Credibile est, Alexandrum quoque tam absurdo errore fuisse dementatum. Et hodie miramur, si quosdam variis praestigiis deludat Sathan, cum vnum ex Pauli comitibus tam horrendo praecipitio periisse cernamus? Caeterùm ambos illi nominat tanquam homines notos. Ego quidem non dubito, quin hic ille sit Alexander, cuius mentionem facit Lucas Actorum 19. qui tumultum ficta et profana sua oratione composuit. Erat autem Ephesius. Et hanc Epistolam diximus maxime scriptam fuisse Ephesiorum causa. Nunc audimus qualem habuerit exitum.  \pend
\section*{I. TIMOTH. }
\marginpar{[ p.16 ]}\pstart Audientes colamus bona conscientia fidei nostrae possessionem, vt salua nobis ad extremum maneat.  \pend\pstart Quos tradidi Sathanae.) Quemadmodum diximus in quintum caput prioris Epistolae ad Corinthios, sunt qui interpretantur extraordinariis flagellis fuisse coercitos: idque referunt ad δυναμίν, cuius meminit Paulus eiusdem Epistolae capite decimo. Nam vt dono sanationis praediti erant Apostoli ad testandam erga pios Dei gratiam et beneficentiam, ita aduersus impios et rebelles armati erant potentia: vel vt vexandos Diabolo traderent: vel vt aliis flagellis in eos animaduerterent. Cuius potentiae specimen edidit Petrus in Anania et Sapphira Paulus vero in Bar-Iesu mago. At mihi de excommunicatione exponere magis placet. Nam vt credamus aliter, quam excommunicatione castigatum fuisse Corinthium illum incestum, nulla probabilis coniectura affertur. Si autem illum excommunican do Sathanae tradidit Paulus: cur non eadem locutio idem hic valebit? Adde quod vim excommunicationis optime explieat. Nam cum in Ecclesia sedem regni sui habeat Christus, extra non est nisi Sathanae dominium. Proinde qui ab Ecclesia exterminatur, eum tantis per sub tyrannide Sathanae degere necesse est, dum Ecclesiae reconciliatus, ad Christum redeat. Vnum excipio, quod, pro mali atrocitate. perpetuum istis anathema denunciare potuit. De quo nihil tamen ausim pro certo asserere. Sed quid sibi vult proxima particula: Vt discant non maledicere? Nam qui ab Ecclesia eiectus est, plus licentiae sibi vsurpat: quia tanquam communis disciplinae iugo solutus, petulantius insurgit. Respondeo, vtcunque exultet eoruni improbitas, aditum tamen praecludi,ne sua contagione gregem inficiant. Nam eo maxime nocent improbi, dum eiusdem fidei praetextu se insinuant. Nocendi ergo facultas illis eripitur, dum notantur publica infamia. vt nemo simplicium ignoret, profanos esse homines et execrandos : ideóque omnes ab eorum communicatione abhorreant. Ipsos etiam interdum contingit irrogata sibi talis ignominiae nota confusos. a proteruia desistere. Ergo etiam si improbiores reddantur aliquando: tamen non semper inutile est ad domandam ferociam temedium.  \pend
\endnumbering\beginnumbering\section{CAPVT II.}
\phantomsection
\addcontentsline{toc}{subsection}{\textit{\huge\textbf{A}\normalsize Dhortor igitur, vt ante omnia fiant deprecationes, obsecrationes, interpellationes, gratiarum actiones pro omnibus hominibus, pro regibus et omnibus in eminentia constitutis, vt placidam et quietam vitam degamus cum omni pietate et honestate. }}
\subsection*{\textit{\huge\textbf{A}\normalsize Dhortor igitur, vt ante omnia fiant deprecationes, obsecrationes, interpellationes, gratiarum actiones pro omnibus hominibus, pro regibus et omnibus in eminentia constitutis, vt placidam et quietam vitam degamus cum omni pietate et honestate. }}
\section*{CAP. II. }
\marginpar{[ p.17 ]}
\phantomsection
\addcontentsline{toc}{subsection}{\textit{ Hoc enim bonum et acceptum coram saluatore nostro Deo, qui omnes homines vult saluos fieri, et ad agnitionem veritatis venire. }}
\subsection*{\textit{ Hoc enim bonum et acceptum coram saluatore nostro Deo, qui omnes homines vult saluos fieri, et ad agnitionem veritatis venire. }}\pstart Adhortor igitur.) Haec pietatis exercitia exercent nos in sincero cultu Dei ac timore fouentque bonam conscientiam de qua dixerat. Quare non abs re, illatiua particula vtitur: quia ex superiore praecepto exhortationes istae pendent. Ac initio quidem de publicis orationibus disserit: quas iubet non pro fidelibus modo concipi, sed pro vniuerso genere humano. Poterant enim nonnulli ita secum reputare: Cur de infidelium salute essemus soliciti, quibuscum nihil est nobis necessitudinis? Nónne satis est, sifratres pro fratribus mutuo oremus, ac commendemus Deo totam suam Ecclesiam? Nihil enim ad nos extranei. Huic sinistrae opinioni occurrit Paulus, ac iubet Ephesios suis precibus complecti omnes mortales, nec eas ad corpus Ecclesiae restringere. Porrd quid inter se differant quatuor species quas enumerat, fateor me non penitus tenere. Puerile quidem est, quod Pauli verba Augustinus ad ritus suo tempore vsitatos detorquet. Sim plicior est eorum expositio, qui deprecationes esse putant, quibus petimus liberari a malis: obsecrationes, quibus vtilia nobis poscimus: interpellationes, quibus apud Deum iniurias nobis illatas deploramus. Quanquam ego non ita subtiliter distinguo, vel saltem aliam distinctionem magis probo. *ροσωυχας sane Graeci vocant omne genus orationis: δεή σας autem eas precationum formas, ybi certum aliquid petitur. Hoc modo, inter se conuenirent haec duo nomina, tanquam genus et species. ιν rιύfας vocare solet Paulus quas alii pro aliis suscipimus preces. Pro quo Latina transsatio habet intercessiones. Quanquam Plato in Alcibiade secundo aliter accipit, nenpe, pro seipso certam precationem concipere. In ipsa autem libri inscriptione et compluribus locis satis quod dixi ostendit, προσευχὴν nomen esse generale. Sed ne plus aequo laboremus in re non necessaria: Paulus, meo iudicio, simpliciter iubet, quoties publicae orationes habentur, supplicare et deprecari pro omnibus: etiam qui in presentia nihil nobiscum habent coniunctionis. In gratiarum actione nihil est obscurum. Nam quemadmodum vult infidelium salutem commendari Deo, ita et propter laetos eorum ac prospe ros successus gratias agi. Ista enim admirabilis Dei bonitas, quam quoridie ostendit, cum Solem suum oriri facit super bonos et malos, digna est quam laude prosequamur. Et charitas nostra vsque ad indignos extendere se debet.  \pend
\section*{I. TIMOTH. }
\marginpar{[ p.18 ]}\pstart Pro Regibus.) Nominatim de Regibus meminit et magistratibus reliquis, quia prae aliis exosi esse poterant apud Christianos. Quotquot enim erant illo tempore magistratus, totidem erant quasi iurati Christi hostes. Poterat igitur obrepere ista cogitatio, non esse pro illis orandum, qui totas vires opesque suas conferrent ad oppugnanduni Christi Regnum, cuius propagatio imprimis optanda est. Occurrit autem Apostolus, ac diserte, iubet pro illis etiam precari. Et certe, non efficit hominum prauitas quominus amanda sit Dei institutio. Proinde cum magistratus ac principes Deus ad conseruationem humani generis creauerit: vtcunque multi degenerent a diuina ordinatione: non tamen cessare proprerea debemus, quin et amemus quod Dei est, et saluum cupiamus. Haec causa est, cur debeant fideles, in quacunque regione degant, non modo legibus et magistratuum imperio parêre:sed in suis etiam precibus eorum salutem commenda re Deo. Dicebat Israelitis Ieremias: Orate pro pace Babylonis: quoniam in pace eius pax vestra. Haec vniuersalis est doctrina, vt ordinatas a Deo potestates cupiamus saluas et pacatas stare.  \pend\pstart Vt tranquillam.) Proposita vtilitate, stimulum nobis addit. Fructus enim enumerat, qui ex principatu rite composito nobis proueniunt. Primus est, tranquilla vita. Gladio enim sunt armati magistratus, vt nos in pace contineant. Nisi compescerent improborum audaciam, latrociniis et caedibus omnia essent plena. Haec igitur pacis tuendae ratio, cum vnicuique redditur quod suum est, nec impune grassatur potentiorum violentia. Secundus est, pietatis conseruatio: dum scilicet incumbunt magistratus, ad fouendam religionem, ad asserendum Dei cultum, ad sacrorum reuerentiam exigendam. Tertius est, cura publicae honestatis. Nam hoc etiam fit magistratuum beneficio, ne ad beluinas foeditates se homines prostituant, aut indecore lasciuiant, sed porius vigeat modestia et moderatio. Haec tria si tollantur qualis erit humanae vitae status? Si qua ergo vel publicae tranquillitatis, vel pietatis, vel honestatis cura nos tangit: meminerimus eorum quoque habendam esse curam, quorum ministerio tam praeclara bona ad nos perueniunt. Vnde colligimus, omnis humanitatis expertes essefanaticos homines, nec aliud quam feram barbariem spirare, qui magistratus e medio sublatos cupiunt. Quantûm enim haec discrepant: vt vigeat ius et honestum, vt floreat religio, orandum esse pro Regibus: et non modo regni nomen, sed totam politiam religioni esse aduersam? Atqui prioris sententiae authorem habemus Spiritum Dei. Secundam ergo a Diabolo esse oportet, Si roget quispiam, an etiam pro Regibus orandum sit, a quibus nibil tale percipimus? Respondeo, huc tendere vota nostra, vt Spiritu Dei gubernati, bonorum, quibus antea nos priuabant, mini\pend
\section*{CAP. II. }
\marginpar{[ p.19 ]}\pstart stri nobis esse incipiant. Non ergo pro iis tantd̀m, qui iam digni sunt, orare nos decet: sed rogandus Deus, vt ex malis bonos faciat. Tenendum enim semper istud principium: tam ad religionis, quam ad tranquillitatis et honestatis publicae custodiam destinatos esse a Deo magistratus: non aliter quam terra procreandis alimentis ordinata est. Quemadmodum itaque pro pane quotidiano orantes, Deum rogamus, vt terram sua benedictione foecundet: fic in illis prioribus beneficiis ordinarium medium, quod sua prouidentia constituit, respicere debemus. Huc accedit, quod si beneficiis istis, quorum dispensationem magistratibus Paulus assignat, destituimur: id contingit nostra culpa. Inutiles enim magistratus nobis reddit ira Dei, non secus, ac terram sterilem. Quare nostrûm est, eiusmodi flagella deprecari, quae peccatis nostris irrogantur. Praeterea hic officii sui vicissim admonentur principes, et quicunque magistratum gerunt. Neque enim satis est, si ius cuique suum reddendo, iniurias omnes coerceant pacemque foueant:nisi et religionem promouere, et honesta disciplina mores componere studeant. Neque enim frustra hortatur Dauid, vt Filium osculentur nec frustra Iesaias denunciat, fore Ecclesiae nutricios. Quare non est quod sibi blandiantur, si ad cultum Dei asserendum adiutores se praebere neglexerint.  \pend\pstart Hoc enim bonum est.) Postquam vtile esse docuit quod praecipiebat, iam validius argumentum proponit: Deo ita placere. Nam vbi constat de eius voluntate, instar omnium rationum esse nobis debet. Bonum accipit, pro recto et legitimo. Quia autem Dei voluntas regula est, ad quam exigenda sunt officia nostra omnia: rectum inde esse probat quia Deo acceptum est. Sequitur deinde huius etiam secundi membri confirmatio quia velit Deus omnes homines saluos facere. Quid autem magis aequum, quam vt huic Dei decreto vota nostra subseruiant? Postremd autem, Deo cordi esse omnium salutem demonstrat, quia omnes ad veritatis suae agnitionem vocet. Argumentum est a posteriori. Nam si Dei potentia est Euangelium in salutem omni credenti: certum est inuitari omnes ad spem vitae aeternae, quibus Euangelium offertur. Denique sieut vocatio documentum est arcanae electionis: ita quos facit Deus Euangelii sui panticipes, eos ad salutis possessionem admittit. Quia Euangelium iustitiam Dei nobis reuelat: quae certus est in vitam ingressus. Hinc apparet, quam pueriliter hallucinentur, qui locum hunc praedestinationi opponunt. Si Deus, inquiunt, omnes indifferenter vult saluos fieri: falsum est, alios ad salutem, alios ad exitium aeterno eius consilio esse praedestinatos. Aliquid forte dicerent, Si Paulus hic de singulis homi nibus ageret. Quanquam tunc quoque non deesset solutio. Nam etsi Dei voluntas non ex occultis ipsius iudiciis aestimanda est, vbi exter  \pend
\section*{I. TIMOTH. }
\marginpar{[ p.80 ]}\pstart nis signis eam nobis patefacit: non tamen propterea sequitur, quin constitutum intus habeat, quid de singulis hominibus fieri velit. Caeterùm id, quia nihil ad praesentem locum facit, praetereo. Nam Apostolus simpliciter intelligit, nullum mundi ordinem a salute excludi: quia omnibus sine exceptione Euangelium proponi Deus velit. Est autem Euangelii praedicatio viuifica. Merito itaque colligit, Deum omnes pariter salutis participatione dignari. At de hominum generibus, non singulis personis sermo est. Nihil enim aliud intendit. quam principes in hoc numero includere. Quod autem his quoque Euangelii doctrinam communem esse velit Deus, constat ex testimoniis iam citatis, et similibus. Neque enim de nihilo dictum est: Nunc Reges intelligite. In summa, indicare Paulus voluit, non esse considerandum, quales tunc essent principes: sed quales esse Deus vellet. Est autem officium charitatis, quoscunque vocatione sua complectitur Deus, eorum salutis curam studiumque suscipere: idque piis votis testari. Atque huc pertinet, quod Deum vocat nostrum saluatorem. Vnde enim nobis salus, nisi ex gratuita Dei beneficentia? Atqui Deus idem qui iam nos salutis fecit compotes, ad eos quoque gratiam suam extendere aliquando poterit. Qui iam nos ad se traxit, poterit illos nobiscum adducere. Fadurumi autem pro confesso sumit: quia Prophetarum vaticiniis ita praedictum erat, de vniuersis tam ordinibus, quam nationibus.  \pend
\phantomsection
\addcontentsline{toc}{subsection}{\textit{Vnus enim Deus, vnus et mediator Dei et hominum, homo Christus Iesus, qui dedit semetipsum precium redemptionis pro omnibus, vt esset testimonium temporibus suis, in quod positus sum praeco et Apostolus. Veritatem dico in Christo, non mentior, Doctor Gentium in fide et veritate. }}
\subsection*{\textit{Vnus enim Deus, vnus et mediator Dei et hominum, homo Christus Iesus, qui dedit semetipsum precium redemptionis pro omnibus, vt esset testimonium temporibus suis, in quod positus sum praeco et Apostolus. Veritatem dico in Christo, non mentior, Doctor Gentium in fide et veritate. }}\pstart Vnus Deus.) Parum valida in speciem foret haec ratio: Deum ideo velle omnes saluos esse, quia vnus est. Chrysostomus, et post eum alii sic accipiunt: non esse plures Deos, sicuti idololatrae com miniscuntur. Sed aliud fuisse Pauli consilium arbitror: vt hic sit tacita comparatio vnius Dei cum toto mundo diuersisque nationibus. Quemadmodum et tertio capite ad Romanos: Nunquid Iudaeorum Deus tantùm? An non et Gentium? Imo vnus Deus, qui iustificat Circuncisionem ex fide, et praeputium per fidem. Ergo vtcunque in hominibus tunc esses diuersitas, quia multi ordines alieni erant a fide, multae que gentes: ad vnitatem Dei fideles reuocat Paulus: vt  \pend
\section*{CAP. II. }
\marginpar{[ p.21 ]}\pstart sciant cum omnibus sibi esse coniunctionem, quia vnus sit omnium Deus: vt sciant non in perpetuum a spe salutis exclusos, qui sub eius dem Dei sunt potestate. Idem sibi vult quod de vno mediatore continuo subiicit. Nam sicuti vnus est Deus omnium creator et pater: ita vnicum mediatorem esse dicit, per quem accessus nobis ad Deum patet. Atque hunc mediatorem non vni tantùm genti vel paucis hominibus certae conditionis esse datum, sed omnibus. Nam sacrificii, quo peccata expiauit, fructum ad omnes pertinere. Particula vniuersalis semper ad hominum genera referri debet non ad personas. Acsidixisset: Non solos Iudaeos, sed Gentiles quoque: non solos plebeios, sed etiam principes redemptos esse morte Christi. Cum itaque commune mortis suae beneficium omnibus esse velit: iniuriam illi faciunt, qui opinione sua quempiam arcent a spe salutis.  \pend\pstart Homo Iesus.) Hominem dum praedicat, non negat esse Deum: sed cum vinculum notare vellet nostrae cum Deo coniunctionis, humanae potius naturae meminit, quam diuinae. Quod sedulo obseruandum est. Hinc enim factum est ab initio, vt homines sibi hos vel illos fingendo mediatores, longius a Deo recesserint, quia hoc errore praeoccupati, Deum procul abesse, quo se verterent nesciebant. Huic malo Paulus medetur, cum nobis Deum quasi praesentem sistit: quia ad nos vsque descendit, ne supra nubes quaerendus esset. Idem ergo hic dicitur, quod quarto ad Hebraeos capite: Non habemus Pontificeni, qui compati infirmitatibus nostris nequeat: cum tentatus fuerit per omnia. Et sane si id omnium animis insideret: porrigi nobis fraternam manum a filio Dei, et naturae societate nobis coniunctum, vt nos ex hac nostra tam abiecta conditione in coelum vsque attollat, quis non recta hanc viam tenere mallet, quam in deuiis salebris vagari? Proinde quoties orandus est Deus, si in mentem venit sublimis illa et inaccessa maiestas: ne eius formidine absterreamur, simul occurrat etiam homo Christus qui comiter nos inuitat, nósque veluti manu prehensat, qui Patrem ex formidabili ac tremendo propitium et facilem nobis reddat. Haec sola clauis est, qua nobis ianua coelestis Regni reseratur, vt appareamus cum fiducia in Dei conspectum. Proinde etiam videmus Sathanam omnibus seculis hunc lapidem voluisse, vt inde homines auerteret. Omitto quam variis artibus, ante Christi aduentum, distraxerit hominum mentes, ad media comminiscenda, quibus ad Deum peruenirent. Iam ab Ecclesiae Christianae exordio, cum nuper apparuisset Christus cum tam praeclaro pignore: et adhuc in terris ex eius ore suauissima illa vox propemodum resonaret: Venite ad me omnes, et c.erant tamen quidam fallendi artifices, qui Angelos pro mediatoribus obtruderent ipsius loco: quemadmodum ex Epistola ad Colossenses colligere promptum est.  \pend
\section*{I. TIMOTH. }
\marginpar{[ p.22 ]}\pstart Sed quod tunc clanculum moliebatur Sathan, sub Papatu ed̀ vsque perduxit, vt vix millesimus quisque mediatorem Christum, vel titulo tenus agnosceret. Atque ita sepultum erat nomen, vt res magis esset ignota. Nunc postquam sanos, piosque Doctores excitauit Deus, qui tanquam postliminio reuocare in hominum memoriam studuerunt, quod vnum ex notissimis fidei nostrae principiis esse debuerat: omnia excogitant Romanenses sophistae, quibus rem clarissimam obscurent. Primùm ita illis nomen est inuisum, vt si quispiam Christi mediatoris faciat mentionem praeteritis sanctis, mox suspitione grauerur haereseos. Quia autem repudiare in totum non audent quod hic a Paulo docetur, insulso comniento eludunt: vnum vocari, non solum. Quasi vnum ex magna turba Deum nominauerit. Cohaerent enim haec duo membra: Vnum esse Deum, et vnum mediatorem. Quare qui Christum vnum aliquem ex multis faciunt, idem ad Deum quoque trahant oportet. An eo prorumperent impudentiae, nisi caecus furor ad opprimendam Christi gloriam eos raptaret? Alii sibi videntur acutiores, dum Christum vnicum statuunt mediatorem redemptionis, sanctos autem intercessionis mediatores nominant. Atqui horum quoque insulsitatem coarguit loci circunstantia quandoquidem hic de precibus ex profe sso tractatur. Pro omnibus, inquam, orare Spiritus praecipit, quia vnicus noster mediator omnes ad se admittat: sicuti morte sua omnes reconciliauit Patri. Et adhuc Christiani censeri volunt, qui tam sacrilega audacia Christum spoliant suo honore. Obiicitur tamen quod speciem repugnantiae habeat. Nam hoc ipso loco Paulus iubet nos pro aliis intercedere: quod octauo ad Romanos, Christo quasi proprium assignat. Respondeo, sanctorum intercessores, quibus se mutuo iuuant apud Deum, minime obstare, quominus tamen vnicum omnes intercessorem habeant. Nemo enim vel pro se, vel pro alio, nisi Christi parrocinio fretus exauditur. Quod itaque alii pro aliis intercedimus, adeo non tollit vnicam Christi intercessionem, vt inde maxime pendeat, côque referatur. Putabit quispiam, nos igitur posse facile conciliari cum Papistis, si vnicae Christi intercessioni subiiciant, quascunque ipsi fanctis tribuunt. Non ita est. Nam ideo tranfferunt intercedendi officium ad sanctos, quoniam alias destitui nos aduocato fingunt. Hoc est vulgare inter ipsos principium: nos patronis opus habere, quia indigni sumus, qui per nos appareamus in Dei conspectum. Ita loquendo, Christum honore suo priuant. Deinde execranda est blasphemia, affingere sanctis dignitatem, quae nobis apud Deum gratiam conciliet. A quo tani longe absunt, tam Prophetae omnes quam Apostoli et Martyres, vsque ad ipsos Angelos, vt patrono ipsi quoque eodem nobiscum opus habeant. Deinde merum est figmentum.  \pend
\section*{CAP. II. }
\marginpar{[ p.23 ]}\pstart natum in eorum cerebris, mortuos pro nobis intercedere. Ergo in eo fundare preces nostras, est fidem penitus a Dei inuocatione auellere. Atqui Paulus regulam Dei rite inuocandi praescribit, fidem ex verbo Dei conceptam Ro.1o. Meritd itaque repudiamus, quod curiosi homines extra Dei verbum imaginantur. Sed ne prolixior sim, quam loci expositio flagitat: summa sit, Christo vno fore contentos, qui vere didicerint eius officium: mediatores sibi propria libidine fabricare eorum esse, qui nec Deum nec Christum nouerunt. Vnde consti ruo, Papistarum doctrinam, quae et Christi patrocinium obscurat, imo fere sepelit, et fictitios inducit Patronos absque vsso scripturae testimonio, tum impiae diffidentiae, tum etiam impiae temetitatis plenam esse.  \pend\pstart Qui dedit semetipsum.) Non est superuacua hoc loco redemptionis commemoratio. Sunt enim res necessario coniunctae, sacrificium mortis Christi, et perpetua intercessio: suntque duae sacerdotii partes. Nam hac significatione sacerdos noster Christus vocatur, quod semel peccata nostra morte sua expiauit, vt Deum nobis propitium redderet, et nunc sanctuarium coeli ingressus, ad impetrandam nobis gratiam coram Patre apparet, vt eius nomine exaudiamur. Quo melius detegitur sacrilega Papistarum impietas, qui dum in hoc officio sanctos Christo socios adiungunt, simul ad eos transferunt sacerdotii gloriam. Lege quartum caput ad Hebraeos circa finem et initium quinti. Reperies quod dico, intercessionem, qua propitiatur nobis Deus in sacrificio fundatam esse. Quinetiam tota veteris sacerdotii ratio hoc demonstrat. Sequitur ergo nihil ex officio intercedendi a Christo transeribi posse ad alios, quin sacerdotii titulo nudetur. Praeterea cum αντιλυτρον appellat, alias omnes satisfactiones eueritit. Quanquam non me latet Papistarum argutia: fingunt enini redemptionis preciumi, quod morte sua persoluit Christus, in Baptismo nobis applicari, vt deleatur originale peccatum: postea sarisfactionibus nos reconciliari Deo. Hoc modo, quod erat vniuersale et perpetuum beneficium, ad exiguum tempus et speciem vnam restringunt. Sed plenam huius rei tractationem ex Institutione pete.  \pend\pstart Vt esset testimonium.) Hoc est, vt haec gratia patefieret tempore constituto. Dixerat, pro omnibus. Quae parcicula quaestionem mouere poterat:cur ergo peculiarem vnum populum elegisset Deus, si omnibus communiter se propitium Patrem exhiberet, ac vna omnibus communis esset redempcio in Christo. Huic quaestioni ansam praecidit, cum reuelandae gratiae opportunitatem refert ad Dei consilium. Nam si hyeme non miramur emarcidas arbores, agros niue coopertos, prata gelu rigentia: vere aucem reuirescere.  \pend
\section*{I. TIMOTH. }
\marginpar{[ p.24 ]}\pstart quae ad tempus quodammodo emortua fuerant, quia ordinatae sunt a Deo temporum vices: cur non ipsius prouidentiae in aliis idem iuris concedemus? An propterea Deum inconstantiae arguemus, quia quod semper apud se decretumfixum est habuit, suo tempore profert in me dium? Quare et si hoc repentinum mundo accidit et minime expectatum, quod Christus Iudaeis ac Gentibus promiscue redemptor apparuit: ne tamen putemus subitum hoc fuisse Deo: sed potius admirabili eius prouidentiae discamus subiicere sensus omnes nostros. Ita fiet, vt nihil quod ab ipso prodierit nobis non sit maxime tempestiuum. Ideo frequenter apud Paulum occurrit haec admonitio: praesertim eum de Gentium vocatione agitur, quae sua nouitate multos tunc percellebat, et quasi reddebat attonitos. Quibus hac solutione non satisfit, Deum occulta sua sapientia distribuisse temporum vices: hi aliquando experientur, quo tempore ociosum fuisse putant, fabricasse inferos curiosis.  \pend\pstart In quod positus sum.) Ne temere, vt multi solent, de re sibi parum comperta asserere videatur: se ad hoc praedicat diuinitus ordinatum, vt Gentes prius a regno Dei alienas ad Euangelii participationem adducat. Nam Apostolatus eius ad Gentes, certum erat diuinae vocationis fundamentum. Et ideo tantopere in ipso asserendo laborat: vt certe apud multos non absque difficultate recipiebatur. Ius iurandum vel obtestationem adhibet, vt in re maxime seria, doctorem se esse Gentium: idque in fide et veritate. Quae duo bonam conscientiam significant. Sed eam quoque fultam esse oportet certitudine diuinae voluntatis. Proinde non tantùm sincero affectu se Gentibus Euangelium praedicare intelligit: sed etiam sana et intrepida conscientia: eo quod nihil agat nisi Dei mandato.  \pend
\phantomsection
\addcontentsline{toc}{subsection}{\textit{Volo igitur orare viros in omni loco, sustollentes puras manus, absque ira, et disceptatione. Consimiliter et mulieres in amictu modesto cum verecundia et temperantia ornare semetipsas, non tortis crinibus, aut auro, aut Margaritis, aut vestitu sumptuoso: sed quod decet mulieres profitentes pietatem per bona opera. }}
\subsection*{\textit{Volo igitur orare viros in omni loco, sustollentes puras manus, absque ira, et disceptatione. Consimiliter et mulieres in amictu modesto cum verecundia et temperantia ornare semetipsas, non tortis crinibus, aut auro, aut Margaritis, aut vestitu sumptuoso: sed quod decet mulieres profitentes pietatem per bona opera. }}\pstart Volo igitur.) Pendet haecillatio ex proxima sententia. Nam, vt in Epistola ad Galatas vidimus, Spiritu adoptionis praeditos esse nos oportet, vt Deum rite inuocemus. Postquam ergo Christi gratiam omnibus exposuit, ideôque se Gentibus datum esse  \pend
\section*{CAP. II. }
\marginpar{[ p.25 ]}\pstart Apostolum meminit, vt eodem redemptionis beneficio promiscue cum Iudaeis fruerentur: omnes pariter ad orandum inuitat. Fides enim inuocationem parit. Vnde et ad Romanos 15. Gentium vocationem probat istis testimoniis: Exultent Gentes cum populo eius. Item: Omnes Gentes laudate Dominum. Item Confitebor tibi inter Gentes. Valet enim reciprocum argumentum a fide ad inuocationem, et ab inuocatione ad fidem: tanquam a causa ad effectum, et ab effectu ad causam. Quod obseruatione dignum est: quia inde admonemur, Deum se nobis patefacere verbo suo, vt a nobis in uocetur: et hoc praecipuum esse exercitium fidei. Quare particula, in omni loco, peraeque hic valet, atque initio prioris ad Corinthios: vt nullum sit iam discrimen inter Gentilem et Iudaeum, inter Graecum et Barbarum: quia Deus communis est omnium Pater.  \pend\pstart sustollentes puras manus.) Acsi diceret: modo adsit bona conscientia, nihil obstare, quominus vbique inuocent Deum omnes Gentes. Sed rei loco signum posuit. Nam purae manus sunt puri cordis indices. Quemadmodum econuerso Iesaias exprobrat Iudaeis, quod manus ad Deum sanguinolentas extollant, dum inuehitur in eorum crudelitatem. Porrô ceremonia haec vsitata fuit omnibus seculis: quia natura hic nobis ingenitus est sensus, cum Deum quaerimus, sursum aspicere. Adeóque semper valuit, vt ipsi quoque idololatrae, cum alioqui Deum affigerent ligneis et lapideis simulacris, morem tamen in coelum tollendi manus retinerent. Discamus ergo, ceremoniam esse verae pietati congruentem, modo respondeat, quae per eam figuratur veritas. Nempe, vt admoniti Dumin caelo quaerendum esse, primùm de ipso terrenum vel carnale nihil imaginemur deinde exuamus aftectus carnales, ne quid impediat, quominus animi nostri supra mundum assurgant. Idololatrae vero et hypocritae, manus inter oranduni attollendo, simiae sunt: quia dum externo symbolo profitentur, se mentes sursum habere erectas, priores in ligno et lapide haerent, quasi illic inclusus esset Deus: secundi aut inanibus curis, aut prauis cogitationibus impliciti in terra iacent. Proinde contrario gestu testimonium aduersum se ferunt.  \pend\pstart Absque ira.) Nonnulli fremitum indignationis exponunt, dum secum tumultuatur conscientia, et quasi cum Deo expostulat. Quod acidere solet, cum res aduersae grauius nos premunt. Tunc enim aegreferimus, non statim Deum nobis succurrere, etimpatientia turbamur. Concutitur etiam fides variis insultibus. Nam quia non apparet eius auxilium, obrepunt nobis dubitationes, an curam nostri habeat, an saluos velit nec ne: et similes. Talem consternationem haesitantis animi disceptationis voce notari putant. Ita secundùm eos sensus esset: pacata conscientia et intrepida fiducia  \pend
\section*{I. TIMOTH. }
\marginpar{[ p.26 ]}\pstart orandum esse. Alii, vt Chrysostomus, requiri putant, non tam ergo Deum quam homines placidos animos et omni perturbatione vacuos. Quia nihil est quod puram Dei inuocationem magis impediat, quam rixae et contentiones. Vnde etiam iubet Christus si quis dissideat cum fratre, prius reconciliari, quam munus ad altare offerat. Horum vtrunque verissimum esse fateor. Sed dum praesentis loci cireunstatiam expendo, non dubito quin respexerit Paulus ad disceptationes, quae inde oriebantur, quod Iudaei Gentes sibi aequari indignabantur, et ideo controuersiam mouebant de earum vocatione, imo eas repudiabant, arcebantque a gratiae confortio. Vult ergo Paulus, sedatis eiusmodi contentionibus, omnes vbiuis gentium ac terrarum Dei filios vnanimes orare. Quanquam nihil prohibet, quin ex particulari sententia eliciamus generalem doctrinam.  \pend\pstart Consimiliter et mulieres.) Quemadmodum viros iussit tollere puras manus, ita nunc praescribit, qualiter mulieres ad rite orandum compararas esse deceat. Ac videtur inter virtutes quas commendat, et externani Iudaeorum sanctificationem tacita subesse antithesis. Nam innuit, nullum esse profanum locum, et unde non ad Deum accessus pateat, modo suis vitiis non arceantur, tam viri quam mulieres. Porrd̀ sumpta occasione, vitium, quo fere laborare solent mulieres, corrigere voluit. Fieri quoque potest, vt Ephesi, in vrbe opulenta et celebri emporio, magis grassantum sit. Est autem nimia ornatus cura etcupiditas. Vult autem vt ad pudorem ettemperantiam compositus sit earum ornatus. Nam inde luxuria et immodicus sumptus, vel quod superbiae, vel quod lasciuiae causa, ostentare se appetunt. Atque inde petenda est regula mediocritatis. Nam quia res est indifferens, vestitus vt sunt res omnes externae. diffieile est praescribere certum modum, quousque liceat. Magistratus quidem leges sumptuarias ferre poterit, quibus libidinem vtcunque coerceat. Sed pii doctores, quorum est regere conscientias, semper debent vsus legitimi finem respicere. Ita hoc saltem erit extra controuersiam quidquid non est cum pudore et temperantia consentaneum in vestitu, esse improbandum. Quanquam ab affectu semper est incipiendum. Nam vbi intus regnat proteruia, illic nullus pudor. Vbi regnat intus ambitio, illic nulla modestia in externo vestitu extabit. Sed quia hypocritae vitiosos suos affectus quibuscunque possunt coloribus fucare folent, necessario protrahendi sunt ad id quod apparet. Magnae improbitatis erit negare, honestis et castis mulieribus pudorem conuenire, tanquam proprium et pe rpetuum ornamentum. Similiter, moderationem ab omnibus seruandam esse. Quid quid ergo erit his contrarium, frustra excusabunt. Certas deinde species excessus nominatim reprehendit: vt sunt capilli complexi,  \pend
\section*{CAP. II. }
\marginpar{[ p.27 ]}\pstart Margaritae, et aurea monilia: non quod vetitus sit in totum auri vel Margaritarum vsus, sed quod haec vbicumque refulgent secum, fere trahunt alia mala quae dixi, et ex ambitionis fonte nascuntur aut impudiciuae.  \pend\pstart Quod decet mulieres.) Nam certe cultum honestae et piae mulie ris a meretricio differre oportet. Illae autem sunt discretionis notae, quas posuit. Et si operibus testanda est pietas, in vestitu etiam castoapparere haec professio debet.  \pend
\phantomsection
\addcontentsline{toc}{subsection}{\textit{Mulier in quiete discat, cum omni subiectione. Docere autem mulieri non permitto, neque imperiosam esse in virum, sed quietam esse. Adam enim creatus fuit prior, deinde Eua. Et Adam non fuit deceptus, sed mulier decepta, transgressionis rea fuit. Seruabitur autem per generationem, si manserint in fide, et charitate, et sanctificatione, cum temperantia. }}
\subsection*{\textit{Mulier in quiete discat, cum omni subiectione. Docere autem mulieri non permitto, neque imperiosam esse in virum, sed quietam esse. Adam enim creatus fuit prior, deinde Eua. Et Adam non fuit deceptus, sed mulier decepta, transgressionis rea fuit. Seruabitur autem per generationem, si manserint in fide, et charitate, et sanctificatione, cum temperantia. }}\pstart Postquam de vestitu locutus est, nunc addit, qua se modestia gerere in sacro coetu debeant mulieres. Ac primum iubet eas placide discere. Nam quies silentium significat, ne loquendi vices sibi vsurpent in publico. Quod statim clarius exponit, dum eas vetat docere. Non vt illis auferat munus instituendae familiae: sed tantùm vt a munere docendi arceat, quod solis viris Deus mandauit, Qua de re tractauimus in epistola ad Corinthios priore. Si quis Debotam et similes obiiciat, quas aliquando gubernando populo Dei mandato praefectas legimus: responsio est facilis: extraordinariis Dei factis non euerti communem politiam, cui nos voluit alligatos . Ergo si aliquando prophetandi et docendi locum tenuerunt mulieres, idque Dei Spiritu excitatae: potuit hoc, qui ab omni lege immunis est. Sed quia hoc singulare est, non pugnat cum perpetua et vsitata politia. Addit quod proximum est officio docendi: ne sibi authoritatem sumant in viros. Nam ista est ratio,eur docere prohibeantur: quia hoc non patitur earum conditio. Subiectae enim sunt docere autem est potestatis, vel superioris gradus. Quanquam videtur ratio non satis firma:cum Prophetae quoque et doctores regibus sint subiecti, aliisque magistratibus. Respondeo, nihil esse absurdi, quin praesit aliquis, et simul pareat, secundùm diuersos repectus. Sed in muliere id non valet, quae natura, hoc est, ordinaria Dei lege, ad parendum nata est. Nam γυναικοκρατιαν omnes prudentes semper instar portenti repudiarunt. Quare, coelum quodammodo terrae miscebitur, si docendi ius arripiant mulieres. Iubet ergo eas quiescere, hoc est,se continere intra suum ordinem.  \pend\pstart Adam enim creatus.) Videtur haec parum valida esse subie\pend
\section*{I. TIMOTH. }
\marginpar{[ p.28 ]}\pstart ctionis ratio. Nam et Ioannes Baptista tempore Christum praecessit, dignitate tamen longe inferior. Sed Paulus circunstantias omnes, quae a Moise referuntur, tametsi non exprimat, voluit tamen a lectoribus expendi. Moises autem docet, ita posteriore loco creatam esse mulierem, vt sit quasi viri accessio: et hac lege fuisse viro adiunctam, vt praesto adsit, ad exhibenda obsequia. Cum ergo non duo capita Deus aequa potestate creauerit: sed viro addiderit adiumentum inferius: merito ad illum creationis ordinem nos Apostolus reuocat, in quo relucet aeterna et inuiolabilis Dei institutio.  \pend\pstart Et Adam non fuit deceptus.) Respicit ad poenam mulieri inflictam: Quia obedisti voci serpentis: sub potestate viri eris, et ad illum appetitus tuus. Nam quia exitiale consilium dederat, digna fuit, quae ab alieno iure et arbitrio pendêre disceret: quia virum abduxerat a Dei imperio, digna fuit quae omni libertate priuata, sub iugum redigeretur. Caeterdm non simplicem aut nudam transgressionis causam hic Apostolus obtendit: sed pronunciata a Deo sententia nititur. Sed videntur haec duo nonnihil pugnare: quod mulieris sub iectio poena sit transgressionis: et tamen illi fuerit a creatione imposita. Nam inde sequitur, prius addictam fuisse seruituti, quam delinqueret. Respondeo, nihil impedire, quominus et naturalis ab initio fuerit conditio seruiendi: et deinde propter peccatum accidentalis esse coeperit: vt iam minus liberalis sit subiectio, quam prius fuisset. Caeterùm locus hic quibusdam occasionem praebuit, vt negarent Adam errore lapsum esse:sed vxoris tantùm illecebris fuisse victum. Putant igitur mulierem tantûm astu diaboli deceptam, vt se et virum diis similes fore crederet. Adae vero nihil fuisse persuasum: sed frustum gustasse, vt morem vxori gereret. Atqui hanc opinionem promptum est refellere. Nisi enini fidem habuisset sathanae mendacio: non illi exprobraret Deus: Ecce Adam quasi vnus ex nobis. Sunt aliae rationes quas taceo: quia non longa refutatione indiget error, nulla probabili coniectura fultus. Paulus enim suis verbis non significat Adam non eadem implicitum fuisse Diaboli fallacia: sed tantum causam vel originem transgressionis ab Eua profectam esse.  \pend\pstart Seruabitur autem.) Quia imbecillitas sexus mulieres reddit magis suspitiosas et timidas: et talis erat superior sententia, quae vehementer percellere et consternare vel maxime viriles animos posset: ideo adhibita consolatione, temperat quod dixit. Neque enim nos ideo accusat vel coarguit Spiritus Dei, vt pudore confusis insultet: sed prostratos mox erigit. Poterat hoc, vt iam dixi, exanimare mulieres: eum sibi imputari audirent totius humani generis exitium. Qualis enim est hic reatus? Prasertim cum in subiectione  \pend
\section*{CAP. II. }
\marginpar{[ p.29 ]}\pstart ira Dei assidue versetur ob oculos. Ergo Paulus, vt illas soletur,ac suam illis conditionem tolerabilem reddat, spem salutis relictam esse admonet, vteunque sustineant poenam temporalem. Duplex iste consolationis fructus notandus est: quod proposita salutis spe retinentur: ne reatus sui commemoratione perterritae, in desperationem ruant: deinde quod assuefiant ad tolerandam aequis et tranquillis animis seruiendi necessitatem, vt libenter se viris submittant, cum admonentur hoc genus obsequii etsibi esse salutare et Deo acceptum. Si ad operum adstruendam iustitiam hic locus torqueatur, vt faciunt Papistae:simplex est solutio: quia hic non disputat Apostolus de causa salutis, non posse nec debere ex eius verbis colligi, quid mereantut opera: sed tantùm ostendi, qua nos via perducat Deus ad salutem cui nos sua gratia destinauit.  \pend\pstart Per generationem.) Ridiculum hoc foret hominibus nasutis, Christi Apostolum non modo hortari mulieres ad dandam operam procreandae soboli: sed hoc opus tanquam pium et sanctum ita vrgere, vt rationem esse dicat obtinendae salutis. Quinetiam videmus, quibus probris tlorum coniugalem infamauerint hypocritae, qui sanctiores videri prae aliis omnibus volebant. Sed aduersus sannas impiorum facilis est defensio. Primùm enim non tantùm de gignenda prole hic tractat Apostolus: sed de perferendis molestiis omnibus, quae et durae sunt etmultiplices tam in partu quam in educatione liberorum. Deinde quidquid hypocritae iudicent vel mundi sapientes, pluris hanc obedientiam facit Deus, cum mulier ad quid vocata sit reputans, impositae sibi a Deo conditioni se subiicit, nec taedium gestationis, nec ciborum fastidium, nec morbos, nec difficultatem pariendi, imo potius dirum cruciatum, nec solicitudinem pro foetu, nec alia quae sunt officii sui detrectat, quam si haeroicis alioqui virtutibus se ostentaret, Dei interim vocationi obsequi detrectans. Adde, quod non potuit aptior afferri consolatio, nec efficacior, quam si rationem et medium, vt ita loquar, salutis consequendae in ipsa poena ostenderet.  \pend\pstart Si manserint in fide.) Quia vetus interpres generationem liberorum posuerat: vulgo receptum fuit, vt ad liberos referrent hoc membrum. Atqui vnica vox est apud Paulum, τεκνογονία. Proinde ad mulieres referri necessarium est. Quod autem plurale verbum est, nomen vero singulare, nihil habet incommodi. Si quidem nomen indefinitum, vbi scilicet de omnibus comimunis est sermo, vim collectiui habet: ideóque mutationem numeri facile patitur. Porro, ne totam mulierum virtutem in coniugalibus officiis includeret:ideo maiora adducit. Imo tunc demum generatio gratum est Deo obsequium,cum ex fide et charitate procedit. His duo\pend
\section*{I. TIMOTH. }
\marginpar{[ p.30 ]}\pstart bus adiungit sandificationem, quae totam vitae puritatem continet, dignam Christianis mulieribus. Postremo sequitur temperantia, cuius pauloante meminerat, de vestitu loquens. Sed nunc eam latius ad alias vitae partes extendit.  \pend
\endnumbering\beginnumbering\section{CAPVT. III.}
\phantomsection
\addcontentsline{toc}{subsection}{\textit{\huge\textbf{C}\normalsize ERTVS sermo: Si quis episcopatum appetit, praeclarum opus desiderat. Oportet ergo episcopum irreprehensibilem esse, vnius vxoris maritum, sobrium, temperantem, compositum, hospitalem, aptum ad docendum, non vinolentum, non percussorem, non turpiter lucri cupidum, sed aequum, alienum a pugnis, alienum ab auaritia, qui domui suae bene praesit, qui filios habeat in subiectione, cum omni reuerentia. Quod si quis propriae domui praeesse non nouit: EcCicsiam Dei quomodo curabit? Non nouitium, ne inflatus in condemnationem incidat Diaboli. Oportet autem illum et bonum testimonium habere ab extraneis:ne in probrum incidat et laqueum Diaboli. }}
\subsection*{\textit{\huge\textbf{C}\normalsize ERTVS sermo: Si quis episcopatum appetit, praeclarum opus desiderat. Oportet ergo episcopum irreprehensibilem esse, vnius vxoris maritum, sobrium, temperantem, compositum, hospitalem, aptum ad docendum, non vinolentum, non percussorem, non turpiter lucri cupidum, sed aequum, alienum a pugnis, alienum ab auaritia, qui domui suae bene praesit, qui filios habeat in subiectione, cum omni reuerentia. Quod si quis propriae domui praeesse non nouit: EcCicsiam Dei quomodo curabit? Non nouitium, ne inflatus in condemnationem incidat Diaboli. Oportet autem illum et bonum testimonium habere ab extraneis:ne in probrum incidat et laqueum Diaboli. }}\pstart Certus sermo.) Quod Chrysostomus clausulam esse putat superioris doctrinae, mihi nomplacet. Solet enim Paulus hanc formulam praefando magis vsurpare. Deinde non erat illic tantae affirmationis vsus. Quod autem nunc dicturus est, plus habet aliquanto ponderis. Sit igitur praefatio, ad notandam rei grauitatem  \pend\pstart habita. Si quis episcopatum.) Quoniam mulieribus interdixerat officium docendi sumpta inde occasione, nunc de episcopatu ipso disserit. Primo, quo melius appareat, quam non sine causa mulieres admittere noluerit ad ipsum exercendum: deinde, ne solas mulieres arcendo viros omnes promiscue admittere videatur: tertio quia expediebat, Timotheum et alios admoneri, quanta in eligendis episcopis religio seruanda esset. Talis ergo estcontextus, mea sententia, acsi diceret Paulus, adeo non esse idontas mulieres obeundae tantae functioni, vt ne viris quidem ipsis indiscriminatim patefieri aditum  \pend
\section*{CAPVT III. }
\marginpar{[ p.31 ]}\pstart oporteat. Proinde admonet, non vulgare esse opus, vt sibi quilibet permittat id suscipere. Nam cum dicit καλόν, non dubito quin ad vetus prouerbium alludat, quod saepius a Platone repetitur: Αυσκολα τα καλὰ. Quo Graeci significant, quae pulchra sunt, eadem ardua et difficilia esse: atque ita difficultatem cum excellentia coniungit. Imo sic argumentatur: non esse cuiusuis Episcopatu fungi, quia res sit eximia. Iam puto sensum Pauli esse satis perspicuum: quem tamen nemo interpretum assecutus est, quod videam. Summa est, habendum esse in recipiendis Episcopis delectum, quia munus sit laboriosum ac difficile: ac diligenter qui ad illud adspirant debere secum reputare, an pares futuri sint ferendo tanto oneri. Audax semper est imperitia: et circunspecta rerum cognitio modestum reddit hominem. Vnde fit, vt tanta confidentia saepe munus gubernandi ambiant qui nihil habent, vel ingenii, vel prudentiae, nisi quod clausis oculis ruunt? Sic dicebat Quintilianus, indoctos fortius dicere: cum trepident summi oratores. Vt in appetendo episcopatu eiusmodi temeritatem cohibeat Paulus, admonet primo, non ociosam esse dignitatem, sed opus. Deinde non quoduis opus, sed eximium, et ideo arduum, plenumque difficultatis: vt est reuera. Neque enim leue est. in tanta administratione sustinere personam filii Dei, vbi agitur de Regno Dei erigendo et propagando, vbi de animarum salute curanda, quas sanguinis sui precio Dominus ipse dignatus est, de regenda Ecclesia, quae est Dei haereditas. Sed non est hic propositum concionari. Et hoc argumentum rursus attinget Paulus proximo capite. Verum quaeritur, an vllo modo episcopatum appetere liceat? Videtur enim absurdum, vt quis voto praeueniat Dei vocationem. Atqui Pau lus, dum temerarium desiderium reprehendit, permittere videtur, vt circunspecte et modeste quis appetat. Respondeo, si aliis in rebus dam natur ambitio, in episcopatu multd grauius damnandam esse. Atqui Apostolus de pio desiderio loquitur, quo student sancti homines quidquid habent doctrinae conferre ad Ecclesiae aedificationem. Nam si docendi munus appetere omnino nefas esset: quorsum se discendo compararent qui totam adulescentiam in sacrae Scripturae lectione consumunt? An non theologicae scholae pastorum sunt seminariat Proinde qui ita sunt instituti, non modo licite possunt, sed etiam debent etiam antequam in munus sint cooptati, voluntaria oblatione se ac suam operam sacrare Deo. Modo tamen non se ingerant ipsi, ac ne voto quidem suo se designent Episcopos: sed tantùm parati sint ad munus obeundum, si eorum opera ftagitetur. Quod si legitimo ordine non fuerint vocati, sciant aliter visum esse Deo, neque indigne ferant se fuisse aliis posthabitos. Verùm ita affecti erunt, qui, nullo sui respectu. Deo et Ecclesiae tantùm seruire optabunt: et simul erunt  \pend
\section*{I. TIMOTH. }
\marginpar{[ p.32 ]}\pstart ea modestia, vt minime inuideant, si alii, tanquam digniores, praeferantur. Si obiiciat quispiam, tantae molis esse Ecclesiae regimen, vt horrorem sanis hominibus incutere potius debeat, quam prouocare ad sui desyderium: respondeo, piorum desyderium non recumbere in fiduciam propriae industriae, aut virtutis, sed Dei auxilio, a quo no stra est sufficientia:quemadmodum alibi tractat Paulus. Caeterùm, obseruandum est, quid Paulo sit episcopatus. Eóque magis, quod veteres, suorum temporum consuetudine a vero sensu abducti fuerunt. Nam cum Paulus generaliter comprehendat omnes pastores: ipsi Episcopum accipiunt, qui ex vnoquoque collegio eligebatur, vt fratribus praeesset. Meminerimus ergo perinde valere hoc nomen, acsi ministros, vel pastores, vel presbyteros nominasset.  \pend\pstart Oportet ergo.) Illatiua particula expositionem, quam posui, confirmat. Nam pro muneris dignitate colligit, hominem requiri praeditum raris dotibus, et non quemlibet ex vulgo. Si esset bonum opus, vt habet communis transsatio, vel honestum, vt reddidit Erasmus illatio non quadraret. Vult auteni Episcopum esse irreprehensibilem. Pro quo ad Titum posuit, ἀνέ γαλκτομ: vtroque nomine significans, non debere vlla infamia notatum esse, quae authoritatem illi deroget. Neque enim inter homines reperietur, qui purus sit omni vitio. Sed aliud est, laborare communibus vitiis, quae existimationem non laedunt, quia cadunt in homines honestissimos: aliud probrosi esse nominis, vel aliqua turpitudine esse foedatum. Ergo ne authoritate careat Episcopus, iubet eligi, qui sit honestae et integrae fa mae, nullique notabili vitio obnoxius. Porro non tantùm praescribit Timotheo, quales eligere debeat:sed etiam singulos admonet, qui ad gradum illum adspirant, vt se vitamque suam diligenter examinent.  \pend\pstart Vnius vxoris maritum.) Puerile est commentum, interpretari vnius Ecclesiae pastorem. Receptior est altera expositio, debere recipi, qui non fuerit duplici coniugio implicitus, iam tamen coelibem, vnica vxore mortua. Sed enim tam hic, quam primo ad Titum capite, verbaita sonant, qui sit non qui fuerit. Et infra capite quinto, cum aget de viduis, distincte vtetur praeteriti temporis participio. Prae terea, hoc modo pugnaret ipse secum: vtpote qui alibi testatur, se talem nolle iniicere laqueum conscientiis. Sola igitur Chrysostomi exposicio vera est, polygamiam hic nominatim in Episcopo damnari, quae tunc apud Iudaeos pro lege fere obtinuerat. Partim hoc habebant ex peruersa patrum imitatione. Nam qui legebant Abraham, Iacob, Dauidem et similes, pluribus simul vxoribus fuisse coniugatos idem sibi quoque licere arbitrabantur. Partim hanc corruptelam ex vicinis gentibus traxerant. Nunquam enim apud Orientales qua  \pend
\section*{CAPVT III. }
\marginpar{[ p.33 ]}\pstart decebat religione et fide culta fuerunt coniugia. Quidquid sit: polygamia apud eos yulgo inualuerat. Itaque non immeritd̀ Paulus maculam hanc ab Episcopi persona alienam esse iubet. Neque tamen eorum opinionem improbo, qui putant Spiritum sanctum Diabolicae superstitioni voluisse occurrere, quae postea obrepsit. Acsi dixisset, adeo non cogendos ad coelibatum Episcopos, vt maxime ipsos deceat coniugium. Hoc modo non exigeret ab illis tanquam rem necessariam: sed tantùm laudaret, tanquam ab officii dignitate minime alienam. Simplicius tamen est ac firmius, quod iam dixitab ordine episcopali arceri polygamiam a Paulo: quia signum sit hominis libidinosi et fidem coniugalem non feruantis. Verùm hic obiicitur, non debuisse, quod in omnibus est vitiosum, solis episcopis prohiberi. Solutio facilis est: non ideo protinus aliis fieri licentiam, cum episcopis nominatim hoc interdicitur. Neque enim dubium est, quin Paulus in vniuersum damnauerit, quod cum perpetua Dei lege pugnabat. Fixum enim est decretum: Erunt duo in carnem vnam. Sed potuit in aliis vtcunque tolerare, quod in episcopo niniis turpe, ideóque non ferendum fuisset. Neque enim hic lex in posterum sancitur, ne episcopus qui vnam habet vxorem, secundam aut tertiam inducat. Sed ab ordine episcopali repellit Paulus, quicunque tale flagitium admiserint. Quod itaque iam factum corrigi non poterat, fert coactus, sed non nisi in plebe. Quid enim remedii? An secundas ac tertias vxores repudiassent, qui sub Iudaismo in polygamiam se coniecerantt Atqui tale repudium non caruisset iniuria. Quod ergo integrum non erat, intactum reliquit. Tantùm excepit ne quis episcopus tali macula aspersus esset.  \pend\pstart Sobrium.) Erasmus reddidit, vigilantem. Quando νκc α̃- atas Graecis vtrunque significat:eligant vtrum volent lectores. σuερον α malui reddere temperantem, quam sobrium: quia συρ ροσυνκ latius patet, quam sobrietas. Compositus est, qui decenter et honeste se gerit. Hospitalitas est erga extraneos: cuius frequentior apud veteres vsus fuit. Turpe enim fuisset honestis hominibus, et praesertim notis, in cauponas diuertere. Hodie alia ratio.Est tamen, semperque erit virtus haec in Episcopo multas ob causas necessaria.  \pend\pstart Aptum ad docendum.) Doctrinae in Epistola ad Titum expressior fit mentio. Hic tantùm breuiter dexteritatem ponit. Neque enim satis est, si quis recondita eruditione polleat: nisi accedat etiam docendi gratia. Multi enim, vel quiaelinguae sunt impeditae, vel quia non satis perspicacis ingemi, vel quia non satis populares, intus suppressum retinent quod norunt. Tales sibi, vt aiunt, et musis canere decet. Qui populi regendi curam habem, ad docendum apposici sint oportet. Ne que verd̀ linguae tantum volubilitas hic requiri\pend
\section*{L. TIMOTH. }
\marginpar{[ p.34 ]}\pstart tur, Nam plerosque videmus disertos, quorum prompta loquacitas nihil habet aedificationis. Sed potius commendatur prudentia, vt scite ad vsum populi accommodetur Dei verbum. Operaeprecium autem est reputare, quam nihil ad se pertinere ducant Papistae, quidquid Apostolus praescribit. Neque vero singulis, vel excutiendis, vel enumerandis insistam. Vnum hoc qua diligentia apud eos seruatur? Et sane dos ista superuacua foret, quando ministerium docendi, quasi vile et sordidum, a se depellunt: quod tamen Episcopi maxime proprium erat. At neminem latet, quantum absit a Pauli regula, episcopali titulo κω φὸν πρόσυπον superbe gloriari: si modo theatrico habitu appareat. Quasi cornuta mitra, annulus bene gemmatus, argenteum pedum, et fimiles nugae cum lusoria pompa, spirituale Ecclesiae regimen constituant, quod nihilo magis a doctrina potest auelli, quam nostrum vnusquisque a propria anima.  \pend\pstart Non vinolentum.) Non tantùm ebrietatem hoc verbo damnat: sed quamuis intemperiem in vino ingurgitando. Non tantùm quia pastori pimis indecora est immodica vini potatio: sed etiam quod fere multa peiora secum trahit: vt sunt, rixae, ineptae gesticu latiomnes, lasciuiae: et quae recensere non est necesse. Quod percussorem accipit Chrysostomus pro eo qui lingua ferit, hoc est, conuitiatur, aut iurgiis proteruius insurgit, non recipio. Nec me mouet eius ratio, nihil fore magnum, si manu non caedat Episcopus. Generaliter enim puto hic reprehendi militarem ferociam, qua nihil minus decet Christi seruos. Quanto enim se ludibrio exponant, notum est, qui ad infligendum pugni ictum, ne dicam ad gladium exerendum sunt promptiores, quam ad pacandas sua grauitate aliorum lites. Percussores ergo vocat minaces et Martii caloris. Turpiter lucri cupidi, sunt auari omnes. Nam vbicunque erit auaritia, simul erit hoc turpe, cuius meminit Apostolus. Qui diues fieri vult, et cito vult fieri, inquit ille. Hinc fit, vt auari omnes, etiam si palam id non appareat, ad foedos et illicitos quaestus animum adiiciant. Quare huic vitio opponit con temptum pecuniae: sicuti nullum est aliud remedium, quo corrigatur. Qui paupertateminquam, aequo libentique animo non feret, nunquam effugiet morbum illum sordidae ac seruilis cupiditatis. Percussoris vocabulo opposuit aequum etalienum a pugnis. AEquus autem est, qui placido et moderato animo iniurias ferre nouit, qui multa condonat, nec ad summum ius omnia exigit. Alienus a pugnis, qui fugit contentiones et certamina.  \pend\pstart Qui domui suae bene praesit.) Hinc apparet, non requiri a Paule Episcopum humanae vitae expertem, sed bonum et probatum patremfamilias. Nam vtcunque vulgo plus habeat admirationis coelibatus et vita philosophica, quae prorsus remota sit a communi mo\pend
\section*{CAPVT III. }
\marginpar{[ p.35 ]}\pstart re: experientia tamen ipsa, prudentibus quidem et consyderatis hominibus, hoc demonstrat, eos ad regendani Ecclesiam magis formatos esse et aptiores, qui communem vitam non ignorant: sed in humanae conuersationis officiis sunt exercitati. Ideóque obseruanda est ratio, quae paulopost additur:minime parem fore gubernandae Ecclesiae Dei, qui familiae praeesse nesciet. Atqui tales sunt permulti, imo fere omnes, qui ex ociosa et solitaria vita, quasi ex latebris, extrahuntur: syluestres scilicet quodammodo et nullius humanitatis. Non tamen hic astutum hominem probat Apostolus, et rei domesticae callentem: sed qui familiam gubernare didicerit honesta disciplina. De liberis tamen praecipue loquitur, qui debent patris referre ingenium. Proinde magnum erit dedecus Episcopi, si filios habeat flagitiosae et infamis vitae. De vxoribus aget postea. Nunc, vt dixi, praecipuam domus partem attingit. Primo capite ad Titum, ostendit quid hic significet per nomen reuerentiae. Nam postquam dixit non oportere praefractos et immorigeros esse filios Episcopi, simul addit, neque obnoxios probro dissolutionis, aut intemperantiae. Vult ergo in summa, moribus esse ad omnem castitatem, modestiam, et grauitatem compositis. Argumentum quod sequitur a minori ad maius per se clarum est: nequaquam fore parem regendo populo, qui ad regendam familiam sit ineptus. Praeterquam enim quod inde constat necessariis virtutibus destitui quid authoritatis habere poterit in populo, eum propria domus contemptibilem eum reddat?  \pend\pstart Non nouitium. Quia multi tunc praestanti ingenio et doctrina adducebantur ad fidem: tales Paulus ad episcopatum, statim atque Christo nomen dederint, admitti prohibet. Nam ostendit quan tum sit in ea re periculi. Ventosos enim fere esse constat, plenôsque ostentatione. Ita fiet, vt arrogantia et ambitio praecipites eos agat. Quod dicit Paulus nos experimur. Neque enim ad audendum modo feruidi sunt nouitii: sed etiam stulta confidentia turgidi: acsi nubes transuolare possent. Proinde non sine causa ab honore episcopatus arcentur, donec temporis successu domita sit ingenii altitudo. Iudicium Diaboli tribus modis interpretari licet. Nam Diabolum quidam pro Sathana accipiunt, alii pro calumniatoribus. Priorem sensum magis amplector: quia rarum est, vt iudicium maledicentiam signi ficet. Sed rursus iudicium Sathanae, tam actiue quam passiue sumi potest. Hunc posteriorem senfum Chrysostomus sequitur, cui ego facile assentior. Est enim elegans antithesis, quae indignitatem auget, si is qui praeficitur Ecclesiae Dei, sua elatione in eandem cum Diabolo condemnationem ruat. Quanquam actiuam significationem non re iiciofore vt Diabolo causam sui aceusandi praebeat. Sed verior Chry sostomi opinio.  \pend
\section*{I. TIMOTH. }
\marginpar{[ p.36 ]}\pstart Testimonium ab extraneis.) Difficillimum hoc esse videtur, vt pius vir testes habeat infideles suae integritatis, qui mentiendi aduersum nos libidine insaniunt. Sed intelligit Apostolus, vt quoad externam consuetudinem etiam infideles cogantur eum agnoscere probum virum. Nam etsi omnibus Dei filiis sine causa obtrectant: non possunt tamen sceleratum iudicare, qui honeste et innocenter inter eos versatus sit. Talem probitatis notitiam hic Paulus notat. Ratio additur: ne in probrum incidat, et c. Quod ego ita expono: infamiae expositus, perfrictae frontis esse incipiat, tantóque maiore licen tia se prostituat ad omnem nequitiam. Quod est, Diaboli plagis se irretire. Quid enim spei restat, vbi nullus est peccandi pudor?  \pend
\phantomsection
\addcontentsline{toc}{subsection}{\textit{Diaconos similiter graues, non bilingues, non multo vino deditos, non turpiter lucri cupidos habentes myster ium fidei in pura conscientia. Et hi probentur primùm: deinde ministrent, vbi irreprehensibiles comperti fuerint. Mulieres similiter graues, non calumniatrices, sobrias, fideles in omnibus. Diaconi sint vnius vxoris mariti, qui honeste praesint filiis et domibus suis. Nam qui bene ministrauerint, gradum sibi bonum acquirunt, et multam libertatem in fide, quae est in Christo Iesu. }}
\subsection*{\textit{Diaconos similiter graues, non bilingues, non multo vino deditos, non turpiter lucri cupidos habentes myster ium fidei in pura conscientia. Et hi probentur primùm: deinde ministrent, vbi irreprehensibiles comperti fuerint. Mulieres similiter graues, non calumniatrices, sobrias, fideles in omnibus. Diaconi sint vnius vxoris mariti, qui honeste praesint filiis et domibus suis. Nam qui bene ministrauerint, gradum sibi bonum acquirunt, et multam libertatem in fide, quae est in Christo Iesu. }}\pstart Diaconos similiter.) Non est, quod variae interpretationes serupulum nobis iniiciant. Certum est, Apostolum de iis tractare, qui publicam in Eeclesia functionem exercent. Quo refellitur eorum opinio, qui ministros domesticos putant hoc nomine designari. Quod autem alii ad presbyteros referunt Episcopo inferiores, caret fundamento. Nam ex aliis locis patet, commune omnibus presbyteris fuisse nomen. Idque fateri omnes coguntur. Praesertim locus in primo capite ad Titum eam significationem esse comprobat. Restat. vt Diaconos intelligamus, quorum mentionem facit Lucas Actorum capite sexto. Ii sunt, qui curam pauperum habebant. Sed de officio eorum qui plura volet, petat ex Institutione. Quatuor primae virtutes, quibus ipsos vult esse praeditos, satis per se notae sunt. Nisi quod sedulo notandum est, eum admonere, ne sint bilingues: quia vi tium est, quod in illa administratione difficulter cauetur: et tamen re motissimum, si quod aliud, esse debebat.  \pend\pstart Habentes mysterium fidei.) Perinde est acsi dixisset: Qui  \pend
\section*{CAPVT III. }
\marginpar{[ p.37 ]}\pstart puram religionis doctrinam teneant, et quidem ex animo, eum serio Dei timore:vel vt rite sit in fide institutus: ne quid ignoret, quod hominibus Christianis est cognitu necessarium. Summam tamen piae doctrinae honorifice mysterium appellat. Vt certe Deus per Euangelium sapientiam terrenis hominibus patefacit, quam in coelo suspiciunt Angeli. Quare nihil mirum, si humanum captum superet. Proinde eam cum summa reuerentia amplectendam esse meminerimus: et quia ad tantam altitudinem nunquam conscenderemus proprio marte, suppliciter Deum rogemus, vt eam Spiritu reuelationis nobis reseret. Rursus cum videmus impios, vel ridere, vel nihil gustare: diuinae esse gratiae agnoscamus, quod quae aliis sunt abscondita, nobis in corde sunt et in oculis, vt inquit Moises. Ergo Diaconos in mysterio fidei vult esse eruditos: quia etsi docendi munus non habent, nimis tamen absurdum foret, publicam in Ecclesia personam sustinere, et rudes esse in fide Christiana. praesertim cum saepe incidat necessitas, monendi et consolandi: nisi partibus suis deesse velint. Additur pura conscientia, quae extenditur ad totam vitam: tum vero vt sciant se Deo seruire.  \pend\pstart Et hi probentur.) Vult non ignotos, sed compertae probitatis eligi, quemadmodum Episcopos. Et hinc apparet, irreprehensibiles dici, qui nullo notabili vitio laborent. Porro haec probatio non est vnius horae: sed in longa experientia consistit.  \pend\pstart Mulieres similiter.) Vxores intelligit tam Diaconorum. quam Episcoporum. Debent enim maritis esse adiutrices in officio: quod fieri nequit, nisi prae aliis bene sint moratae. Et quoniam vxorum meminerat, idem praecipit de Diaconis, quod ante de Episcopis. Nemipe, vt vnica vxore quisque eorum contentus, casti et honesti pa trisfamilias sit exemplum: filionsque et totam domum in sancta disci plina contineat. 9l8 Hifl22d 2 ROU  \pend\pstart Qui bene ministrauerint.) Quia vno, aut altero seculo post Apostolorum mortem inualuerat vsus, vt ex Diaconorum ordine eligerentur presbyteri: vulgo hunc locum exposuerunt de tran situ ad gradum superiorem quasi Apostolus ad honorem presbyterii vo cet, qui fideles se Diaconos praestiterint. Ego, etsi non nego Diaconorum ordinem, interdum seminarium esse posse, ex quo sumantur presbyteri: tamen simplicius accipio Pauli verba, qui probe defuncti fuerint hoc ministerio, non paruo honore dignos esse: quia non sit sordidum aliquod exercitium, sed honorificum imprimis munus. Porro hac particula significat, quantùm Ecclesiae intersit, hoc munus a viris selectis administrari, quia existimationem et reuerentiam con ciliat sancta administratio. Caeterùm plusquam ridiculi sunt Papistae, dum in Diaconis suis facere se obcendunt, quod Paulus praescribit.  \pend
\section*{I. TIMOTH. }
\marginpar{[ p.38 ]}\pstart Initiô quorsuni Diaconos suos creant, nisi vt in pompa calicem gestent et ludicris nescio quibus spectaculis pascant oculos imperitorum? Adde quod ne id quidem seruant. Neque enim ab annis quingentis Diaconus vnus creatus est, nisi vt hoc gradu facto mox ad presbyrerium conscenderet. Quantae igitur impudentiae est, quod iactant ad gradum altiorem eos euehi, qui probe ministrauerint, cum eos tantùm sacerdotio suo dignentur qui nullam prioris officii partem vnquam attigerunt?  \pend\pstart Et libertatem in fide.) Merito hoc addidit. Nihil enim est, quod aeque libertateni pariat, vt bona conscientia. Quemadmodum econuerso timidos esse oportet, qui male sibi conseii sunt. Quod si licentia interdum exultant, neque aequabilis est ac perpetua, nec quidquam habet ponderis. Ideo etiam exprimit libertatis speciem: In fide, inquit, quae est in Christo. Quo scilicet maiore fiducia Christo seruiant. Sicuti rursus os habent quasi obturatum et constrictas manus, suntque ad bene agendum inhabiles, qui turpiter in sua functione versati sunt: quia nihil illis habetur fidei, nihil authoritatis tribuitur,  \pend
\phantomsection
\addcontentsline{toc}{subsection}{\textit{Haec tibi scribo, sperans breui ad te venire. Quod si tardauero, vt videas quomodo oporteat in domo Dei versari, quae est Ecclesia Dei viuentis, columna et firmamentumn veritatis. Et sine controuersia, magnum est pietatis mysterium: Deus manifestatus est in carne, iustificatus in spiritu, visus Angelis, praedicatus Gentibus, fidem obtinuit in mundo, receptus est in gloria, }}
\subsection*{\textit{Haec tibi scribo, sperans breui ad te venire. Quod si tardauero, vt videas quomodo oporteat in domo Dei versari, quae est Ecclesia Dei viuentis, columna et firmamentumn veritatis. Et sine controuersia, magnum est pietatis mysterium: Deus manifestatus est in carne, iustificatus in spiritu, visus Angelis, praedicatus Gentibus, fidem obtinuit in mundo, receptus est in gloria, }}\pstart Spem aduentus sui facit Timotheo, partim vt eum animet, partim vt eorum proteruiam comprimat, qui sua absentia magis insolescebant. Neque tamen ficte quidquam Timotheo promittit, vel per simulationem alios territat. Omnino enim sperabat se venturum: vt venisse probabile est, si hanc Epistolam scripsit, quo tempore Phrygiam peragrabat: sicuti refert Lucas Actorumdecimo octauo. Hoc autem nobis argumento sit, quanta solicitudine curauerit Ecclesias, qui praesenti malo remedium in breue tempus differre non sustinuerit. Quanquam continuo post addidit, Epistolam instruendo Timotheo scriptam esse, si forte sibi contingeret, moram facere opinione sua longiorem.  \pend\pstart Quomodo versari in domo Dei.) Hac locutione commen  \pend
\section*{CAPVT III. }
\marginpar{[ p.39 ]}\pstart dat officii pondus ae dignitatem: quia sint pastores veluti oecononn quibus regendam fuam domum commisit Deus. Si quis amplae familiae curam sustineat, noctes ac dies anxio studio iucumbit, ne quid sua negligentia, vel imperitia, vel socordia pessum eat. Si hominibus tantùm datur: quanto plus Deo dabitur? Porro Ecclesiam suam hoc nomine non frustra Deus ipse dignatur. Vtpote qui non modo nos adoptionis gratia in filios recepit: sed etiam in medio nostri habitat. Accedit non vulgaris amplificatio, ex eo quod vocat columnam et firmamentum veritatis. Quid enim magnificentius potuit dici? Quid aeterna illa veritate augustius magisque sacrosanctum, qua et Dei gloria et hominum salus contineturai in vnum cumulum congerantur omnes prophanae philosophiae laudes, quibus a suis sectatoribus ornatur: quid hoc, ad dignitatem coelestis huius sapientiae? quae sola et lux, et veritas, et doctrina vitae, et via, et regnum Dei praedicari meretur. Atqui ea solius Ecclesiae ministerio in terris conseruatur. Quantùm ergo onus Pastoribus incumbit, qui tam inaestimabilis thesauri custodiae praesunt? Impudenter autem 7 nugantur Papistae, qui colligunt ex Pauli verbis sua omnia delyria pro Dei oraculis habenda esse: quia sint columna veritatis, ideôque errare nequeant. Primo videndum est, quorsum Paulus Ecclesiam insignat tam splendido titulo. Proculdubio voluit, proposita officii magnitudine, admonitos esse pastores, quanta illud fide, diligentia, reuerentia administrare debeant. Etenim quam horribilis futura est vltio, si eorum culpa intercidat illa veritas, quae imago est diuinae gloriae lux mundi, salus hominum? Haec sane cogitatio trepidationem alliduam debet incutere pastoribus: non vt eos exanimet, sed ad maiorem vigilantiam acuat. Inde colligere promptum est, quo sensu Paulus ita loquatur. Ecclesia enim ideo columna est veritatis, quia suo ministerio eam tuetur ac propagat. Deus non descendit ipse e coelo ad nos, neque angelos quotidie mittit ad suam veritatem promulgandam sed vtitur Pastorum opera, quos in hunc finem ordinauit. Vt erassius exprimam: Nónne Ecclesia mater est piorum omnium, quae ipsos regenerat Dei verbo, quae educat alitque tota vita, quae confirmat, quae ad solidam perfectionem vsque perducit? Eadem quoque ratione columna veritatis praedicatur: quia doctrinae administrandae munus, quod Deus penes eam deposuit, vni cum est instrumentum retinendae veritatis, ne-ex hominum memoria pereat. Ergo elogium hoc ad ministerium verbi refertur: quo sublato, concidet Dei veritas. Non quia per se minus firma sit, nisi hominum humeris fuleiatur: quemadmodum etiam iidem Papistae blatterant. Nam ista est execrabilis blasphemia: incertum esse Dei verbum, donee ab hominibus certitudinem quasi precarid mutuetur.  \pend
\section*{I. TIMOTH. }
\marginpar{[ p.40 ]}\pstart simpliciter intelligit Paulus quod aliis verbis tradit io ad Romcapite: quoniam fides est ex auditu, nullam fore, nisi sit praedicatio. Itaque hominum respectu sustinet Ecclesia veritatem, quia eam praeconio suo celebrat, quia puram et sinceram retinet, quia transmittit ad posteros. Quod si non resonet Euangelii doctrina, si nulli sint pii ministri qui sua praedicatione veritatem a tenebris et obliuione vindicent: statim mendacia, errores, imposturae, superstitiones, omneque genus corruptelae regnum vindicabunt. Denique silentium in Ecclesia, est exilium oppressióque veritatis. In hac expositione quid omnino est coactum: Nunc, quando tenemus Pauli mentem, redeamus ad Papistas. Primùm, dum ad se transferunt hoc encomium improbe faciunt, alienis plumis se vestiendo. Nam vt euehatur Ecclesia supra tertium coelum: nego id totum ad eos vllo modo pertine re. Quinetiam locum praesentem aduersus eos retorqueo. Nam si Ecclesia columna est veritatis: sequirur non esse apud eos Ecclesiam, vbi non modo sepulta iacetveritas, sed horrendum in modum diruta et euersa pedibus calcatur. An hoc est vel aenigma vel cauillum? Pau Ius Ecclesiam non vult agnosci, nisi in qua excelsa et conspicua stat Dei veritas. In Papatu nihil tale apparet: sed disiectio tantùm et ruinae. Ergo genuina Ecclesiae nota illie non extat. Sed inde hallucinatio, quia quod praecipuum erat, non consyderant, sustineri Dei veritatem pura Euangelii praedicatione: neque hominum inge nio aut sensu niti fulturam, sed ex alto magis pendêre:nempe si a sim plici Dei verbo non disceditur.  \pend\pstart Magnum pietatis mysterium.) Alia rursus amplificatio:e minoris, quam decebat, hominum ingratitudine fieret Dei veritas. Ergo extollit eius precium, cum admonet, magnum esse pietatis arcanum: quia scilicet non agatur de rebus abiectis, sed de reuelatione filii Dei, in quo absconditi sunt thesauri omnes sapientiae. Ex tantarum rerum magnitudine munus suum aestimare debent Pastores, quo plus religionis et timoris adhibeant ad illud exercerdum.  \pend\pstart Deus manifestatus.) Vulgaris interpres, omisso Dei nomi ne, ad mysterium refert quae sequuntur. Inscite omnino et perperam:vt clarius, ex ipsa tractatione patebit. quanquam Erasmum habet suffragatorem, qui tibi ipse sibi fidem eleuat: vt mea refutatione opus non sit. Graeci certe omnes in hanc lectionem consentiunt: Deus manifestatus in carne. Verùm, vt demus nome Dei non fuisse expres sum a Paulo, subaudiendum tamen esse Christi nomen fatebitur, quisquis prudenter omnia expendet. Ego tamen sine difficultate, receptam a Graecis lectionem sequor. Quod manifestationem Christi, qualem postea describit, vocat magnum arcanum, ratio in promptu est. Haec enim est altitudo, profunditas, et latitudo sapientix.  \pend
\section*{CAPVT III. }
\marginpar{[ p.41 ]}\pstart cuius meminit ad Ephesios primo, ad quam sensus omnes nostros ob stupescere necesse est. Nunc ordine excutiamus singula. Non potuit magis proprie de Christi persona loqui, quam his verbis: Deus manifestatus in carne. Nam primo, disertum hic testimonium habemus vtriusque naturae. Verum enim Deum, et verum hominem simul pronunciat. Secundo, notatur inter duas naturas distinctio, cum hine vocat Deum, inde manifestationem in carne ponit. Tertio, vnitas personae significatur, cum vnum eundemque esse tradit, qui cum Deus esset, manifestatus fuit in carne. Ita hoc vno testimonio vera et orthodoxa fides aduersus, Arrium, Marcionem, Nestorium, et Eutichem, egregie munitur. Caeterùm magnam emphasin continent ista antitheta, Deus, in carne. Quantùm enim interest inter Deum et hominem? Et tamen immensam Dei gloriam sic videmus in Christo coniunctam cum hac nostra carnis putredine, vt vnum efficiant.  \pend\pstart Iustificatus in Spiritu.) Sicuti carnem induendo se exina niuit filius Deisita et in ipso spiritualis apparuit virtus, quae Deum esse testaretur. Porro locum hunc varie exponunt. Sed ego germa num Apostoli sensum vt a me quidem percipitur, contentus explicasse, nihil praeterea addam. Primùm iustificatio, hic approbationem diuinae potentiae significat. Vt Psal.1 9. Iudicia Dei iustificata, hoc est, exquisite et ad vnguem perfecta. Et Psal.51.iustificari Deum, pro eo quod est, laudem iustitiae eius conspicuam reddi. Sic et Matthaei vndecimo, eum dicit Christus iustificatam esse sapientiam a filiis suis, intelligit, honorem illi suum exhibitum. Item Lucas septimo capite, cum publicanos iustifieasse Deum commemorat, intelligit, debita reuerentia et gratitudine prosecutos esse Dei gratiam quam in Christo cernebant. Eigo tantundem valet quod hic legimus, ac si dixisset Paulus, eundem, qui humana carne vestitus apparuit, simul tamen declaratum fuisse Dei filium, vt carnis infirmitas nihil derogauerit eius gloriae. Spiritus nomine, comprehendit, quidquid in Christo diuinum fuit ac supra hominem. Idque duas ob causas. Nam quia in carne humiliatus fuerat, nunc exaduerso, illustrationem gloriae ponens, Spiritum opposuit. Deinde, gloria illa, vnigenito Dei filio digna, quam Ioannes in Christo spectatam fuisse docet, non externa pompa, non terreno splendore constabat, sed spritualis fere tota erat. Eadem loquendi forma vsus est primo ad Romanos: Qui factus est ex semine Dauid secundùm carnem, declaratus autem filius Dei in potentia Spiritus. Nisi quod speciem vnam illic adiicit, nempe resurrectionem.  \pend\pstart Visus Angelis, praedicatus Gentibus.) Mira omnia et stupenda: quod Gentes, quae hactenus in caecitate mentis errauerant,  \pend
\section*{I. TIMOTH. }
\marginpar{[ p.42 ]}\pstart dignatus est Deus reuelatione filii sui, quae coelestes quoque Angelos latuerat. Nam quod Angelis visum esse dicit, intelligit tale spectaculum fuisse, quod tam nouitate quam praestantia sua Angelos in se conuerterit. Quam verd̀ noua res et insolens fuerit Gentium vocatio, diximus secundo adEphesios capite. Nec mirum est, Angelis fuisse nouum spectaculam: qui rametsi de humani generis redemptione cognouerant: modum tamen non tenuerunt ab initio: et celatos ideo oportuit quo plus admirationis haberet tanta Dei bonitas.  \pend\pstart Fidem obtinuit in mundo.) Iam illud erat imprimis admirabile, quad eiusdem reuelationis, et Gentes profanas, et Angelos perpetuos regni sui haeredes aeque compotes fecerat Deus. Sed tanta etiam Euangelii praedicati efficacia, miraculum fuit non vulgare, cum Christus tot superatis impedimentis in obedienriam fidei subegit, qui videbantur esse prorsus indomabiles. Nihil certe minus credibile videbatur. Adeo aditus omnes praeclusi erant et obserati, Fides tamen eluctata est, sed incredibili victoriae genere. Postremo, in gloriam receptum esse dicit, nempe ab hac mortali et aerumnosa vita. Ergo sicuti in mundo quoad fidei obedientiam, ita et in Christi persona mira fuit conuersio, dum ex tam abiecta serui conditione euectus est ad dexteram Patris, vt flectatur omne genu.  \pend
\phantomsection
\addcontentsline{toc}{subsection}{\textit{CARVT. IIII. }}
\subsection*{\textit{CARVT. IIII. }}\pstart \huge\textbf{S}\normalsize PIRITVS autemclare dicit, quod in posterioribus temporibus desciscent quidam a fide, attendentes spiritibus impostoribus, et doctrinis daemoniorum in hypocrisi falsiloquorum, cauterio notatam habentium conscientiam, prohibentium matrimonia contrahere, iubentium abstinere a cibis, quos Deus creauit ad percipiedumcum gratiarum actione fidelibus, et qui cognouerunt veritatem: quod omnis creatura Deibona: et nihil reiiciendum, quod cum gratiarum actione sumatur. Sanctificatur enim per sermonem Dei et precationem.  \pend\pstart Spiritus autem.) Diligenter admonuerat de multis Timotheum. lam necessitatem ostendit, quia occurrendum sit periculo, quod imminere Spiritus sanctus denunciat. Venturos scilicet  \pend
\phantomsection
\addcontentsline{toc}{subsection}{\textit{}}
\subsection*{\textit{}}
\section*{CAPVT IIII. }
\marginpar{[ p.43 ]}\pstart falsos dostores, qui nugas pro fidei dostrina obtrudant: qui totam san ctitatem collocantes in externis exercitiis, spiritualem Dei cultum, qui solus est legitimus, obscurent. Et sanê fuit hoc seruis Dei perpetuum certamen cum talibus, quales hic depingit Paulus. Quia enim natura homines ad hypocrisin feruntur, facile illis persuadet Sathan ceremoniis et externa disciplina rite coli Deum. Imo absque magistro hanc omnes fere persuasionem animis infixam habent. Accedit postea Sathanae astus, ad errorem confirmandum. Inde factum est, vt seculis omnibus extiterint impostores, qui fictitios cultus commendarent, quibus obruebatur vera pietas. Porro haec pestis alteram parit, quod in rebus liberis imponitur necessitas. Nam mundus facile sibi prohiberi sinit, quod per Deum licitum erat, vt impune Dei leges sibi transgredi liceat. Hic ergo Paulus in persona Timothei, non Ephesios modo, sed omnes vbique terrarum Eccle sias praemonet, qui fictitios cultus instituendo, et conscientias illaqueando nouis legibus verum Dei cultum adulterent, et corrumpant puram fidei doctrinam. Hic proprius est scopus, quem notari, imprimis est necesse. Porro quo maiore attentione excipiant omnes quod dicturus est, praefatur, certum esse et minime obscurum ora culum Spiritus sancti. Non est quidem dubium, quin reliqua ex codem Spiritu hauserit. Verum vteunque semper audiendus sit tanquam Christi organum:tanien in causa magni ponderis, peculiariter voluit hoc testatum, nihil se proferre, nisi ex Spiritu prophetiae. Solenni itaque praeconio nobis hanc prophetiam commendat. Nec to contentus, addit esse claram, nec vllo aenigmate implicitam.  \pend\pstart In posterioribus.) Hoc vero nemo tunc expectasset, vt in tanta Euangelii luce deficerent. Sed hoc est scilicet, quod dicit Petrus, falsos doctores, sicut populo Israelitico molesti olim fuerunt, ita Christianae Ecclesiae semper fore infestos. Perinde ergo est, ac si diceret: Nunc floret Euangelii dostrina: sed non diu quiescet SaF than, quin zizaniis suffocare studeat purum semen. Proderat haec admonitio seculo Pausi, vt sedulo instantes purae doctrinae, tam Pastores quam reliqui, non circunuenirentur. Nobis hodie non minus est vtilis, cum nihil cernimus accidisse, quod non fuerit diserto Spiri tus vatieinio praedictum. Animaduertere praeterea hic licet, quantam Ecclesiae suae curam agat Deus, dum ita mature occurrit periculis. Multiplices quidem artes habet Sathan quibus in errorem nos inducat: et miris insidiis nos adoritur. Verdm ex aduerso satis mu nit nos Dominus: nisi sponte decipi, nobis liberet. Quare non est quod queramur, potentiores esse tenebras luce, aut superari veritatem a mendacio: sed potius nostrae incuriae ignauiǽque poenas damus, cum abducimur a recta salutis via. Sed obiiciunt, qui sibi in  \pend
\section*{1. TIMOTH. }
\marginpar{[ p.44 ]}\pstart iuis erroribus indulgent, vix posse discerni quosnam aut quales desi gnet Paulus. Quasi vero Spiritus hanc prophetiam de nihilo ediderit, et tantoante promulgauerit. Nam si nullum est certum discrimen, superuacua foret, adeôque ridicula, tota praesens admonitio. Sed absit, vt frustra nos territari putemus a Dei Spiritu, periculi denunciatione: non simul ostendi cauendi rationem. Atque illam calumniam satis per se refutant Pauli verba. Nam quasi digito ostendit malum, quod fugiendum monet. Neque enimin genere de pseudoprophetis verba facit: sed clare speciem exprimit falsae doctrinae. Nempe quae pietatem alligando externis elementis, spiritua lem Dei cultum peruertit ac profanat: vt iam dixi.  \pend\pstart Discedent a fide.) Incertum de magistris ne loquatur, an de auditoribus. Malo tamen ad posteriores referre: quia magistros deinde vocat Spiritus impostores. Et hoc est εμφατικότερον non modo fore, qui spargant impia dogmata, et fidei puritatem inficiant: sed discipulos etiam non defore, quos trahant in suam sectam. Hinc autem nascitur maior perturbatio, cum ita mendacium inualescit. Porro non leue est vitium quod notat, sed atrocissimum crimen, apostasia a fide. Et tamen, prima specie, non videtur in ipsa doctrina tantùm esse mali. Quid enim? An propter ciborum vel coniugii prohibitionem funditus euersa est fides? Sed consideranda est altior ratio: quod peruersum Dei cultum homines, pro libidine comminiscuntur, quod regnum inuadunt in conscientias, quod vsum bonarum rerum a Domino permissum vetare audent. Vbi semel vitiatus est cultus Dei, nihil amplius integrum aut sanum. Imo tota fides periit. Quare, rideant nos Papistae licet, cum tyrannicas eorum leges exagitamus, de obseruationibus externis: scimus tamen, causam maxime seriain et grauissimam nos agere:quia dissipatur fidei doctrina, simul atque talibus corruptelis inficitur Dei cultus. Neque enim de carne aut piscibus lis est, neque de nigro aut cineritio colore, neque de die veneris aut mercurii: sed de insanis ho minum superstitionibus, qui Deum talibus nugis placare volunt, et cultum illi carnalem fingendo, idolum, Dei loco, sibi fingunt. Quis hanc a fide discessionem esse neget?:  \pend\pstart Spiritibus.) Prophetas aut doctores intelligit, quos ideo sic nominat, quia Spiritumiastant, atque hoc titulo se venditant ad populum. Est quidem hoc perpetud̀ verum, qualescunque sint homines, spiritus instigatione loqui. Sed non idem Spiritus omnes instigat. Nam Sachan aliquando Spiritus est mendax in ore pseudoprophetarum, ad decipiendos infideles, qui decipi merentur. Quisquis autem debitam Christo gloriam tribuit, ex Spiritu Dei loquitur teste Paulo. Verdm locutio ista, de qua nunc agimus, hinc initio  \pend
\section*{CAPVT IIII. }
\marginpar{[ p.45 ]}\pstart manauit, quod serui Dei, vt res erat, se habere ex Spiritus reuelatione profitebantur, quaecunque in medium afferebant. Vnde illis inditum fuit Spiritus nomen cuius erant organa. Sathanae verd ministri, falsa aemulatione, tanquam simiae, idem postea iactare coeperunt: et nomen quoque falso ad se traxerunt. Secundùm hanc rationem dicit Ioannes: Probate Spiritus, nunquid ex Deo sint. Et Paulus seipsum exponit, addendo, doctrinis daemoniorum. Perinde enim est, acsi dixisset, attendentes pseudoprophetis, et diabolicis eorum dogmatibus. Iterum obserua, quam non leuis sit error nec dissimulandus, cum alligantur hominum figmentis conscientiae, et simul vitiatur cultus Dei.  \pend\pstart In hypocrisi loquentium.) Si ad daemonia referas: tunc significabit hoc nomen, homines diaboli instigatione fallentes. Potest tamen etiam subaudiri, hominum loquentium. Porrô, ad speciem nunc descendit, cum dicit eos mentiri in hypocrisi, et cauterio notatam habere conscientiam. Et quidem sciendum est, haec duo ita esse coniuncta, vt prius ex secundo nascatur. Nam malae conscientiae, et seelerum cauterio adustae semper ad hypocrisin, tanquam ad prom ptum asylum, confugiunt. hoc est, fucos excogitant, quibus fallantur Dei oculi. Quid enim aliud façiunt, qui student eum placare externis obseruationum fucis? Ergo hypocrisis pro loci praesentis cir cunstantia definienda est. Nam primo ad doctrinam referri debet: deinde cam doctrinae speciem significat, quae spiritualem Dei cultum exercitiis corporalibus mutando, genuinam eius puritatem adulterat. Ita factitios omnes placandi aut promerendi Dei modos comprehendit. Summa haec sit: Diaboli instigatione ferri, quicunque fucatam sanctimoniam inducunt: quia externis ritibus nunquam rite Deus colitur. Nam veri cultores eum adorant in spiritu et veritate. Secundô, hoc esse inutile medicamentum, quo dolores suos mitigant hypocritae, vel potius emplastrum, quo malae conscientiae vulnera, nullo profectu, summa sua pernicie occultant.  \pend\pstart Prohibentium. Speciali designatione posita, nunc duo quoque indiuidua notat: coniugii scilicet prohibitionem, et ciborum. Sunt enim ex illa hypocrisi, quae relicta vera sanctitate, alienos colores in fucum accersit. Nam qui ab ambitione, odio, auaritia, crudelitate aliisque similibus non abstinent, earum rerum abstinentia, quarum libertatem permisit Deus, iustitiam acquirere sibi tentant. Quorsum enim istis legibus onerantur conscientiae? Nisi quia perfectio extra Dei legem quaeritur. Hoc autem non fit nisi ab hypocritis, qui vt impune transgrediantur cordis iustitiam quam lex requirit, interiorem suam nequitiam obseruatiunculis istis tegere, quasi velis oppositis, tentant. Clara fuit haec periculi denunciatio vt non  \pend
\section*{I. TIMOTH. }
\marginpar{[ p.46 ]}\pstart fuerit difficilis cautio: fiquidem aures Spiritui sancto tam diserte monenti homines praebuissent. Videmus tamen in plurimis praeualuisse Sathanae tenebras: vt nihil tanta lux insignis et memorandi oraculi profuerit. Etenim non multapost Apostoli mortem exor ti sunt Encratitae, qui nomen sibi a continentia indiderunt. Taciani, Cathari, Montanus cum sua secta, et tandem Manichaei, qui ab esu car nium, et coniugio abhorrerent: et tanquam res profanas damnarent. Quanquam autem propter suam arrogantiam, quod alios suis placitis obnoxios facere volebant, ab Ecclesia fuerunt repudiati:eos tamen ipsos, qui illis restiterunt, plus aequo in eorum errorem declinasse constat. Nolebant quidem legem imponi Christianis. Sed interca, plusquam par esset, superstitiosis obseruationibus deferebant: si quis fugeret coniugium si quis carnem non gustaret. Tale est ingenium mundi. Semper carnali more Deum, quasi carnalis esset, fibi colendum somniat. Paulatim, rebus in deterius collapsis, obtinuit haec tyrannis, ne sacerdotibus aut monachis fas esset matrimonium contrahere: ne quispiam carnem gustare certis diebus auderet. Quare non immerito Papistis hodie obiicimus hanc prophetiam: quando coelibatum et ciborum abstinentiam seuerius vrgent, quam vllum Dei praeceptum. Ipsi autem ingenioso cauillo effugere se putant, eum Pauli verba torquent aduersus Tacianos, aut Manichaeos, ac similes. Quasi verd̀ non idem patuisset Tacianis effugium, reiiciendo etiam Pauli sententiam ad Cathaphryges, et authorem illorum Montanum. Quasi non promptum fuisset Cataphrygibus, Encratitas supponere veluti reos suo loco. At Paulus non de personis hic agit, sed de reipsa. Quamobrem etiam si centum diuersae sectae proferaneur, quae tamen eadem laborent hypocrisi in cibis prohibendis, omnes eundem sustinebunt reatum. Vnde sequitur, nullo profectu obtendi a Papistis veteres haereticos, quasi soli illi taxentur. Semper videndum est, an sint in eadem noxa. Excipiunt, se Encratitis et Manichaeis esse dissimiles: quia non simpliciter vsum coniugii, et car nium interdicunt: sed certis tantum diebus cogunt ad carnis abstinentiam, solos autem manachos et sacerdotes cum monialibus ad votum coelibatus cogunt. Verdm haec quoque nimis friuola est exeusatio. Nam sanctimoniam nihilominus in his rebus locant: deinde falsum et adulterinum Dei cultum instituunt: postremo consci entias alligant necessitati, a qua debebant esse liberae. Extat apud Eusebium libro quinto, fragmentum quoddam ex scriptis Apollonii, vbi inter caetera Montano exprobrat, primum fuisse, qui nuptias soluerit, et ieiuniorum leges imposuerat. Non dicit in vniuersum prohibuisse matrimonium aut cibos. Satis est, si quis religionem con lescssdnu.n. elande las unl rerum obseruatione. Deum colere pra\pend
\section*{CAPVT IIII. }
\marginpar{[ p.47 ]}\pstart cipiat. Nam rerum liberarum prohibitio, siue generalis sit, sine specialis, semper est tyrannis Diabolica. Melius tamen id de cibis verum esse, ex proximo membro patebit.  \pend\pstart Quos Deus creauit.) Notanda ratio, debere nos in vsu ciborum concessa a Deo libertate esse contentos, quia eos creauerit in hunc finem. Hoc piis omnibus inaestimabile gaudium est, cum sciunt alimenta omnia, quibus vescuntur, sibi de manu in manum a Domino porrigi:vt purus sit ac legitimus illorum vsus. Quanta vero hominum arrogantia est, eripere quod Deus largitur? An cibos ipsi crearunt? An Dei creationem possunt irritam facere? Semper itaque nobis in mentem veniat, eum qui creauit, liberum quoque vsum fecisse, quem frustra impedire conantur homines. Creauit, inquam,cibos Deus ad percipiendum, hoc est, vt illis fruamur. Hunc finem euertere nunquam poterit humana authoritas. Addit autem, cum gratiarum actione: quia nihil Deo pro sua liberalitate rependere possumus, quam gratitudinis testimonium. Porro non potest esse gratiarum actio sine sobrietate et temperantia. Neque enim Dei beneficentiam vere agnoscit, qui ea perperam abutitur,  \pend\pstart Fidelibus et iis qui cognouerunt.) Quid ergo An non solem suum quotidie oriri facit Deus super bonos et malos? An non eius iussu terra impiis panem producit? An non eius benedictiont etiam pessimi aluntur? Est enim vniuersale illud beneficium, quod Dauid Psalmo centesimo quarto decantat. Respondeo, Paulum de vsu licito hic agere, cuius ratio coram Deo nobis constat. Huius minime compotes sunt impii, propter impuram conscientiam, quae omnia contaminat: quemadmodum habetur primo capite ad Titum, Et sane, proprie loquendo, solis filiis suis Deus totum mundum, et quidquid in mundo est, destinauit. Qua ratione etiam vocantur mundi haeredes. Nam hac conditione constitutus initio fuerat Adam omnium Dominus, vt sub Dei obedientia maneret. Proinde rebellio aduersus Deum, iure, quod illi eollatum fuerat, ipsum vna cum posteris spoliauit. Quoniam autem subiecta sunt Christo omnia:eius be neficio in integrum restituimur: idque per fidem. Quare veluti alienum furantur aut praedantur infideles, quodeunque vsurpant. Posteriore membro definit quos vocet fideles. Nempe qui notitiam habent sanae doctrinae. Neque enim fides, nisi ex verbo Dei:ne falso putemus confusam esse opinionem, quemadmodum fingunt Papistae. Quia omnis creatura.) Ciborum vsus aestimari debet partim ex illorum substantia, partim ex cius persona qui yescitur. Vtroque igitur argumento pugnat Apostolus. Nam quod ad cibos attinet, puros esse asserit: quia sint a Deo conditi. Consecrari vero nobis illorum vsum fide et precatione. Bonitas ereaturarum euius me\pend
\section*{I. TIMOTH. }
\marginpar{[ p.48 ]}\pstart minit, relationem habet ad homines, et eam quidem non ad corpus, aut valetudinem:sed ad conscientias. Ne quis scrupulose extra loci circunstantiam philosophetur. Nam vno verbo significat Paulus, non esse immunda nobis coram Deo, vel polluta, quaecunque ex Dei mann proueniunt, et nobis in vsum destinantur, quin illis vesci, quoad conscientiam, liceat. Si quis obiiciat, multa animalia olim pro nunciata fuisse immunda in lege: et fructum quem arbor scientiae bo ni et mali proferebat exitialem homini fuisse: responsio est, non ideo tantùm puras vocari creaturas, quia Dei fint opera: sed quia eius beneficentia nobis sint datae. Semper enim respicienda est Dei ordinatio: et quid iubeat, aut prohibeat.  \pend\pstart Sanctificatur enim.) Confirmatio est proximae particulae, si eum gratiarum astione sumatur. Et est argunientum a contrariis. sunt enim contraria, sanctum et profanum. Nunc videamus, qualis sit omnium bonorum, quae ad sustinendam praesentem vitam pertinent, sanctificatio. Eam Paulus sermone Dei et oratione constare, testatur. Sed notandum est, fide oportere apprehendi hunc sermonem. vt prosit. Tametsi enim Deus ipse solo Spiritu oris sui omnia sanctificat:nos tamen id beneficium, non nisi fide, percipimus. Huc accedit oratio:quia et a Deo panem nostrum quotidianum petimus ex Chri sti mandato: et gratiarum actione prosequimur eius bonitatem. Vnde colligendum est, impuram esse vsurpationem omnium donorum Dei, nisi adsit vera cognitio, et inuocatio nominis Dei. Ac beluinum esse vescendi morem, cum homines sine vlla precatione ad mensam se ingerunt, ac bene saturi, sepulta Dei mentione, alio se proripiunt. Quod si in alimentis communibus, quae vna cum ventre corruptioni sunt obnoxia, requiritur talis sanctificatio: quid de spirituasibus sacramentis est sentiendume Si verbum absit, et inuocatio Dei ex fide quid restat nisi profanum? Hic notanda est distinctio inter benedictionem mysticae et communis mensae. Cibos enim, quibus ad corpus pascendum vescimunin hunc finem benedicimus, vt pura sit ac legitima perceptio. At panem et vinum in mystica Coena augustiore modo conseeramus, vt nobis sint corporis ac sanguinis Christi pignora.  uta proupmolobiu 111  \pend
\phantomsection
\addcontentsline{toc}{subsection}{\textit{2 2 Haec suggerens fratribus, bonus eris Iesu Christi minister, innutritus sermonibus fidei, et bonae doctrinae, quamsecutus es. Profanas autem et aniles fabulas deuita. Quin potius exerce teipsum ad pietatem. Namcorporalis exercitatio pauluIum habet vtilitatis: at pietas ad omnia vtilis est, vt }}
\subsection*{\textit{2 2 Haec suggerens fratribus, bonus eris Iesu Christi minister, innutritus sermonibus fidei, et bonae doctrinae, quamsecutus es. Profanas autem et aniles fabulas deuita. Quin potius exerce teipsum ad pietatem. Namcorporalis exercitatio pauluIum habet vtilitatis: at pietas ad omnia vtilis est, vt }}
\section*{CAPVT IIII. }
\marginpar{[ p.49 ]}
\phantomsection
\addcontentsline{toc}{subsection}{\textit{quae promissiones habeat vitae praesentis ac futurae. Indubitatus sermo, dignus que qui modis omnibus approbetur. Nam in hoc et laboramus, et probris afficimur, quod spem fixam habemus in Deo viuente, qui seruator est omnium hominum, maxime fidelium. }}
\subsection*{\textit{quae promissiones habeat vitae praesentis ac futurae. Indubitatus sermo, dignus que qui modis omnibus approbetur. Nam in hoc et laboramus, et probris afficimur, quod spem fixam habemus in Deo viuente, qui seruator est omnium hominum, maxime fidelium. }}\pstart Haec suggerens.) Hoc verbo ad frequentem earum rerum commonefactionem hortatur. Quod iterum paulopost et tertio repetet. Sunt enim res eiusmodi, quas saepius in memoriam reuocari expediat. Ac notanda est antithesis, quae subest. Non enim doctrinam, quam commendat, falsis opponit, aut impiis dogmatibus: sed nihili argutiis, quae non aedificant. Eas vult longua obliuione sepultas esse, cum Timotheum aliis suggerendis intentum esse iubet.  \pend\pstart Bonus eris minister.) Quiduis potius saepenumero affectant homines, quam se Christo probare. Hinc fit, vt ingenii, facundiae, reconditae scientiae laudem plures aucupentur. Atque eadem ratio est, cur posthabeantur res necessariae, quae ad conciliandam vulgi admirationem minus valent. At Paulus hoc vno contentum esse Ti motheum iubet, vt fidelis sit Christi minister. Et profecto titulus hic longe honorificentior nobis esse debet, quam si millies Seraphici, subtilesque vocaremur. Meminerimus ergo, sicuti summus est honos pii pastoris, censeri bonum Christi seruum: ita nihil aliud in tota sua administratione captandum esse. Nam quisquis alio tendet, fieri porerit, vt plausum reportet ab hominibus, sed Deo non placebit. Quare, ne priuemur tanto bono: discamus nihil aliud expetere, aut preciosum ducere, quin prae hoc vno vilescant nobis omnia.  \pend\pstart Enutritus.) Quia participium est medium, tranfferri etiam astiue posset, enutriens. Sed cum nulla addita sit constructio: videtur id mihi esse nonnihil coactum. Quare malo passiue accipere, vt confirmet proximam exhortationem ab educatione Timothei. Acsi diceret: Quemadmodum ab incunabulis fuisti recte institutus in fide: et quasi vna cum lacte imbibisti sanam doctrinam: et continuos in ea progressus hucvsque fecisti: da operam vt fideliter ministrando, talem esse te comprobes. Huc etiam conuenit verbi compositio. Fides hic pro summa doctrinae Christianae capitur. Et quod de sana doctrina mox addidit, exegeticum est. Nam significat aliis omnibus doctrinis nihil sanitatis inesse: quantumuis sint plausibiles. Particula haec, quam secutus es, perseuerantiam significat. Nam multi, qui a pueritia Christum pure didicerant, postea degenerant temporis progressu: qui\pend
\section*{I. TIMOTH. }
\marginpar{[ p.50 ]}\pstart bus dissimilem facit Timotheum.  \pend\pstart At pietas ad omnia vtilis.) Hoc est, qui pietatem habet, illi nihil deest, etiam si careat istis adminiculis. Nam pietas se sola conten ta est ad solidam perfectionem. Principium est, medium et finis Chri stianae vitae. Ideo vbi integra est, nihil est mutilum. Christus non adeo austeram vitae rationem tenuit, quam Ioannes Baptista. An ideo pilo deteriore Summa sit incumbendum esse prorsus in vnam  \pend
\section*{CAPVT IIII. }
\marginpar{[ p.51 ]}\pstart pietatem: quia vbi semel eam assecuti fuerimus, nihil Deus in noonamplius desideret.  \pend\pstart Promissiones habens.) Summa consolatio, quod Deus nusquam deesse piis velit. Nam vt perfectionem nostram in pietate statuerat, ita nunc eam facit verae foelicitatis complementum. Quia autem initium foelicitatis est in hac vita, ideo ad hanc quoque exten dit promissiones diuinae gratiae, quae solae nos efficiunt beatos, et sine quibus sumus miserrimi. Nam Deus in hac quoque vita se nobis fore patrem testatur. Verdm meminerimus, inter bona vitae praesentis et futurae discernere. Nam eatenus nobis benefacit Deus in hoc mundo, vt gustum duntaxat suae bonitatis praebeat: et tali gustu nos alliciat ad caelestium bonorum desyderium, vt in illis acquiescamus. Ita fit, vt bona praesentis vitae non modo permixta sint plurimis erumnis: sed quasi obruta. Neque enim nobis expedit hic abundare, ne luxuriemur. Porro ne quis operum merita hinc efferat: tenendum est quod iam diximus, pietatem non modo bonam conscientiam erga homines ac Dei timorem, sed fidem quoque et inuocationem com  \pend\pstart HRI Fidelis sermo.) Non abs re tanta asseueratione vtitur. Est enim paradoxum valde pngnans cum sensu carnis: Deum suis in hoc mundo suppeditare, quae ad beate et foeliciter viuendum pertinent: cum saepenumero destituti sint bonis omnibus, ideóque puten tur a Deo abiecti. Ergo non contentus simplici doctrina, hoc veluti clypeo contrarias omnes tentationes propulsat.  \pend\pstart In hoc enim laboramus.) Anticipatio est, qua soluit illamt quaestionem: fideles, quia premantur omne genus aerumnis, omnium esse miserrimos. Ergo vt ostendat, non debere ab externa specie aestimari eorum conditionem, separat eos a reliquis, primùm in causa: deinde in exitu. Vnde sequitur, nihil illis decedere ex promissionibus quas dixit, cum rebus aduersis vexantur. Summa est, non esse miseros fideles in afflictionibus: quia bona conscientia eos sustentat: et prosper laetusque exitus eos manet. Quia autem foelicitas vitae praesentis duobus potissimùm membris constat, honore et commoditatibus: ponit duo contraria, labores et probrum: priore vocabulo significans incommoda quaeuis et molestias: vt sunt paupertas, frigus, nuditas, fames, exilia, spoliationes, vincula, verbera, et aliae persecutiones. Consolatio liaec, speramus in Deum viuum, ad causam refertur. Nam adeo non sumus miseri, cum patimur propter iustitiam, vt potius ea sit iusta materia gratulationis. Adde quod afflictiones nostrae cuni spe in Deum viuum coniunctae sunt. Imo spes est quasi fundamentum. Atqui ea nunquam pudefacit. Sequitur ergo, piis in lucro deputanda esse omnia.  \pend
\section*{I. TIMOTH. }
\marginpar{[ p.52 ]}\pstart Qui seruator est.) Secunda consolatio, quae tamen ex priore dependet. Nam liberatio, de qua loquitur, est quasi fructus spei. Quod vt clarius pateat, sciendum est, esse argumentum a minori ad maius. Nam σοτῆρος nomen, hic est generale, pro eo qui tuetur ac conseruat. Intelligit Dei beneficentiam ad omnes homines peruenire. Quod si nemo est mortalium, qui non sentiat Dei erga se bonitatem, eiusque sit particeps: quanto magis eam experientur pii, qui in eum sperant? An non peculiarem ipsorum gerer curam? An non multo liberalius se in eos effundet An non saluos praestabit:  \pend
\phantomsection
\addcontentsline{toc}{subsection}{\textit{Praecipe haec et doce. Nemo tuam iuuentutem despiciat. Sed esto exemplar fidelium, in sermone, inconuersatione, incharitate, in spiritu, in fide, in castitate. Donec venio attende lectioni, exhor tationi, doctrinae. Ne donum, quod in te est, negligas: quod tibi datum est per prophetiam, cum impositione manuum presbyterii. Haec cura, in his esto: vt profectus tuus in omnibus manifestus fiat. Attendetibi ipsi et doctrinae. Permane in his. Hoc enim si feceris, et teipsum seruabis, et eos qui te audiunt. }}
\subsection*{\textit{Praecipe haec et doce. Nemo tuam iuuentutem despiciat. Sed esto exemplar fidelium, in sermone, inconuersatione, incharitate, in spiritu, in fide, in castitate. Donec venio attende lectioni, exhor tationi, doctrinae. Ne donum, quod in te est, negligas: quod tibi datum est per prophetiam, cum impositione manuum presbyterii. Haec cura, in his esto: vt profectus tuus in omnibus manifestus fiat. Attendetibi ipsi et doctrinae. Permane in his. Hoc enim si feceris, et teipsum seruabis, et eos qui te audiunt. }}\pstart Praecipe haec.) Intelligit eius generis esse doctrinam, cuius nullum debeat esse fastidium, etiam si quotidie audiatur. Nam et alia docenda sunt. Sed emphasin habet demonstratiuum haec: quia signi ficat, non esse res leuis momenti, quas semel quasi transeundo attigisse sufficiat: sed potius quotidiana repetitione dignas esse, quia non possint nimium inculeari. Prudentem itaque pastorem decet expendere, quae maxime sunt necessaria vt illis insistat. Nec est quod vereatur fastidium. Nam quisquis Dei est, libenter identidem audiet, quae toties dicere expedit.  \pend\pstart Nemo tuam iuuentutem.) Hoc dicit, tam aliorum respectu, quam ipsius Timothei. Quantum ad alios, non vult Timothei aetatem esse obstaculo, quominus eam, quam meretur, reuerentiam obtineat: modo se alias ita gerat, vt decet Christi ministrum. Et simul ipsum Timotheum admonet, vt quod deest aetati, compenset grauitate morum. Aesi dixisset: Fac vt morum grauitate, tantum reuerentiae tibi concilies, ne quid aetas tua iuuenilis, quae alioqui contemptui obnoxia esse solet, de tua authoritate minuat. Vnde agnoscimus, et  \pend
\section*{CAPVT IIII. }
\marginpar{[ p.53 ]}\pstart iuuenem adhuc fuisse Timotheum, qui tamen inter multos pastores longe eminebat: et perperaui eos facere, qui annorum numero metiuntur, quantum homini debeant. Porro quae vera sint decora, mox praescribit. Non externae scilicet laruae: vt sunt lituus, infulae, annulus, pallium, et eiusmodi ineptiae, aut puerorum crepundia:sed doctrinae sanitas et sanctitas vitae. Sermone et conuersatione cum dicit, hoc perinde valet, aesi dixisset, dictis et factis adeôque tota vita. Quae sequuntur, partes sunt piae conuersationis: charitas, spiritus, fides, etc. Nomine ipiritus intelligo feruorem zeli Dei. Si quis tamen generalius malit accipere, non repugno. Castitas non tantùm libidini opponitur, sed puritatem totius vitae significat. Hinc discimus, ineptos esse, et ridiculos, qui nihil honoris sibi exhiberi conqueruntur, cum nihil habeant laude dignum: sed potius contemptui se exponant, tum sua inscitia, tum impuro vitae exemplo, aut leuitate, aut aliis sor dibus. Vnica enim haec via est comparandae reuerentiae, si virtutum praestantia nos a contemptu vindicemus.  \pend\pstart Donec venio.) Nouerat Timothei diligentiam. Et tamen illi commendat assiduam Scripturae lectionem. Quid enim alios docebunt pastores, nisi discendo sint intenti? Quod si tantus vir admonetur, vt proficere in dies studeat: quanto magis in nos competit ista admonitio? Vae igitur eorum socordiae, qui Spiritus oracula non reuoluunt dies et noctes, vt inde instruantur ad munus suum exercendum. Caeterùm ne ociosa lectio sufficere credatur, simul ostendit, in vsum esse explicandam: cum doctrinae et exhortationi instare praecipit. Acsi praecipe ret ipsum discere, quod aliïs communicaret. Notandus etiam est hic ordo, quod lectionem doctrinae et exhortationi praeponit. Nam certe fons omnis sapientiae est Scriptura, vnde haurire debent pastores quidquid proferunt apud gregem.  \pend\pstart Ne donum quod in te est.) Hortatur, vt in Ecclesiae aedificationem conferat gratiam qua praeditus est. Neque enim perirevult Dominus talenta, aut inutiliter sub terra defodi, quae apud vnumquemque depo suit, vt lucrum afferant. Negligere donum, est neglectim suppressum habere per desidiam, ita vt, quasi rubigine contracta, nullo vsu, deteratur. Ergo consyderet quisque nostrum quid facultatis habeat: vt id sedulo, in vsum applicet. Dicit gratiam illi collatam esse, per prophetiam. Quomodo? Quia seilicet, vt ante diximus, Spiritussanctus oraculo Timotheum destinauerat, vt in ordinem pastorum cooptaretur. Neque enim delectus tantùm fuerat hominum iudicio, vt fieri solet: sed praecesserat spiritus nuncupatio. Dicit collatam esse, cuui impositione manuum. Quo significat, vna cum ministerio, necessariis etiam dotibus ornatum fuisse. Vsitatum fuit ac solenne Apostolis, mi nistros per impositionem manuummordinare. Ac de hoc quidem ritus  \pend
\section*{I. TIMOTH. }
\marginpar{[ p.54 ]}\pstart eiusque origine et significatione aliquid prius attigi: et reliqua ex Institurione petere licet. Presbyterium qui hic collectiuum nomen esse putant, pro collegio presbyterorum positum, recte sentiunt, meo iudicio. Tametsi oninibus expensis, diuersum sensum non male quadrare fateor: vt sit nomen officii. Ceremoniam pro ipso actu ordinationis posuir. Itaque sensus est, Timotheum, cum Prophetarum voce ascitus fuit in ministerium, et deinde solenni ritu ordinatus, simul gratia Spiritussancti instructumn fuisse ad functionem suam exequen dam. Vnde colligimus, non inanem fuisse ritum: quia consecrationem, quam homines impositione manuum figurabant, Deus Spiritu suo impleuit.  \pend\pstart Haec cura.) Qua plus est difficultatis in fideli Ecclesiae administratione, eo magis instare sibi debet pastor, neruósque omnes intendere, idque non ad exiguum modo tempus: sed perseueranter. Ergo admoner Paulus, non esse locum oscitantix vel remissioni: sed requiri summam diligentiam et assiduitatem. Cum addit, vt profectus tuus apparear: significat, in hoc esse elaborandum, vt per eius manum magis ac magis promoueatur Ecclesiae aedificatio: atque dignum operaeprecium extet. Est enim non vnius diei opus. Quare ad quotidianos progressus enitendum est. Alii referunt ad Timothei personam vt in melius proficiat. Sed malo de ministerii effectu interpretari. In omnibus, tam masculinum esse potest, quam neutrum. Ita duplex erit sensus vel, vt omnes perspiciant, quae ex eius laboribus prouenient, increuienta: vel, vt omni ex parte, aut modis omnibus, quod idem est, se ostendant.  \pend\pstart Attende tibi ipsi.) Duo sunt curanda bono pastori: vt docendo inuigilet, ac seipsum purum custodiat. Neque enim satis est, si vitam suam componat ad omnem honestatem, sibique cauear, ne quod edat malum exemplum, nisi assiduum docendi studium adiungat sanctae vitae. Et parum valebit doctrina, si non respondeat vitae honestas et sanctitas. Non ergo abs re Paulus Timotheum incitat, vt tani sibi attendat quam doctrinae. Et rursus commendat illi constan tiam, ne vnquam defatigetur. Multa enim identidem accidunt, quae possint nos abducere a recto cursu: nisi fixo pede fortiter stemus ad resistendum.  \pend\pstart Hoc si feceris.) Hie non leuis est stimulus ad excĩtandam pastorum solicitudinem, eum audiunt suam et populi salutem in eo consistere, si muneri suo gnauiter et continenter incumbant. Nec absurdum videri debet, quod Timotheo has partes assignat, saluandae Ecclesiae. Nam certe quidquid Deo acquiritur, saluum est. Praedicatione autem Euangelii colligimur ad Christum. Et sicuti pastoris infidelitas aut negligentia Ecclesiae est exitialis: ita merito fidei et dili\pend
\section*{CAPVT V. }
\marginpar{[ p.55 ]}\pstart gentiae adscribitur salutis causa. Solus quidem Deus saluat. Cuius gloriae ne minimam quidem portionem ad homines transferri, fat est. Sed nihil gloriae suae Deus derogat, eum ad salucem administrandam vtitur hominum opera. Dei ergo solius beneficium est salus nostra quia et ab ipso solo est, et sola eius virtute peragitur. Ideóque illi tanquam vnico authori ferenda est accepta, Sed non excluditur propte rea hominum ministerium: nec id totum obstat, quin salutare sit regimen illud, quo constare Eeclesiae integritatem alibi Paulus docet 4. ad Ephes.cap. Quinetiam et hoc quoque in solidum Dei est opus. Quia et ipse est, qui format bonos Pastores, et agit suo Spiritu, et eorum labori ne irritus sit, benedicit. Quod si suis auditoribus saluti est bonus Pastor:sciant mali et ignaui, eorum quibus praesunt. imputandum sibi esse exitium. Nam sicuti gregis salus pastoris est corona: ita quidquid perierit, ab ignauis pastoribus, reposcetur. Seruare autem se pastor dicitur, cum fideliter munus sibi iniunctum exequendo, vocationi suae seruit. Non modo quia horrendam illam vltionem effugit, quam Dominus per Ezechielem minatur: sanguinem de manu tua reposcam: sed quia illud vsitatum est, salutem suam fideles peragere, cum in cursu salutis suae pergunt: de qua locutione diximus in secundum ad Philip. cap.  \pend
\endnumbering\beginnumbering\section{CAPVT. V.}
\phantomsection
\addcontentsline{toc}{subsection}{\textit{\huge\textbf{S}\normalsize ENIOREM ne aspere obiurges: sed hortare vt patrem, iuniores vt fratres, mulieres natu grandiores vt matres, iuniores vt sorores cum omni castitate. Viduas honora, quae vere sunt viduae. Porro si qua vidua liberos aut nepotes habet, discant primùm Doceit propriam domum pie tractare, et mutuum rependere progenitoribus. Hoc enim acceptum est coram Deo. }}
\subsection*{\textit{\huge\textbf{S}\normalsize ENIOREM ne aspere obiurges: sed hortare vt patrem, iuniores vt fratres, mulieres natu grandiores vt matres, iuniores vt sorores cum omni castitate. Viduas honora, quae vere sunt viduae. Porro si qua vidua liberos aut nepotes habet, discant primùm Doceit propriam domum pie tractare, et mutuum rependere progenitoribus. Hoc enim acceptum est coram Deo. }}\pstart Nunc lenitatem et moderationem Timotheo commendat in corrigendis vitiis. Medicina est correctio, quae semper amari aliquid habet, ideóque minus adlubescit. Deinde quia iuuenis erat Timotheus, minus tolerabilis fuisset eius seueritas, nisi admodum tempera ta. Praecipit ergo, vt maiores natu reprehendat, tanquam parentes. Imo vtitur mitiore verbo hortandi scilicet. Fieri autem nequit, quin reuerentia tangamur, vbi repraesentamus nobis ob oculos Patris vel  \pend
\vspace{0.5cm}\noindent
\vspace{0.2cm}\rule{1cm}{0.2pt}\\ 
\hspace{0.2cm}\textit{mg}
\footnotesize 
\normalsize| 
\section*{I.TIMOTH. }
\marginpar{[ p.56 ]}\pstart matris personam. Ita fit, vt in locum durioris vehementiae mox succedat modestia. Notandum tamen, quod non ita parci vult senibus aut indulgeri, vnimpune et sine correctione peccent: tantùm vult eorum aetati haberi monnihil honoris, quo aequiore animo se admoneri sinant. Quinetiam, in minori aetate vult moderationem adhiberi: tametsi non aequali modo. Semper enim oleo teniperandum est acetum. Sed hoc interest, quod maioribus deferenda est reuerentia: aequales fraterna mansuetudine excipi debent. Hinc admonentur Pastores, non muneris tantùm habendam esse rationem, sed peculiariter etiam videndum esse singulis, quid aetatem suam deceat. Neque enim eadem omnibus conueniunt. Veniat ergo illud in mentem: si decorum seruant in scena histriones, non esse pastoribus in tam excelso loco negligendum. Parcicula, cum omni castitate, ad iuniores pertinet. Namin illa aetate semper timen da est omnis suspitio. Neque tamen vetat Paulus, ne quid Timotheus agat libidinose, aut ne lasciuiat cum iuuenculis. Hac enim pro hibitione non fuit opus. Sed tantam praecipit cauendum, ne improbis detur obloquendi ansa. Quod vt fiat, requirit castam grauitatem, quae reluceat in omni congressu et colloquio: quo liberius cum iuuenibus communicet, absque vlla sinistra fama.  \pend\pstart Viduas honora.) Honoris nomen hoc loco, non quanuis obseruantiam significat, sed peculiarem curam quam suscipiebant episcopi in veteri Ecclesia. Recipiebantur enim viduae in Ecclesiae fidem, vt ex publico alerentur. Tantundem ergo valet haec locutio, ac si dixisset: In eligendis viduis, quae sint sub tua et diaconorum cura, earum tibi habenda est rario, quae vere sunt vidux. Qualis autem fuerit earum condirio, fusius deinde tractabimus. Hse tamen notanda est ratio, eur Paulus non adnuittar inisi prorsus viduas et orbas simul, quae sobole sint destitutae. Namça conditione se Ecclesiae addicebant, vt abdicarent se omni priuata familiae cura: et omnia impedimenta abiicerent. Proinde merito Paulus, matresfamilias recipi prohibet. Cum nominat vere viduas, allusio est ad nomen Graecum: quod est ἀπο το͂ κορωοται. Quod est destitui et priuari. DR  \pend\pstart Si qua vidua.) Varie exponitur hie locus. Inde autem nascitur ambiguitas, quod posterius membrum tam ad viduas, referri porest, quam ad earum liberos. Nec obstat, quod verbum, discant, est plurale: cum de vidua locutus sit Paulus in nuniero singulari. Numeri enim mutatio satis vsitata est in oratione indefinita: hoc est, vbi de toto genere, non de indiuiduo, tractatur. Qui de viduis accipiunc, hunc pucant esse sensum: Discant pia familiae administra tione bene ficium educationis, quod a maioribus acceperunt, erga  \pend
\section*{CAPVT V. }
\marginpar{[ p.57 ]}\pstart minores repende re. Ita Chrysostomus, et alii quidam. Simplicius tamen aliis videtur de liberis et nepotibus accipere. Matremi itaque eorum iudicio, aut auiam, materiam esse docet, in qua exerceant suam pietatem. Nihil enim magis est naturale, quam ἀντιπολαργια. Quare minime aequum est, eam ab Ecclesia impediri. Prius ergo quam oneretur Ecclesia, faciant illi suum officium. Hactenus retuli, quid alii sentiant. Velim tamen lectores expendere, an non aprior sit ista translatio: doceant pie se domi gerere. Nam μαντάναν, vtrunque significat, docere et discere. Porro verbum ευσεβ ἀν, adiue omnes fere accipiunt, quia sequitur accusatiuus. Verùm illa ratio non cogit: quia Graecis familiare est subaudire praepositionem. Apte igitur contextui quadraret haec expositio: Doceant matronae liberos aut nepores, pietarem erga familiam exercere, ac mutuam yicem rependere majoribus.  \pend\pstart Hoc enim acceptum.) Monstrum esse fatentur omnes, si quis ingratum se maioribus praebeat. Id enim dictat sensus naturae. Nec modo ingenita est omnibus haec persuasio: secundum pietatis gradum esse erga parentes: sed ipsae quoque Ciconiae gratirudinem suo exemplo nos docent. Vnde et nomen ἀντιπαλαγρίας. Verùni Paulus eo non contentus, idem a Deo sancitum esse pronunciat. Acsi diceret: Non est quod putet quispiam,id natum esse ex hominum opinione: sed Deus ita ordinauit.  \pend
\phantomsection
\addcontentsline{toc}{subsection}{\textit{Porro quae vere vidua est ac desolata, sperat in Deo, et perseuerat in orationibus et obsecrationibus noctu et die. Quae autem in delitus versatur, viuens mortua est. Et haec praecipe, vt irreprehen sibiles sint. Quod si quis suis, et maxime familiaribus non prouidet, fidem abnegauit: et est infideli deterior. }}
\subsection*{\textit{Porro quae vere vidua est ac desolata, sperat in Deo, et perseuerat in orationibus et obsecrationibus noctu et die. Quae autem in delitus versatur, viuens mortua est. Et haec praecipe, vt irreprehen sibiles sint. Quod si quis suis, et maxime familiaribus non prouidet, fidem abnegauit: et est infideli deterior. }}\pstart Quae vidua est ac desolata. Explicatius quam antea, loqui tur:quia vere eas demum viduas esse docet, quae solitariae sint, et familia orbatae. Tales in Domino sperare dicit:non quod omnes id faciant, vel solae. Plaerasque enim viduas cernere licet orbas ac ont ni cognatione destitutas, quae tamen et superbae sunt et proteruae, et tam animo quam vita prorsus impiae. Et qui progenie abundant, non impediunpur quominus spem habeant in Deo reposiram. Vt Iob, et Iacob, et Dauid. Alioqui maledicta esset πολυτεανί α, quam scriptura passint inter eximias Dei benedictiones recenset. Verum  \pend
\section*{I. TIMOTH. }
\marginpar{[ p.58 ]}\pstart Paulus hoc loco viduas in Deo sperare dicit, quemadmodum septimo prioris ad Corinthios capite scribit, coelibes tantùm studere, vt Deo placeant: quia non sint instar coniugatorum diuisi. Sensus est igitur: nullo eas auocamento distrahi, quin respiciant in solum Deum: quia in mundo nihil reperiunt, quo fulciantur. Hoc argumento ipsas commendat. Nam vbi deficiunt humanae opes omniaque praesidia, Ecclesiae officium est manum porrigere ad opem ferendam.  \pend\pstart Perseuerat in orationibus.) Secundum argumentum commendationis: quia assidue precibus incumbunt. Vnde sequitur, iuuandas esse, ac sustentandas Ecclefiae impensis. Interea his duabus notis discernit inter dignas et indignas. Perinde enim valent haec verba, ac si tantum eas recipi praeciperet, quae nihil expectant opis ab hominibus, sed pendent a solo Deo: et depositis aliis curis et occupationibus, intentae sunt ad precandum. Alias minime idoneas esse aut commodas. Haec quoque precandi assiduitas requirit ab aliis curis vacationem. Nam quae circa familiam regendam occupantur, minus habent libertatis et ocii. Omnes quidem iubemur assidud̀ orare. Sed videndum, quid ferat vniuscuiusque conditio, dum ad orandum secessus quaeritur, et ab omnibus aliis negociis immunitas. Quod in viduis laudar Paulus, Lucas de Hanna filia Phanuelis praedicat. Atqui non idem omnibus conueniret: quia diuersa esset vitae ratio. Erunt ineptae muliereulae, Hannae simiae, magis quam imitatrices, quae circum altaria cursitabunt, et vsque ad meridiem susurris suis ac murmuribus dabunt operam. Hoc praetextu ab omnibus domesticis molestiis se eximent. Reuertae domum, nisi omnia ad nutum composita offenderint, insanis clamoribus exagitabunt totam familiam, prosilient interdum ad verbera. Meminerimus ergo, viduis et orbis non fine causa hoc quasi proprium tribui, vacare dies et nostes precibus, quia solutae sunt a iustis impedimentis, quae non idem permitterent familiam regentibus. Neque tamen monachis aut monialibus suffragatur haec sententia, qui murmura sua vel boatus ita vendunt, vt ociosi alantur. Quales olim fuerunt ἰυχοται, vel Psalliani. Nihil enim ab Euchitis differunt monachi et sacrifici papales, nisi quod illi, continenter precando, solos se pios ac sanctos putabant: hi aurem minori assiduitate se alios quoque sanctificare fingunt. Nihil tale in mentem venit Paulo: sed tantùm indicare voluit, quanto liberius vacare queant precibus, quae nihil aliunde habent distractionis.  \pend\pstart Quae autem in delitiis.) Postquam sua nota veras viduxs insigniuit, nunc alias opponit, quae reiiciendae sint. Participium, quo vtitur οπαταλωσα, significat eam, quae sibi indulget, moliterque ac delicate viuit. Quare eas, meo iudicio, taxat Paulus,  \pend
\section*{CAPVT V. }
\marginpar{[ p.59 ]}\pstart quae viduitate sua in hoc abutuntur, vt solutae iugo maritali, et omni molestia vacuae, suauem in ocio vitam degant. Multas enim videmus libertati suae consulere et commoditati, seque nimium hilariter tractare. Tales Paulus viuendo mortuas esse pronunciat. Quod aliqui de infidelitate intelligunt:cui opinioni minime assentior. Mortuam enim congruentius dici puco, quae inutilis sit ac nihili. Quorsum enim viuimus, nisi vt frugi sit nostra opera? Et quid, si esset emphasis in participio, viuens? Nam qui desidem vitam appetunt, que viuant commodius: habent illud in ore dicterium: Non est viuere. sed valere vita. Ita sensus esset :si bearae sibi videntur, dum agunt ex animi sui voto, atque hanc demum vitam reputant, in quiete delitiarisego mortuas esse pronuncio. Quia ramen sensus hic nimium forte argutus esset: tantum volui obiter attingere, nihil asserens. Cer tùm quidem est, ignauiam hic damnari a Paulo,vt mortuas vocet, quae nulli sunt vsui.  \pend\pstart Et haecpraecipç Significat, non tantùm se praescribere Timotheo quid sequi debeat: sed diligenter admonendas esse ipsas mulieres, ne eiusmodi viriis laborent. Neque enim se tantùm pastor opponere debet prauis studiis aut ambitioni eorum qui se praeter rationem ingerunt: sed doctrina et continuis admonitionibus omnia pericula, quantum in se est, praeuenire.  \pend\pstart Quod si quis.) Erasmus transtulit in foeminino genere. Mihi generalis sententia magis probatur. Solet enim Paulus, etiam cum aliquod particulare argumentum tractat, probationes ex generalibus principiis afferre: et rursus elicere ex particularibus sententiis vniuersalem doctrinam. Et certe plus habebit hoc ponderis, si conueniat tam viris, quam mulieribus. Dicit ergo eos fidem abnegasse, qui suos omnes, maximeque domum propriam non curant. Et merito. Nulla enim in Deum est pietas, vbi quis ita humanitatis sensum exuere potest. An fides, quae filios Dei nos facit, deteriores redderet brutis pecudibus? Manifestus ergo Dei contemptus est talis inhumanitas, et fidei abnegatio. Necdum eo contentus Paulus crimen exaggerat, dicens, infideli esse peiorem, qui suorum obliuiscitur. Quod duabus de causis verum est. Nam quo plus quisque in cognitione Dei profecit, ed minus habet excusationis. Ergo infide Iibus sunt peiores, qui in clara Dei luce caecutiunt. Deinde hoc genus officii est, quod natura ipsa dictat. Sunt enim στοργαι βυσικα. Quod si natura duce infideles vltro propensi sunt ad suos amandos: quid de iis sentiendum, qui nullo tali affectu tangunturt Nónne impios ipsos feritate superant? Quaeritur autem, cur domesticos filiis ipsis praeferat Apostolus. Respondeo, cum suos, et maxime domesticos nominat, vtroque loco de filiis et nepotibus ipsum loqui. Nam vtcunque  \pend
\section*{I. TIMOTH. }
\marginpar{[ p.60 ]}\pstart emancipati sint liberi, aut in alienam familiam per coniugium transierint, aut quouis modo emigrarint ex aedibus parentum: nunquam tamen prorsus extinguitur ius naturae, quin debeant maiores natu tanquam sibi a Deo commissos, regere, aut saltem consulere, quoad possunt. In domesticis autem arctior est obligatio. Nam eos duplici nomine curare debent, eo quod ipsorum sunt sanguis, et simul pars familiae cui praesunt.  \pend
\phantomsection
\addcontentsline{toc}{subsection}{\textit{Vidua deligatur non minor annis sexaginta quae fuerit vnius viri vxor, in operibus bonis habens testimonium, si liberos educauit, si fuit hospitalis, si sanctorum pedes lauit, si afflictis subministrauit, si in omni bono opere fuit assidua. Porro iuniores viduas reiice. Cumenim lasciuire coeperint aduersus Christum, nubere volunt, habentes condemnationem, quod primam fidem reiecerint. Simul autemt et ociosae discunt circuiredomos: nec folum ociosae, verùm etiam garrulae et curiosae, loquentes quae non oportet. }}
\subsection*{\textit{Vidua deligatur non minor annis sexaginta quae fuerit vnius viri vxor, in operibus bonis habens testimonium, si liberos educauit, si fuit hospitalis, si sanctorum pedes lauit, si afflictis subministrauit, si in omni bono opere fuit assidua. Porro iuniores viduas reiice. Cumenim lasciuire coeperint aduersus Christum, nubere volunt, habentes condemnationem, quod primam fidem reiecerint. Simul autemt et ociosae discunt circuiredomos: nec folum ociosae, verùm etiam garrulae et curiosae, loquentes quae non oportet. }}\pstart Vidua deligatur.) Iterum quales viduae recipiendae sint in Ecclesiae curam, ostendit:et quidem explicatius, quam antea. Ac primo aetatem designat, sexaginta annos. Nam quia publicis sumptibus alebantur: decebat iam senio confectas esse. Deinde erat altera etiamnum validior ratio. Nam se in ministerium Ecclesiae consecrabant. Quod minime ferendum erat, si adhuc ad coniugium ido neae fuissent. Ea conditione recipiebantur, vt eorum inopiae subueniret Eeclesia: et suam interim operam pauperibus impenderent, quoad patiebatur valetudo. Ita mutua erat obligatio inter ipsas et Ecclesiam. Robustiores, et quae per aetatem adhuc erant vegetae, iniquum erat aliïs esse oneri. Praeterea timendum erat, ne mutato animi proposito, ad nouas nuptias adspirarent. Hae sunt duae causae, cur non velit, nisi sexagenarias recipi. Quoad nubendi cupiditatem, satis occursum est periculo, vbi annum sexagesimum excessit mulier. Si praesertim vnum duntaxat maritum tota vita experta fuerit. Est enim quoddam veluti pignus continentiae et pudoris, cum mulier ad eam vsque aetatem peruenit vnico marito contenta. Non quod secundum coniugium improbet aut ignominiae notam aspergat bis nuptis. Quin potius iuniores viduas ad nubendum hortatur. Sed  \pend
\section*{CAPVT V. }
\marginpar{[ p.61 ]}\pstart quia solicite cauere voluit: ne vllis necessitatem imponeret coelibatus, quae maritis opus haberent. Qua de re mox plura.  \pend\pstart In operibus bonis.) Quae sequuntur dotes, partini ad hono rem, partim ad onus respiciunt. Non dubium est, quin fuerint honesta viduarum collegia, plenaque dignationis. Ideo non vult Paulus illuc cooptari, nisi quae ornatae sint egregio testimonio totius anteactae vitae. Iam non ad segne et ignauum ocium destinabantur: sed vt pauperibus et aegrotis ministrarent, donec viribus defectae, tanquam emeritae quiescerent. Itaque vt ad obeundam functionem sint aptiores, in officiis omnibus, quae sunt illius propria, longo vsu exercitatas esse vult. Qualia sunt, labor et sedulitas in liberis educandis, hospitalitas, ministeria erga pauperes, et reliqua beneficentiae opera. Si quis hic interroget: Ergô ne repudiabuntur omnes steriles, quia nullos vnquam liberos sustulerint? respondendum, non ste rilitatem hic damnari a Paulo, sed matrum delitias, quae sobolis alendae taedia deuorare dum recusant, satis ostendunt se minus officiosas fore erga extraneos. Et simul hanc piis matronis honorariam mercedem pro ponit, quae sibi non pepercerint, vt vicissim aetas earum senilis Ecclesiae gremio fouenda excipiatur. Per lotionem pedum, synecdochicons obsequia intelligit, quae sanctis praestari solent. Moris enim tunc erat pedes lauare. Hoc officii genus in speciem sordidum et prope seruile videri poterat. Itaque hac quasi nota insignit sedulas et minime delicatas aut molles foeminas. Quod proxime sequitur, ad liberalitatem pertinet. Atque idem postremo generalius exprimit, cum dicit, si in bono opere fuit assidua. Nam hic de beneficentia loquitur.  \pend\pstart Iuniores viduas reiice.) Non iubet excommunicari ab Ecclesia, aut aliqua ignominia notari. Sed tantùm in illo honore, cuius supra meminit, impertiendo, habendam esse earum rationem ne gat. Quod si Dei Spiritus per os Pauli mulierem nullam sexagenaria minorem ordine illo dignam esse pronunciat: quod in illa aeta te periculosus sit coelibatus: cuius audaciae fuit legem perpetui coelibatus iuuenculis, in summo aetatis feruore, postea imponere? Non admittit, inquam, coelibatum Paulus, nisi in aetate iam effoeta, et quae perieulum omne incontinentiae iam euasit. In velandis virginibus ad quadraginta deinde annos ventum est, postea ad triginta. Tandem promiscue, sine vlla exceptione aetatis, velare omnes coeperunt. Obtendunt, continentiam virginibus quae nunquam virum expertae sunt, multo faciliorem esse quam viduis. Sed non efficient, quin timendum sit periculum, quod hic Paulus cauet, ac cauendum prae cipit. Temere itaque adeóque crudeliter laqueum iniiciunt teneris adhuc puellis, quae coniugio futurae erant aptiores.  \pend
\section*{I. TIMOTH. }
\marginpar{[ p.62 ]}\pstart Cum enim lasciuire.) Lasciuire eas aduersus Christum dieit, quae conditionis ad quam vocatae sunt oblitae, hilarius, quam decet, exultant. Debuerant enim se continere sub modestiae iugo: vt grauibus matronis dignum est. Proinde niollior et solutior vita, quasi lasciuia est aduersus Christum, cui fidem suam obstrinxerant. Quoniam aliquot eiusmodi exempla viderat Paulus: generali remedio occurrit: ne qua recipiatur, quam aetatis ratio vnquam solicitare possit ad nubendi appetitum. Quot flagitiorum monstra quotidie in papatu coactus ille monialium coelibatus parit? Quae claustra non refringit libido? Ergo etiam si res initio placuisset: tot tamen atque tam grauibus experimentis edocti, aliquando Pauli consilio obtemperare debuerant. Atqui hoc adeo non faciunt, vt in dies magis ac magis iram Dei sua obstinatione prouocent. Nec de solis monialibus loquor. Ad perpetuum coelibatum sacerdotes quoque et monachos cogunt. Ebulliunt interea foedae libidines, vt vix decimus quisque caste se habeat. In monasteriis autem minimum malorum est vulgaris scortatio. Si Deum per os Pauli loquentem audire in animum inducerent: mox confugerent ad hoc remedium, quod praescribit. Sed tanta est superbia, vt furiose persequantur, quicunque huc eos reuocant. Aliqui particulam, aduersus Christum, coniungunt cum verbo nubendi. Quanquam autem ad Pauli memtem non multum refert: prior tamen sensus est probabilior.  \pend\pstart Habentes condemnationem.) Habere damnationem aliqui expo nunt, reprehensionem mereri. Ego tamen seuerius dictum accipio. vt aeternae mortis damnatione eas terreat Paulus. Ac si hoc ipsum reprehenderet Paulus, eximium illum ordinem, qui eas propius Chri sto vnire debuerat, illis damnationis esse materiam. Et ratio additur, quia a fide baptismi et Christianismo prorsus deficiant. Scio ab aliis secus intelligi. Nempe quod fallant datam Ecclesiae fidem nubendo:cum prius essent pollicitae, se ad mortem vsque coelibes victuras. Quod nimis est frigidum. Deinde, cur vocaret primam fidem? Vehemetius itaque insurgit Paulus, et mali grauitatem amplificat, quod non tantùm ignomi nia Christum et Ecclesiam afficerent resiliendo a suscepta conditione, sed primam quoque fidem impia defectione fallerent. Ita enim fieri solet, vt qui semel pudoris fines transiliit, ad omnem impudentiam se prostituat. Male eas habebat, probrosam esse suam leuitatem apud pios, et suam salacitatem reprimi, aut saltem vituperationi esse obnoxiam. Haec illis proteruius exultandi erat occasio, donec Christianismum abiicerent. Optime quadrat ista amplificatio. Quid enim magis absurdum, quam dum volumus personis consulere, ianuam aperire Christi abnegationit Quod votum perpetui coelibatus hinc stabilite volunt Papiste,  \pend
\section*{CAPVT V. }
\marginpar{[ p.63 ]}\pstart friuolum est. Vt demus, solitant fuisse disertis verbis interponi talem pactionem id tamen nihil ad ipsos. Primùm consyderandus est finis. Non ideo coelibes se fore promittebant olim viduae, vt sanctiorem agerent vitam, quam in coniugio. Sed quod non poterant maritis et Ecclesiae simul esse addictae. Sub Papatu autem vouetur continentia, quasi virtus per se accepta esset Deo. Deinde ea aetate se abdicabant tunc nubendi libertate, qua desierant esse nubiles. Sexagenarias enim ad minimum esse oportebat: et quae vnico matrimonio conten tae, specimen iam castitatis suae dedissent. Nunc autem vota in hoc eduntur apud Papistas, vt coniugiis, vel ante tempus, vel in medio aetatis feruore, renuncient. Duabus autem maxime de causis tyrannicam de coelibatu legem improbamus quod scilicet cultum esse fin gunt coram Deo meritorium: et quod vouendi temeritate praecipitant animas in exitium. Neutrum erat in veteri instituto. Non edebant directum continentiae votum: quasi minus Deo placeret vita coniugalis: sed tantùm, quatenus exigebat functio ad quam eligeban tur, coniugii vinculo se in totam vitam soluebant. Nec se priuabant nubendi libertate, nisi cum nubere, vel maxime liberis, iam absurdum erat, et intempestiuum. Denique tantùm a monialibus differebant istae viduae, quantùm Hanna prophetissa a Claudia vestali.  \pend\pstart Simul et ociosae.) Nihil est, quod mulieres magis deceat, quam domus custodia. Vnde apud veteres, Testudo honestae et probae4 matrisfamilias imago fuit. Atqui multae contrario morbo laborant. Nihil enim est, quod magis eas delectet, quam huc illuc cursitandi licentia. Praesertim vbi exoneratae familia, nihil habent domi quod agant. Adde quod istis viduis, honoris praetextu, quod veluti publicam personam gerebant, facilior quouis aditus parebat. Hanc opportunitatem nactae, beneficio Ecclesiae abutebantur ad desidiam. Deinde, vt fieri solet, ex ocio nascebatur curiositas: quae ipsa garrulitatis est mater. Verissimum enim est illud Horatii: Percontatorem fugito: nam garrulus idem est. Omni enim fide curiosos, vt inquit Plutarchus, carere aequum est: qui simul atque aliquid hauserunt, nunquam cessant, donce effutiuerint. Praesertim mulieribus hoc contingit, quae natura iam propensae sunt ad loquacitatem, nulliusque arcani capaces. Ergo non abs re haec tria simul coniuncta sunt a Paulo: ocium, curiositas, et garrulitas.  \pend
\phantomsection
\addcontentsline{toc}{subsection}{\textit{Volo igitur iuniores nubere, liberos gignere, domum administrare, nullam occasionem dare aduersario, vt habeat maledicendi causam. Nonnullae enim iam deflexerunt post Sathanam. Quod }}
\subsection*{\textit{Volo igitur iuniores nubere, liberos gignere, domum administrare, nullam occasionem dare aduersario, vt habeat maledicendi causam. Nonnullae enim iam deflexerunt post Sathanam. Quod }}
\section*{I. TIMOTH. }
\marginpar{[ p.64 ]}\pstart Volo iuniores nubere.) Rident nasuti homines hoc Apostoli praeceptum. Quasi verô, inquiunt, calcar addendum sit nimis citatae earum cupiditati. Quis enim nescit omnes fere viduas sponte nupturire? Superstitiosi autem homines iudicarent hanc nuptialem doctrinam Christi Apostolo parum congruere. Verùm, omnibus rite expensis, agnoscent sani homines, Paulum nihil hic, nisi necessarium, et maxime salutare tradere. Nam et viduitas multis liberius proteruiendi dat facultatem: et ab altera parte subinde emergunt spiritus in hypocrisi mentientes, qui sanctimoniam constituunt in coeIibatu: quasi angelica esset perfectio: matrimonium vero, aut dannant in totum, aut vilipendunt, acsi carnis inquinamentum saperet. Pauci sunt, tam ex viris, quam mulieribus, qui vocationem suam respiciant. Quotusquisque est virorum, qui onus regendae vxoris liben ter suscipiat? Quia scilicet res est innumeris molestiis implicita. Quam aegre iugum subit mulier? Paulus ergo, cum iuniores nubere iubet, non ad nuptialem laetitiam eas inuitat: nec hortatur vt libidini indulgeant, cum iubet eas liberos gignere: fed dum sexus infirmitatem, et aetatem lubricam reputat, eas ad castum coniugium hortatur: et simul ad subeunda sandi coniugii onera. Atque id praesertim facit, ne contumeliosa videatur ipsarum reiectio ab ordine viduarum. Significat enim, earum vitam non minus Deo probatum iri, si oeconomiae dent operam, quam si in viduitate maneant. Et sane, Deus superstitiosas hominum opiniones nihil moratur: sed hanc obedientiam pluris aestimat quam omnia, vbi suae vocationi obsequimur potius, quam ferri nos sinamus animi nostri voto. Hac consolatione audita, nihil habent causae, cur conquerantur sibi fieri iniuriam, aut moleste ferant, se excludi vno genere honoris. Audiunt enim in vita coniugali non minus se placere Deo, quia eius vocationi pareant. Cum de liberis gignendis loquitur, molestias omnes sub vno verbo comprehendit, quae perferendae sunt in educanda sobole. Quemadmodum, sub familiae gubernatione, omnia oeconomiae officia.  \pend\pstart Nullam occasionem dare.) Nam sicuti maritus vxoris quasi velamen est: ita viduitas multis sinistris superstitionibus est obnoxia. Quorsum autem attinet, nulla necessitate calumniis armare Euangelii hostes? Atqui difficilimum est, mulierem viduam, florenti aetate, tam circunspecte sibi cauere, quin reperiant improbi, aliquem obloquendi colorem. Proinde si illis est cordi aedificatio, vt maledicis  \pend
\phantomsection
\addcontentsline{toc}{subsection}{\textit{}}
\subsection*{\textit{}}
\section*{CAPVT V. }
\marginpar{[ p.65 ]}\pstart hominibus os obstruant, genus vitae eligant minus suspitiosum. Com munes hic Euangelii aduersarios potius intelligo, quam priuatos cu iusque mulieris.  \pend\pstart Nonnullae enim.) Certum est, nullum esse tam sanctum institutum, ex quo non aliquid prodeat mali, hominum prauitate. Verùm quae necessaria sunt, debent fixa manere, quidquid inde contingat: etiam si coelum rueret. At vbi liberum est, vtramuis partem eligere: prout hoc vel illud esse vtile, experientia compertum fuerit, etiam quae ante placuerant, abrogare, prudentiae est. Sicuti ircausa praesenti. Minus prouectae aetatis mulieres in ordinem viduarum allegi, nihil necesse erat. Vsus docebat, rem esse periculosam et noxiam. Meritoitaque Paulus in posterum caueri vult, ne quid tale contingat. Quod si quarundam defectio satis validum illi argumentum fuit, vt remedium vniuersale quaereret: quam multa hodie Papistis argumenta essent abolendi impuri sui coelibatus, si quem aedificationis haberent respectum? Sed crudelibus impii ac Diabolici decreti laqueis malunt strangulari innumeras animarum myriades, quam nodum vnum laxare. Vnde constat, quam aliena sit a sancto Pauli affectu tyrannica eorum saeuitia.  \pend\pstart Post sathanam.) Notanda locutio: quia nemo vel tantillum potest a Christo deflectere, quin Sathanam sequatur. Nam regnum in omnes habet qui Christi non sunt. Hinc admonemur, quam exitialis sit deflexio a recto cursu.  \pend\pstart Si quis fidelis.) Quia fieri solet, vt singuli propriae onera libenter reiiciant in totam Ecclesiam: ideo nominatim id cauendum praecipit. De fidelibus loquitur, qui suas viduas alere debeant. Nam quae impiae cognationi renunciauerant, eas ab Ecclesia suscipi, aequum erat. Quod si peccant, qui dum sibi pareunt, Ecelesiam patiuntur sumptibus grauari: hinc colligamus, quanto sacrilegio se implicent, qui fraude aut rapinis profanant, quod semel dicatum erat Eccleliae.  \pend
\phantomsection
\addcontentsline{toc}{subsection}{\textit{Presbyteri, qui bene praesunt, duplici honore digni habeantur. Maxime qui laborant in verbo et doctrina. Dicit enim Scriptura: Non obligabis os boui trituranti. Et, Dignus est operarius mercede sua. Aduersus presbyterum accusationem ne admittas, nisi sub duobus aut tribus testibus. Peccantes coram omnibus argue: vt et caeteri timorem habeant. Contestor coram Deo, et Domino Iesu }}
\subsection*{\textit{Presbyteri, qui bene praesunt, duplici honore digni habeantur. Maxime qui laborant in verbo et doctrina. Dicit enim Scriptura: Non obligabis os boui trituranti. Et, Dignus est operarius mercede sua. Aduersus presbyterum accusationem ne admittas, nisi sub duobus aut tribus testibus. Peccantes coram omnibus argue: vt et caeteri timorem habeant. Contestor coram Deo, et Domino Iesu }}
\section*{I. TIMOTH. }
\marginpar{[ p.66 ]}\pstart Christo, et electis Angelis, vt haec custodias absque praecipitatione iudicii, nihil faciens alteram in partem declinando.  \pend\pstart Hoc quoque ad retinendum Ecclesiae ordinem in primis est necessarium: non negligi presbyteros, sed eorum haberi rationem. Quid enim inhumanius, quam eos non esse curae, qui totius Ecclesiae curam gerunt? Neque enim hic aetatis nomen est πρεσβύτεροι. sed officii. Quod duplicem honorem interpretatur Chrysostomus victum et reuerentiam, non impugno. Sequatur eius opinionem, qui volet. Mihi tamen probabilius videtur, hic fieri comparationem inter viduas et presbyteros. Prius iusserat Paulus, honorem haberi viduis. Atqui honore digniores sunt presbyteri, quam illae. Quare his duplex honor, earum respectu, exhibendus. Sed vt ostendat, se laruas non commendare, addit, Qui bene praesunt: hoc est, qui fideliter et strenue incumbunt suo muneri. Nam vt centies quis locum obtineat, et titulum iactet, nisi munus simul exerceat, nullo iure stipendiis Ecclesiae se ali postulabit. Denique significat non titulo deberi honorem, sed operae quam impendunt, qui sunt in officio constituti. Praefert tamen eos, qui in verbo et doctrina laborant: hoc est, qui docendo verbo sunt intenti. Nam duae istae voces rem vnam significant. nempe verbi praedicationem. Sed ne verbi nomine ociosum stu dium, et speculatiuum, vt vocant, quispiam intelligeret, adiecit doctrinam. Colligere autem hinc licet, duo fuisse tunc presbyterorum genera: quia non omnes ad docendum ordinabantur. Nam aperte verba sonant, quosdam bene et honeste praefuisse, quibus tamen non erant corumissae docendi partes. Et sane, ex populo deligebantur graues et probati homines, qui vna cum pastoribus communi consilio et authoritate Ecclesiae disciplinam administrarent: ac essent quasi censores moribus corrigendis. Hunc morem Ambrosius obsoleuisse conqueritur Doctorum ignauia, vel potius superbia, dum soli voIunt eminere. Nunc, vt ad Paulum redeam, victum praecipue suppeditari iubet pastoribus, qui docendo sunt occupati. Ea enim est mundi ingratitudo, vt de fouendis verbi ministris non multum sit solicitus. Et Sathan hac arte tentat doctrina priuare Ecelesiam, dum inopiae et famis metu plurimos absterret, ne id oneris suscipiant.  \pend\pstart Non obligabis os.) Praecepsum est politicum, quod nobis in genere aequitatem commendat: sicuti in priorem ad Corinthios Epistolam diximus. Nam si prohibet, ne inhumani simus brutis animalibus: quantoplus humanitatis exigit erga ipsos homines? Ergo perinde valet haec sententia, acsi generaliter diceretur, non esse labo\pend
\phantomsection
\addcontentsline{toc}{subsection}{\textit{}}
\subsection*{\textit{}}
\section*{CAPVT V. }
\marginpar{[ p.67 ]}\pstart te alieno abutendum. Triturandi mos hodie multis Galliae partibus est ignotus, vbi flagellis triticum excutiunt. Soli Prouinciales norunt quid sit triturare. Sed hoc nihil ad sensum. Quia tantundem de aratione dici potest. Quod subiieit continuo post, operarium dignum esse mercede sua, non citat quasi Scripturae testimonium: sed quasi di ctum prouerbiale, quod omnibus dictat communis sensus. Quemadmodum et Christus cum idem dicebat Apostolis, nihil aliud quam sententiam proferebat omnium consensu approbatam. Ergo crudeles sunt et omnis aequitatis immemores, qui iumenta, nedum homines, esurire sinunt, quorum sudorem exugunt in suum commodum. Quantd aucem minus ferenda eorum ingratitudo, qui victum pastoribus suis negant, quibus nullam satis dignam mercedem possunt rependere?  \pend\pstart Aduersus presbyterum.) Postquam stipendia iussit erogarr pastoribus, Timotheum etiam admonet, ne sinat eos infamari calum niis: nec vlla criminatione grauari, nisi quae nitatur legitimo testimonio. Sed posset hoc videri absurdum, quod legem, quae omnium est communis, presbyteris quasi peculiarem facit. Hoc iuris in causis omnibus Deus praescribit: vt ex ore duorum aut trium testium definiantur. Cur ergo hoc quasi singulari priuilegio solos presbyteros munit Apostolus, vt a falsis delationibus tuta sit eorum innocentia? Respondeo, hoc esse necessarium remedium, aduersus hominum malignitatem. Nulli enim calumniis et obtrectationibus magis sunt obnoxii, quam pii Doctores. Praeterquam enim quod muneris difficultas facit, vt interdum, vel succumbant, vel titubent, vel claudicent, vel hallucinentur, vnde multas carpendi ansas arripiunt improbi: huc etiam accedit, quod quantumuis exacte fungantur suis partibus, vt ne minimum quidem erratum admittant, nunquam mil le repreliensiones effugiunt. Atque hic est astus Sathanae, alienare hominum animos a ministris, vt doctrina paulatim veniat in contemptum. Ita non tantùm fit iniuria innocentibus, quod immerito laeditur eorum existimatio ( quod tamen ipsum perquam indignum est in tam honesto ordine) sed detrahitur authoritas sacrosanctae Dei doctrinae. Atque id praecipue, vt dixi, molitur Sathan. Neque enim hic tantùm valet illud Platonis, Malignum esse vulgus ac praestantioribus inuidere. Sed quo sincerius quisque pastorum promouere studet regnum Christi, eo maiore oneratur inuidia, et proteruiores aduersus eum fiunt iusultus. Iam vero, simul atque aliqua de verbi ministris criminatio exiit, perinde acsi iam conuicti forent, certa fide excipitur. Id accidit, non tantùm quia maior ab illis exigitur integritas: sed quia plerosque fere omnes Sathan ad nimiam credulitatem instigat, vt nulla inquisitione habita, cupide damnent suos pa\pend
\section*{I. TIMOTH. }
\marginpar{[ p.68 ]}\pstart stores, quorum tueri famam potius debuerant. Ergo tantae iniquiteti non frustra se Paulus opponit: ac prohibet subiici improborum ma ledicentiae presbyteros: donec legitimis testimoniis sint conuicti. Ne que enim mirum est, si multos habeant hostes, quorum officium est, omnium vitia reprehendere, prauis omnium cupiditatibus aduersari, seueritate sua cohibere quoscunque errare vident. Quid ergo fiet, si promiscue audientur, quaecunque de illis sparsae fuerint calumniae?  \pend\pstart Peccantes argue coram omnibus.) Quoties bonis cauetur, hoe mali ad immunitatem suam arripiunt. Itaque quod de repellen dis iniustis accusationibus dixerat Paulus, sic temperat hac sententia, ne cui tamen impune hoc praetextu peccare liceat. Et sane videmus, quantis et quam multiplicibus priuilegiis suum clerum munierit Papatus: vt cum nequissime viuant, sint tamen ab omni reprehensione immunes. Certe si cautiones obscruentur quas colligit Gratianus causa secunda, quaest. 4. item, q. 7. non erit periculum, ne vnquam ad reddendam vitae rationem cogantur. Vbi enim sepruaginta duo restes reperientur, quos in Episcopo damnando requirit pu tidum edictum sub nomine syluestri? Deinde cum totus laicorum ordo ab accusatione arceatur: in clero etiam minores maioribus prohibeantur molestiam facessere: quid impedit, quominus securi iudicia omnia rideant? Proinde sedulo haec moderatio tenenda est: vt coerceantur linguae petulantes, ne fictis criminibus presbyteros infament: et tamen vt quisquis inter cos male se geret, seuere corrigatur. Nam de presbyteris hoc dictum accipio: palam arguendos esse, qui dis solute viuunt. Quare? Vt reliqui, dum ne illis quidem parci vident, qui gradu et dignitate sunt superiores, tali exemplo moniti magis ti meant. Nam sicuti honestae vitae exemplo praeire aliis debent presbyteri:ita si delinquunt, seueritatem disciplinae exerceri erga ipsos conuenit, quae omnibus sit exemplo. Cur enim plus illis indulgeatur licentiae, quorum peccata longe grauius, quam aliorum nocent? IntelIige autem quod dicit Paulus de sceleribus, aut flagitiosis peccatis, quae publicum offendiculum secum ferunt. Nam si quis ex presbyteris absque flagitio deliquerit :priuatim magis admonendum esse certum est, quam palam redarguendum.  \pend\pstart Contestor coram Deo.) Haec obtestatio interposita fuit a Paulo, non tantùm vt in re maxime seria: sed etiam propter summam difficultatem. Nihil enim difficilius, quam tanta aequitate exer cere iudicia, vt neque gratia vnquam flectaris, nec suspitionibus des aditum, nec delationibus mouearis, nec rigore sis nimius: nec aliud in omnibus causis quam ipsam causam respicias. Nam inde ius aequabile, cum personis clausi sunt oculi. Meminerimus autem, hic sub Timothei persona admoneri omnes pastores, et Timotheum armari  \pend
\section*{CAPVT V. }
\marginpar{[ p.69 ]}\pstart veluti clypeo, aduersus prauas omnes solicitationes, quae non raro plurimum negocii solent, vel optimis quibusque exhibere. Timotheo igitur Deum ponit ante oculos: vt sciat, non minori religione sibi administrandum esse suum munus, quam si ageret in conspectu Dei et Angelorum. Postquam Deum nominauit, addit etiam Christum. Nam is est, cui Pater omnem iudicandi potestatem detulit: et coram cuius tribunali semel apparebimus. Angelos Christo adiungit, non tanquam iudices, sed tanquam futuros testes negligentiae, aut temeritatis, aut ambitionis, aut malae fidei. Adsunt enim spectatores, propter iniunctam sibi Ecclesiae curam. Et sane plusquam stupidum et saxeum esse oportet, cui non excutiat torporem et oscitan tiam sola haec cogitatio: theatrum Dei et Angelorum esse Ecclesiae gu bernationem. Vbi vero accedit solennis ista contestatio, timorem ac solicitudinem duplicari necesse est. Electos vocat Angelos, non tantùm vt a reprobis discernat sed excellentiae causa, vt plus reuerentiae, habeat eorum testimonium.  \pend\pstart Absque praecipitatione.y πρόκριμα Graecis, si verbum de verbo exprimas, Latine sonat praeiudicium. Verdm potius significat festinationem nimiam, dum re nondum bene perspecta, temere pronunciamus: aut significat immodicum fauorem, cum deferimus personis plusquam aequum est, vel alios, tanquam eximios, aliis praeferimus. Quod in iudiciis semper est iniquum. Ergo hic vel leuitatem, vel personarum acceprionem damnat Paulus. Eodem spectat fere, quod sequitur: non esse huc vel illuc declinandum. Dici enim vix potest, quam difficulter se retineant qui iudicandi munus sustinent, inter tot ac tam varios impulsus. Nonnulli codices habent, π gόοολνσιη. Sed prior lectio magis conuenit.  \pend
\phantomsection
\addcontentsline{toc}{subsection}{\textit{Manus cito ne cui imponas: neque communices peccatis alienis. Temet ipsum purumcustodi. Ne posthac bibas aquam: sed paululo vino vtere, propter stomachum tuum, et crebras tuas infirmitates. Quorundam hominum peccata ante manifesta sunt, festinantia ad iudicium: in quibusdam verô etiam subsequuntur. Simil’iter et bona opera ante manifesta sunt, et quae secus habent, latere nequeunt. }}
\subsection*{\textit{Manus cito ne cui imponas: neque communices peccatis alienis. Temet ipsum purumcustodi. Ne posthac bibas aquam: sed paululo vino vtere, propter stomachum tuum, et crebras tuas infirmitates. Quorundam hominum peccata ante manifesta sunt, festinantia ad iudicium: in quibusdam verô etiam subsequuntur. Simil’iter et bona opera ante manifesta sunt, et quae secus habent, latere nequeunt. }}\pstart Non dubium est, quin depellere a Timotheo inuidiam voluerit, et obuiam ire multis querimoniis, quae subinde oriuntur aduersuis pios Christi seruos, qui ambitiosis quorumnis precibus obsequi  \pend
\section*{I. TIMOTH. }
\marginpar{[ p.70 ]}\pstart recusant. Nam alii morositatis eos insimulant, alii inuidiae: nonnulli etiam crudeles esse clamitant: quia non statim admittant, qui se venditant qualieunque commendatione. Quod hodie plus saris experimur. Paulus ergo Timotheum hortatur, ne a grauitate discedat, nec se patiatur vinci praeposteris studiis: non tam quod opus haberet Timotheus eiusmodi admonitione, quam vt sua authoritate compescat, qui alias ipsi Timotheo molesti esse poterant. Primùm, impositio manuum ordinationem significat: hoc est, signum pro reipsa capitur. Prohibet enim, ne quem minus adhuc probatum, nimia facilitate admittat. Sunt enim qui nouitatis studio promoueri cuperent, vel igno tissimum quenque, simul atque specimen vnum vel alterum ediderit, quod arrideat. Huic importunae cupiditati grauem et prudentem Episcopum resistere conuenit: quemadmodum hic Paulus iubet Timotheum facere.  \pend\pstart Neque communices. Significat, eodem implicari reatu, qui illicitae promotioni assentitur, atque ipsos praecipuos authores. Quanquam ira exponunt quidam, si indignos admittat: quidquid postea deliquerint, culpam illi aut culpae partem imputatum iri. Sed mihi hoc videtur simplicius: etiam si in hanc temeritatem alii prorumpant: noli te iilis socium addere, ne simul reus fias. Saepe enim accidit, vt nos, etiam recte alioqui sentientes, abripiat aliorum stultitia et leuitas. Huc etiam refero, quod continuo post sequitur: Temetipsum purum custodi:acsi dixisset: Si quid perperam fiat ab aliïs, vide ne vlla ex consensu, vel approbatione contagio ad te perueniat. Si impedire non potes, quominus se polluant: saltem tua consilia ab illis constanter separata esse conuenit, quo purum te serues. Si quis generalem esse malit sententiam, fruatur suo sensu. Sed magis consentaneum est, meo iudicio, ad cireunstantiam praesentis loci restringere.  \pend\pstart Ne posthac bibas aquam.) Quoniam liaec sententia contextum abrumpit, quidam non esse Pauli suspicantur. Atqui videmus Paulum non adeo fuisse curiosum in texenda orationis serie: atque hoc illi non esse infrequens, vt diuersas sententias absque ordine permisceat. Deinde fieri potest, vt quod, scripta iam Epistola, appositum fuerat extra ordinem versuum, deinde librariorum errore in hune locum obrepserit. In eo tamen non ita anxie laborandum est, si illum, quem dixi, Pauli norem consyderamus. Nempe quod res varias interdum miscere solet. Porrd̀ summa est, vt Timotheus aliquantulum vini bibere assuescat, valetudinis sustentandae causa. Non enim prohibet in totum, ne bibat aquam: sed ne vtatur hoc ordinario potu. Nam verbum 'υ Αροποτ α̃ν hoc significat. Sed cur non simpliciter consulit, vt bibat vinum? Videtur enim, eum pauxillum  \pend
\section*{CAPVT V. }
\marginpar{[ p.71 ]}\pstart addit, velle obuiam ire intemperantiae: quod non fuit verendum in Timotheo. Respondeo, hoc potius fuisse expressum, vt occurreret improborum calumniis, quibus promptum alioqui fuisset eius consilium subsannare, hoc scilicet colore, aut simili: Qualis istaec est philosophia, quae ad potandum vinum incitat? Haecne via est, qua in coelum conscenditur? Vt eiusmodi sannas praeueniat Apostolus, testatur se necessitati dumtaxat consulere: et medioeritatem simul commendat. Apparet autem Timotheum non frugalem tantùm in victu, sed etiam austerum fuisse: vt qui ne valetudini quidem suae parceret. Ac certum est, neque ambitione id ab eo, neque superstitione fuisse factum. Vnde colligimus, non modo a luxu et delitiis fuisse alienissimum: sed quo solutior esset ad opus Domini, aliquid etiam ex ordinariis alimentis rescidisse. Neque enim natura, sed continentiae studio erat abstemius. At quam pauci hodie sunt, quibus necesse sit aqua interdicere: quam multi vero qui restringendi sint ad sobrium vini potumt  \pend\pstart Quorundam hominum peccata.) Quoniam nihil est, quod magis exeruciet fideles Ecclesiae ministros, quam vbi malis corrigendis nul lam vident rationem: coguntur ferre hypocritas, quorum nequitiam incelligunt: multos, qui sunt noxiae pestes, neque arcere ab Ecclesia valent: ac ne impedire quidem, quominus clandestinis artibus venenum suum spargant: hac consolatione sustentat Paulus Timotheum, fore vt aliquando, vbi Deo visum fuerit, protrahantur in medium. Ita ad tolerantiam eum confirmat: quia aequo animo expectanda est opportunitas, quam Deus, pro sua sapientia, praestituit. Mire etiam probos et sanctos pastores torquet altera indignitas: quod vbi summa integritate officio perfuncti sunt, multis iniquis sermonibus per stringuntur, onerantur multiplici inuidia, vident rapi in contrariam partem quae laudem merebantur. Hac etiam ex parte subuenit Paulus, eum admonet Timotheum, quibusdam bonis operibus lucem in aliud tempus reseruari. Itaque si laus eorum quasi sub terram defodiatur hominum ingratitudine, id quoque placido animo esse sustinendum, donec reuelationis tempus adueniat. Quanquam non his solùm malis medetur. Sed quia nos saepe falli contingit in eligendis ministris, dum se indigni astute insinuant, boni autem nobis sunt ignoti:deinde etiam si non erremus ipsi iudicando, nontamem aliis probare iudicium nostrum possumus, sed qui optimi sunt repudiantur nobis contra nitentibus, mali autem vel obrepunt, vel se obtrudunt: fieri non potest, quin nos magnopere angat nostra et Ecclesiae canditio. Itaque Paulus in offensione vel tollenda, vel saltem subleuanda laborat. Venio nunc ad verba, postquam rem breuiter complexus sum. Cum dicit, peccata quorundam ante manifesta esse intelligit mature  \pend
\section*{I. TIMOTH. }
\marginpar{[ p.72 ]}\pstart patefieri, et prodire in notitiam hominum. Et hoc similitudine exprimit, quasi currendo festinent ad suum iudicium. Videmus enim vt multi praecipites ruant, sponteque accersant sibi damnationem: etiam si totus mundus saluos cupiat.  \pend\pstart In quibusdam subsequuntur.) Quod Erasmus vertit, quosdam subsequuntur, non recipio. Tametsi videtur id Graecae constructioni magis quadrare: sesus tamen postulat, vt subaudiatur particula ’εν. Casus enim mutatio non tollit antithesin. Nam sicuti dixerat, quorundam peccata celeriter ad iudicium suum festinare: ita nunc econuerso addit, quorundam vel aliorum tardius venire in notitiam. Sed loco geniciui, posuit datiuumin aliis. Significat autem, vtcunque diutius, quam optemus, lateant quorundam peccata, et tarde proferantur in lucem, non tamen perpetuo fore occulta. Nam et ipsa tempus suum habebunt.  \pend\pstart Similiter et bona.) significat interdum pietatem et alias virtutes mature ac tempestiue suam laudem apud homines habere: vt in precio sint boni. Quod si secus accidat, non tamen passurum Dominum, vt perpetuo opprimatur innocentia et integritas. Sepae enim calumniis tanquam nubibus obscuratur: sed tandem impletur illud, Deum facere vt refulgeat instar aurorae. verum animo tranquillo opus est ad sustinendum. Homines enim sumus. Itaque semper videndum, quem modum habeat, nostra cognitio, ne vltra progredia mur. Nam id esset, Dei officium nobis arrogare.  \pend
\endnumbering\beginnumbering\section{CAPVT. VI.}
\phantomsection
\addcontentsline{toc}{subsection}{\textit{\huge\textbf{Q}\normalsize VICVNQVE sub iugo sunt serui: suos Dominos omni honore dignos exi stiment: vt ne Dei nomen et doctrina blasphemetur. Qui autem fideles ha'bent dominos, ne despiciant, eo quod fratres sunt. Sed magis seruiant, quod fideles sint et dilecti, et beneficent iae participes. }}
\subsection*{\textit{\huge\textbf{Q}\normalsize VICVNQVE sub iugo sunt serui: suos Dominos omni honore dignos exi stiment: vt ne Dei nomen et doctrina blasphemetur. Qui autem fideles ha'bent dominos, ne despiciant, eo quod fratres sunt. Sed magis seruiant, quod fideles sint et dilecti, et beneficent iae participes. }}\pstart Apparet, initio Euangelii, seruos, quasi dato manumissionis signo, animos sustulisse. Nam passim multuni laborat Paulus in cupiditate illa frenanda. Et sane, tam dura erat seruitutis conditio, vt non mirum sic, fuisse imprimis odiosam. Porro vsitatum est, arripere in carnis nostrae commodum quidquid habet minimam spe ciem. Ita cum audiebant, omnes nos esse fratres, protinus infere\pend
\section*{CAPVT VI. }
\marginpar{[ p.73 ]}\pstart bant indignum esse, vt fratrum essent mancipla. Verùm, vt nihil tale venisset illis in mentem: semper tamen miseris hominibus necessaria est consolatio, quae temperet aerumnarum acerbitatem. Deinde non sine diffieultate adduci poterant, vt libenter et alacri animo tam aspero iugo subiicerent ceruices. Huc igitur spectat praesens doctrina.  \pend\pstart Quicunque sub iugo.) Quia sibi quisque praestantiam falsa opinione arrogat: nemo est, qui aequo animo ferat, alios sibi imperare. Qui effugere necessicatem nequeunt, parent illi quidem inuiti superioribus. sed intus fremunt et indignantur, quia sibi putant fieri iniuriam. Omnes eiusmodi disputationes vno verbo praecidit Apostolus, cum voluntariam subiectionem exigit ab omnibus, qui sub iugo sunt. Significat enim, non esse inquirendum, sint ne digni tali for tuna an meliore: quia sufficiat, hac conditione esse obstrictos. Et cum praecipit, vt dignos honore existiment, quibus seruiant: non modo fidem et diligentiam in praestandis obsequiis requirit: sed vt eos suspiciant et venerentur ex animo, tanquam in altiore gradu collocatos. Nemo enim vel principi, vel domino praestat quod debet, nisi eminentiam inspiciens, ad quam illos Deus euexit, honore ideo prosequitur, quia illis subest. Nam; vtcunque saepe indigni sint: ipsa tamen praefectura, qua dignatur ipsos Deus, honorem semper me retur.  \pend\pstart Ne nomen Dei.) Semperin nostrum commodum, plusquam oporteret, ingeniosi sumus. Ita seruis, si infideles habeant dominos, prompta est obiectio, indignum esse, vt qui diabolo seruiunt, imperent filiis Dei. Paulus autem in contrariam partem retorquet argumentum, ideo infidelibus dominis parendum esse, ne male audiat nomen Dei et Euangelium: quasi Deus ille queni colimus, ad rebellandum nos incitet: quasi Euangelium contumaces reddat et praefra ctos, qui aliis subiecti esse debent.  \pend\pstart Qui aucem fideles.) Fraternum nomen videtur aequalitatem statuere, atque ita tollere dominium. Paulus exaduerso ratioci natur, eo libentius debere seruos fidelibus dominis se subiicere, quia filios Dei agnoscunt, et eiusdem secum gratiae consortes. Neque enim vulgaris est ista dignatio, quôd Deus eos dominis, in eo quod praecipuum est, aequauit. Nam communem simul adoptioneni habent. Hoc ergo non leue illis incitamentum esse debet, ad seruitu tem aequo animo ferendam.  \pend\pstart Haec doce et exhortare. Si quis aliter docet, * Aunec acquiescit sanis sermonibus Domini nostri Iesu  \pend
\vspace{0.5cm}\noindent
\vspace{0.2cm}\rule{1cm}{0.2pt}\\ 
\hspace{0.2cm}\textit{mg}
\footnotesize 
\normalsize| 
\section*{I. TIMOTH. }
\marginpar{[ p.74 ]}
\phantomsection
\addcontentsline{toc}{subsection}{\textit{Christi, et ei quae secundùm pietatem est doctrinae: inflatus est, nih il sciens, sed languens circa quaestiones et pugnas verborum: ex quibus oritur inuidia, con tentio, maledicentiae, suspitiones malae, superuacuae conflictationes hominum mente corruptorum, et qui veritate priuatisunt, existimantium quaestum esse pietatem. Seiungere a talibus. }}
\subsection*{\textit{Christi, et ei quae secundùm pietatem est doctrinae: inflatus est, nih il sciens, sed languens circa quaestiones et pugnas verborum: ex quibus oritur inuidia, con tentio, maledicentiae, suspitiones malae, superuacuae conflictationes hominum mente corruptorum, et qui veritate priuatisunt, existimantium quaestum esse pietatem. Seiungere a talibus. }}\pstart Haec doce.) Intelligit haec esse, quibus docendis insistere debeat. Et doctrinae vult exhortationes coniungi. Perinde ergo est, ac si dixisset, quotidie inculcandum esse hoc doctrinae genus. Nec docendos modo esse homines: sed exhortationibus quoque assiduis stimulandos ac vrgendos.  \pend\pstart Si quis aliter. Quia verbum est compositum, potest etiam non male transferri, alia docet. Quanquam in sensu nulla est ambiguitas. Damnat enim eos omnes, qui in hac docendi ratione non conueniunt: etiam si non aperce neque ex professo oppugnent sanam doctrinam. Fieri enim porest, vt qui nullam alioqui profitetur impiam doctrinam, futili tamen garrulitate se ventitando, pietatis doctrinam adulteret. Nam vbi in ipsa doctrina profectus nullus apparet, nullaque aedificario, iam a Christi instituto disceditur. Cum autem non de professis impiorum dogmatum authoribus loquatur Paulus, sed de vanis doctoribus ac profanis, qui ambitione sua vel auaritia, simplicem et genuinam pietatis doctrinam, deformant: videmus tamen quam acriter ac seuere inuehatur aduersus eos. Nec mirum. Dici enim vix potest, quantùm noceat fucosa, et prorsus ad ostentationent et inanem pompam composita praedicatio. Verdm quos perstringat, melius patet ex contextu. Nam quod proxime sequitur: et non acquiescit sanis sermonibus: exegeticum est. Fieri enim solet, vt eiusmodi homines, praepostera curiositate abrepti, quidquid vtile est ac solidum, despiciant. Itaque instar petulantium equorunnsine ratione lasciuiunt. Id verd̀ quid aliud est, quam respuere sanos Christi sermones? Nan sani vocantur ab effectu, quia sanitatem, nobis conferant. Idem significat doctrina, quae secundùm pietatem est. Neque enim pietati consentanea erit, nisi in timote cultuque Dei nos instituat, nisi fidem nostram aedificet, nisi ad patientiam, humilitatem, omnia charitatis officia nos eru diat. Ergo quisquis vtiliter docere non studet, secus docet quam oportebat. Imo nec pia, nec sana est doctrina, quocunque colore fulgeat, quae ad profectum audientium non tendit. Nulla excogitari po\pend
\section*{CAPVT VI. }
\marginpar{[ p.75 ]}\pstart test aptiot ad eastigandos ambitiosos poena, quam dum inscitiae dam natur totum id, quo sibi placent. Itaque Paulus eos nihil scire pronunciat, vtcunque multis argutiis turgeant. Nihil enim habent solidum, sed meros tantùm flatus Et simul pios omnes admonet, ne se ventosa ista ostentatione efferri patiantur, quominus in Euan gelii simplicitate defixi maneant.  \pend\pstart Sed languens.) Est obliqua antithesis inter sanitatem doctrinae Christi, ac languorem istum. Nam cum multùm ditque se fatigauerint subtiles isti quaestionarii: quem tandem fructum reportant ex suo labore? Nisi quod subinde crescit morbus. Ita non frustra tantùm se macerant: sed languorem contrahunt ex peruersa sua curiositare. Sequitur ergo longe abesse, vt rite proficiant, quemad modum oportebat Christi discipulos. Quaestiones et verborum pugnas non abs re simul coniungit. Priore enim vocabulo non quasuis quaestiones significat, quae vel ex sobrio moderatóque discendi studio nascuntur: vel ad puram rerum vtilinm explicationem faciunt: sed quales hodie in Sorbonicis scholis agitantur, ad ostentandum acumen. Illic quaestio quaestionem parit; quia nullus est modus, vbi quispiam plus scire, quam expediat, appetens, suae vanitati in dulget. Hine postea emergunt infinitae rixae. Quemadmodum enim crassae in calore nubes non sine tonitru discutiuntur ita necesse est, vt spinosae istae quaestiones in certamina ebulliant. Caeterum λογομαχίας, nominat contentiosas disputationes de verbis magis quam de rebus: vel, vt vulgo loquuntur, sine materia, aut subiecto. Nam si quis diligenter expendat, quaecunque intet sophistas feruent contentiones, deprehendet non ex rebus eas prouenire:sed ex nihilo eas excudi. In summa: quaestiones omnes damnare voluit Paulus, quae nos ad nihili contentiones acuunt.  \pend\pstart Ex quibus oritur inuidia.) Ab effectis demonstrat, quantopere fugienda sit ambitiosa sciendi cupiditas. Mater enim inuidiae est ambitio. Vbi autem regnat inuidia, illic simul feruent conuitia, rixae, et reliqua mala, quae hic a Paulo numerantur. Addit, talibus peruersam esse mentem, et esse veritate priuatos. Porro certum est, hic taxari sophistas qui, posthabita aedificatione, verbum Dei in artem subtiliter disserendi et friuolas argutias conuertunt. Si tantùm ostenderet Apostolus, inutilem reddi hoc modo salutis doctrinam: iam ista esset minime tolerabilis profanatio. Sed longe grauior ac vehementior est ista reprehensio, quod tam pernitiosa mala, quae inde manant, et tam noxias pestes recenset. Ex hoc itaque loco discamus sophisticen detestari, vt rem Ecclesiae Dei magis pestiferam quani satis credi possit.  \pend\pstart Quaestum esse pietatem.) Sic resolue. Pietatem esse quae\pend
\section*{I. TIMOTH. }
\marginpar{[ p.76 ]}\pstart stum, vel artem quaestuosam: quia scilicet totum Christianismum luero metiuntur. Perinde enim ac si non alium in finem prodita fuissent Spiritussancti oracula, quam vt eorum auaritiae seruirent, tanquam merce venum exposita, negociantur. Cum talibus vetat Paulus seruos Christi quidquam habere commercii. Nam et si palam non aduersentur Euangelio, quin potius falsa eius professione se iactent: summa tamen danda est a nobis opera, vt omnes intelligant adeo nos esse ab illis alienos, vt nihil nobiscum habeant commune.  \pend
\phantomsection
\addcontentsline{toc}{subsection}{\textit{Est autem quaestus magnus pietas, cum sufficientia. Nih il enim intulimus in mundum: certum quod neque efferre quidquam possumus. Habentes autem alimenta et tegmina, his contenti erimus. Nam qui volunt ditescere, incidunt in tentationem et laqueum, et cupiditates multas stultas et noxias, quae demergunt homines in exitium et inte ritum. Radix enim omnium malorum est auaritia, cui addicti quidam aberrarunt a fide, et seipsos implicuerunt doloribus multis. }}
\subsection*{\textit{Est autem quaestus magnus pietas, cum sufficientia. Nih il enim intulimus in mundum: certum quod neque efferre quidquam possumus. Habentes autem alimenta et tegmina, his contenti erimus. Nam qui volunt ditescere, incidunt in tentationem et laqueum, et cupiditates multas stultas et noxias, quae demergunt homines in exitium et inte ritum. Radix enim omnium malorum est auaritia, cui addicti quidam aberrarunt a fide, et seipsos implicuerunt doloribus multis. }}\pstart Est autem quaestus.) Eleganter non sine aronica correctione in contrarium sensum eadem verba mox retorquet: acsi dixisset: Perperam illi et nequiter, qui venalem habent Christi doctrinam: quasi vero pietas esset quaestus. Quanquam certe, si quis recte aestimet, pietas magnus est ac vberrimus quaestus. Ideo autem sic vocat, quod plenam et absolutam beatitudinem nobis affert. Sacrilegi ergo illi, qui lucro nummario addicti, pietatem conferunt ad suum quaestum. Nobis autem pietas satis magnum per se lucrum est, quia per eam consequimur, non tantùm vt simus haeredes mundi, sed etiam vt fruamur Christo omnibusque eius diuitiis.  \pend\pstart Cum sufficientia.) Tam ad affectum animi, quam ad rem ipsam referri potest. Si ad animum referas, sensus erit: pios, dum nihil appetunt, sed contenti sunt sua tenuitate, satis magnum luerum esse consecutos. Si rerum sufficientiam accipias (quod mihi non mi nus arridet) erit promissio, qualis illa Psalmi 3 4. Leones esuriunt, et famelici vagantur: qui autem Dominum inquirunt, non destituentur omni bono. Semper enim adest suis Dominus, et quantùm illorum necessĩtati sufficit, quidquid est bonorum, erogat. Ita vera foelicitas in pietate sita est: haec vero sufficientia veluti quoddam est auctarium.  \pend
\section*{CAPVT VI. }
\marginpar{[ p.77 ]}\pstart Nihil enim intulimus. Hoc ad finiendum sufficientiae mo dum addidit. Gurges est inexplebilis, nostra cupiditas, nisi coerceatur. Hoc autem optimum frenum, si nihil plus appetamus, quam postulat vitae huius necessitas. Nam ideo fines transilimus, quia ad mille vitas, quas falsa imaginatione concipimus, solicitudo nostra se extendit. Nihil est magis vulgare, nihil etiam magis confessum, hac Pauli sententia. Verum, simul atque omnes confessi fuerint, quod etiam oculis quotidie cernitur: vnusquisque votis immensa latifundia non secus absorbet, quam si aluum haberet dimidii mundi capacem. Atque hoc est quod habetur Psasmo 49. Cum stultitia patrum appareat, in eo quod sperabant hic se perpetuo mansuros: nihilominus posteros illorum viam approbare. Ergo, vt constet nobis sufficientia: discamus animum ita compositum habere, vt nobis sufficiat, quod ad tolerandam vitam necessarium est. Cum alimen ta et vestiarium nominat, delitias et abundantem copiam excludit. Natura enim paruo contenta est. superfluum autem, quidquid exuperat naturalem vsum. Non quod liberalior vsus per se damnari de beat. Sed appetitus semper vitiosus est.  \pend\pstart Qui volunt ditescere.) Postquam ad continentiam hortatus est, ac diuitiarum contemptum, nunc admonet, quam pernitiosa sit habendi cupiditas: ac praesertim in Ecclesiae ministris, de quibus nominatim hic tractat. Non sunt autem diuitiae malorum, quae hic commemorat, causa: sed appetitus, etiam si quis alias sit pauper. Nec tantùm quid plaerunque accidere soleat, sed quid perpetuo accidere necesse sit, docet. Quisquis enim hoc habet propositum, vt ditescat, Diabolo se captiuum tradit. Verissimum est illud ethnici poetae: Qui diues fieri vult: et cito vult fieri. Vnde sequitur, omnes qui feruntur ad diuitias expetendas, praecipites ruere. Hine etiam stultae illae, imo insanae cupiditates, quae tandem eos in exitium demergunt. Hoc quidem malum est vniuersale: sed in Ecclesiae pastoribus magis conspicuum eminet. Ita enim eos dementat auaritia, vt nihil quantumuis absurdum refugiant, simul atque auri fulgor vel argenti oculos perstrinxit.  \pend\pstart Radix enim malorum.) Non est, cur in aliorum vitiorum comparatione nimis scrupulose laboremus. Certum est, ambitionem ac superbiam saepe deteriores proferre fructus. Atqui non prouenit ex auaritia ambitio. De libidine venerea idem potest dici. Sed Pauli consilium non fuit, quaecunque nominari possint vitiorum species, eas sub auaritia includere. Quid igitur? Simpliciter intellexit, infinita mala inde oriri. Quemadmodum vulgo loquimur: nihil esse malorum, quod non pariat, vel discordia, vel gula, vel ebrie tas, vel aliud quoduis simile vitium. Et profecto de scelerata haben  \pend
\section*{I. TIMOTH. }
\marginpar{[ p.78 ]}\pstart di cupidine, verissime hoc praedicare licet: nullam esse malorum speciem, quae non ab ea quotidie scaturiat. Vt sunt innumerae fraudes, fallaciae, periuria, circunuentiones, rapinae, crudelitas, corruptelae in iudiciis, lites, odia, veneficia, caedes. Denique omnes fere nocendi artes. Similes passim sententiae apud profanos scriptores occurrunt. Quare improbe faciunt quidani hyperbolice hoc a Paulo dictum insimulantes, qui Horatio, vel Ouidio applauderent sic loquenti. Vtinam vero non comprobaret quotidiana experientia.rem ipsam, vt est, simpliciter fuisse indicatam. Sed meminerimus, eadem scelera, quae ab auaritia manant, posse etiam vel ab ambitione, vel ab inuidia vel ab aliis prauis affectibus oriri:vt certe oriuntur.  \pend\pstart Cui addicti quidam.) Graeca locutio impropria est, cum dieit, appeti auaritiam: sed quae sensum nihil obscurat. Malum enim omnium atrocissimum ab auaritia nasci pronunciat, defectionem a fide. Nam qui laborant hoc morbo, paulatim degenerant, donec abiiciant penitus fidem. Hinc dolores illi, quos commemorat: quo nomine intelligo dira conscientiae tormenta, quae solent desperatis hominibus accidere.  \pend
\phantomsection
\addcontentsline{toc}{subsection}{\textit{Tu vero, o homo Dei, haec fuge. Sectare verô iustitiam, pietatem, fidem, charitatein, patientiam mansuetudinem. Certa bonum certamen fidei. Apprehende vitam aeternam, ad quam etiam vocatus es, et confessus bonam confessionem coram multis testibus. Denuncio tibi coram Deo, qui viuificat omnia, et Christo Iesu, qui testificatus est bonam confessionem coram Pontio Pilato, vt serues mandatum immaculatus et irreprehensibilis, vsque ad reuelationem Domini nostri Iesu Christi: quam suis temporibus manifestabit beatus et solus princeps, Rex regnantium, et Dominus dominantium, qui solus habet immortalitatem, qui lumen habitat inaccessum, quem vidit nullus hominum, nec videre potest, cui honor et potentia aeterna, Amen. }}
\subsection*{\textit{Tu vero, o homo Dei, haec fuge. Sectare verô iustitiam, pietatem, fidem, charitatein, patientiam mansuetudinem. Certa bonum certamen fidei. Apprehende vitam aeternam, ad quam etiam vocatus es, et confessus bonam confessionem coram multis testibus. Denuncio tibi coram Deo, qui viuificat omnia, et Christo Iesu, qui testificatus est bonam confessionem coram Pontio Pilato, vt serues mandatum immaculatus et irreprehensibilis, vsque ad reuelationem Domini nostri Iesu Christi: quam suis temporibus manifestabit beatus et solus princeps, Rex regnantium, et Dominus dominantium, qui solus habet immortalitatem, qui lumen habitat inaccessum, quem vidit nullus hominum, nec videre potest, cui honor et potentia aeterna, Amen. }}\pstart Tu verd.) Hominem Dei vocando, pondus addit cohortationi. Et remedium opponit corrigendae auaritiae, cum praecipit quales appetere diuitias debeat. Nempe spirituales. Nominat autem  \pend
\section*{CAPVT VI. }
\marginpar{[ p.79 ]}\pstart quasdam virtutum species exempli causa, sub quibus alias etiam intelligere conuenit. Ergo quisquis iustitiae studio tenebitur, et ad pietatem, fidem, charitatem adspirabit, tolerantiam et mansuetudinem sectabitur, fieri non poterit, quin ab auaritia eilisque fructibus abhor reat.  \pend\pstart Certa bonum certamen.) Quemadmodum in proxima Epistola dicet: qui militiae nomen dedit, non implicatur alienis a sua professione negociis. Itaque vt Timotheum retrahat a nimia rerum terrestrium cura, certandum esse admonet. Quo autem eum animet ad obeundum fortiter eiusmodi certamen, bonum appellat: hoc est foelix, et ideo minime fugiendum. Nam si dubio euentu, et cum mortis periculo pugnare non dubitant terreni milites: quanto meliore animo certare nos conuenit sub Christi auspiciis et vexillo, postquam de victoria certi sumus? Praesertim cum nos praemium maneat, non quale militibus suis persoluere solent alii Imperatores: sed gloriosa immortalitas, et coelestis beatitudo. Indignum sane foret, tali spe proposita, nos fatigari vel labascere. Atque id est, quod continuo post subiicit: Apprehende vitam aeternam. Perinde enim est, acsi diceret: Deus ad vitam aecernam te vocat. Igitur, contempto mundo, illuc enitere. Apprehendere autem cum ipsam iubet, vetat in medio cursu cessare aut pigressere. Vocationis meminit, ad mcitandam alacritatem. Nihil est enim, quod plus animi facere nobis debeat, quam dum audimus, nos a Deo vocatos. Nam inde colligimus, non temere nos nec frustra laborare, vbi Deum habemus instituti nostri authorem. Et ista esset grauis exprobratio, Deivocationem repudiasse. Quare hic acerrimus stimulus esse debet: Deus te ad vitam aeternam vocat. Caue ne alid abstraharis, aut alioqui deficias, priusquam eam fueris consecutus.  \pend\pstart Et confessus bonam confessionem.) Commemoratione etiam superioris vitae magis ad pergendum stimulat. Turpius enim est, deficere postquam recte coeperis, quam nunquam coepisse. Ergo hoc argumento acuit Timotheum, quod strenue et cum laude hactenus se gesserit: vt primis vltima respondeant. Confessionem intelligo, non verbis conceptani, sed potius reipsa editam: neque id semel duntaxat, sed in toto ministerio. Sensus ergo est: Habes praeclarae tuae confessionis multos testes cum Ephesi, tum in aliis regionibus, qui te fideliter et serio in Euangelii professione versantem spectarunt. Quamobrem, tali specimine edito, iam tibi non licet, nisi cum summo pudore et dedecore, aliuni te praestare, quam egregium Christi mi litem.  \pend\pstart Denuncio tibi.) Tanta obtestationum vehementia, qua Paulus vtitur, argumento est, quam difficilis sit ac rara virtus, persta\pend
\section*{I. TIMOTH. }
\marginpar{[ p.80 ]}\pstart re in Euangelii ministerio, vt decet, vsque ad extremuni. Tametsi enim alios quoque hortatur sub Timothei persona: simul tamen eum conipellat. Porro quae de Christo et Deo praedicat, ex cireunstantia praesentis causae dependent. Nam cum Deo hoc tribuit, quod viuificet omnia, vult occurrere crucis offendiculo, quae nihil quam mortis speciem praesefert. Oculos ergo claudendos esse significat, vbi mortem ostentant ac minantur impii: vel potius in Deum solum esse defigendos: quia ipse est, qui mortuos restituit in vitam. Summa huc tendit, vt omisso mundi intuitu, in solum Deum respicere discamus. Quod de Christo mox subiicit, continet eximiam con firmationem. Admonemur enim, non sub Platone aliquo nos philosophari qui ociosas disputationes in vmbra tractet: sed doctrinam, quam profitemur, morte filii Dei sancitam esse. Non enim multis verbis Christus suam confessionem sub Pilato, sed reipsa, voluntariam scilicet mortem obeundo, asseruit. Hanc vocat bonam confessionem. Nam Socrates quoque mortuus est. Neque tamen mors eius legitima fuit doctrinae, quam tenebat, comprobatio. At cum audimus fusum esse sanguineni filii Dei:illud sacrosanctum est sigillum, quod omni dubitatione nos eximit. Ergo quoties vacillant animi, meminerimus a Christi morte confirmationem esse subinde petendam. Cuius enim ignauiae foret, talem ducem nobis praeeuntem deserere?  \pend\pstart Vt serues mandatum.). Nomine mandati significat, quae hactenus de officio Timothei disseruit. Quorum haec fuit summa: vt se fidelem Christo et Ecclesiae ministrum praeberet. Quid enim opus est, hoc extendere ad totam legem? Nisi forte simpliciter accipere quis malit pro functione illi iniuncta. Neque enim constituimur Ecclesiae ministri, quin simul nobis Deus praescribat, quidnam agere nos velit. Ita seruare mandatum nihil aliud esset, quam munus sibi commissum bona fide exequi. Omnino certe ad ministerium Timothei refero. Quae sequuntur epitheta, quod ad casum et terminationem spectat, Graece tam mandato, quam personae Timothei conueniunt. Sensus tamen, quem reddidi, multo est aptior.  \pend\pstart Vsque ad reuelationem.). Dici non potest, quam fuerit piis omnibus necessarium illo tempore, mentes habere in diem Christi prorsus defixas: cum in mundo quaquauersum occurrerent innume ra scandala. Vndique impetebantur: omnibus odio erant et execrationi: omnium ludibriis expositi erant: quotidie premebantur nouis aerumnis: nullus interea tot laborum et molestiarum fructus apparebat. Quid igitur erat reliquum, nisi vt cogitatione transuolarent ad beatum illum diem redemptionis nostrae? Quanquam eadem hodie quoque ratio apud nos valet: imo fere omnibus seculis est communis. Quam multa enim assidue oculis nostris ingerit Sathan, quae a\pend
\section*{CAPVT VI. }
\marginpar{[ p.81 ]}\pstart lioqui nos a recto proposito centies abducerent? Omitto ignes, gladios, exilia, et omnes furiosos hostiunr impetus. Omitto calumnias, et alias vexationes. Quam multa sunt intus longe peiora? Insurgunt ambitiosi homines: subsannant Epicurei et Lucianici:insultant proterui:fremunt hypocritae qui secundùm carnem sapiunt oblique nos mordent: solicitamur variis artibus huc et illuc. Denique, ingens est miraculum, cum in tam arduo et periculoso munere quispiam infractus perseuerat. His tot difficultatibus vnicum est remedium, oculos coniicere in Christi apparitionem et illinc semper pendêre. Quia autempraecipites in votis nostris solemus ferri: et parum abest, quin diem et horam praestituamus Deo, ne praestare differat si quid promisit: ideo nimiam festinationem mature cohibet, in expectando Christi aduentu. Nam hoc sibi volunt haec verba: quam Deus suis temporibus repraesentabit. AEquiore enim animo expectant homines, cuius nondum opportunum adesse tempus sciunt. Qui fit, vt in sustinendo naturae ordine simus tam patientes, nisi quod nos retinet ista cogitatio: intempestiue nos facturos, si contra desideriis nostris luctemur? Ita sciamus, Christi reuelationi suum esse tempus destinatum, quod patienter expectare nos oporteat.  \pend\pstart Beatus et solus princeps.) Istis elogiis ideo extollit Dei principatum;ne principum huius seculi fulgor oculos nobis perstrin gat. Ac praecipue illo tempore necessaria fuit eiusmodi admonitio. Quanta enim tunc fuit regnorum omnium magnitudo et potentia, tantundem Dei maiestas et gloria obscurabantur. Nam quotquot rerum potiebantur, non modo infestissimi erant regni Dei hostes: sed superbe insultabant Deo, sacrumque eius nomen pedibus calcabant. Ac quo maiore fastu despiciebant veram religionem, eô beatiores sibi videbantur. Quis non ex tali rerum facie iudicasset, Deum misere victum et oppressum? Videmus quam insolenter efferatur Cicero aduersus Iudaeos, propter afflictas eorum fortunas, in oratione pro Flacco. Dum ita rebus secundis inflatos esse impios vident boni, interdum exanimantur. Quare Paulus, vt ab euanido isto splendore piorum oculos auertat, beatitudinem, principatum et regnum vindicat vni Deo. Solum principem cum vocat, non euertit politicum ordinem: quasi nulli in mundo esse debeant magistratus aut reges: sed eum esse solum significat, qui a seipso regnat, ac sunpte virtute. Quod etiam ex sequentibus patet, quae vice expositionis addidit: Rex regnantium. Dominus dominantium. Summa est, omnes mundi potestates summo eius imperio subesse, ab ipso pendêre, stare denique et cadere eius nutu. Hoc vero imperium esse incomparabile: quia ad eius gloriam nihil sint reliqua omnia:et cum haec caduca sint et breui interituraillud perpetuo duret.  \pend
\section*{l. TIMOTH. }
\marginpar{[ p.82 ]}\pstart Qui solus habet immortalitatem.) In hoc incumbit Paulus, vt nihil extra Deum foelicitatis, nihil dignitatis aut excellentiae, nihil vitae reliquum faciat. Ideo nunc dicit solum Deum esse immortalem, vt sciamus, nos cum omnibus creaturis non proprie viuere, sed vitam ab ipso mutuari duntaxat. Hinc sequitur, dum adspiramus ad Deum fontem vitae immortalis, hanc praesentem vitam pro nihilo ducendam esse. Sed obiicitur, suam humanae animae et An gelis esse immortalitatem: ideo hoc non vere de solo Deo praedicari. Respondeo, non hic negari, quin Deus in quas velit creaturas immor talitatem diffundat, cum eam solus habere dicitur. Perinde enim hoc valet, acsi dixisset Paulus, solum Deum non a seipso tantum esse immortalem, et suapte natura: sed immortalitatem in potestate habere. vt in creaturas non competat, nisi quatenus suam illis virtutem inspirans, eas vegetat. Nam si vim Dei quae indita est hominis animae, tollas: statim euanescet. Ideni et de Angelis sentiendum. Ergo, proprie loquendo, non subest animarum, vel Angelorum naturae immortalitas, sed aliunde manat:nempe ab arcana Dei inspiratione. Iuxta illud: In ipso viuimus et moueniur, et sumus. Si quis plura et subciliora desideret: legat duodecimum Augustini librum De ciuitate Dei.  \pend\pstart Lumen habitat.) Duo significat: Deum nobis esse absconditum: nec tamen causam obscuritatis in ipso esse, quasi in tenebris lateat, sed in nobis, quibus propter imbecillam mentis aciem, vel hebetudinem potius, non patet accessus ad eius lumen. Caeterum inaccessum Dei lumen, intellige, si quis proprio marte ad ipsum contendat. Nisi enim sua gratia patefaceret nobis Deus, aditum non diceret Propheta: Qui ad eum accedunt illuminantur. Quanquam verum est, nunquam, dum hac mortali carne circundamur, nos eovsque penetrare in vlcima adyta Dei,vt nihil nos fugiat. Ex parte enim cognoscimus: et videmus quasi per speculum et in aenigmate. Ingredimur ergo per fidem in lucem Dei: sed ex parte. Interea verum est illud: esse lucem illam homini inaccessam.  \pend\pstart Quem nullus hominum vidit.) Maioris explicationis causa hoc additum est: vt discant homines fide in eum respicere, cuius speciem oculis non cernunt:adeóque nec mentis intellectu. Neque enim hoc tantùm ad corporeos oculos refero: sed ad ipsam quoque animae facultatemi. Videndum semper quorsum tendat Apostolus. Difficile est, omnia, quorum praesenti eonspectu fruimur, posthabere et negligere, vt ad Deum enitamur, qui nusquam apparet. Nam subinde obrepit haec cogitatio: Qui seis an Deus sit, cum tantùm audias ipsum esse, non autem videas? Aduersus hoc periculum nos praemunit Apostolus, cum aestimandum id esse negat ex sensu nostro, quia id captum nostrum superet. Ideo enim non videmus, quia non est a\pend
\section*{CAPVT VI. }
\marginpar{[ p.83 ]}\pstart cies, quae in tantam altitudinem conscendat. Longa est disputatio apud Augustinum de hac sententia: quia videtur pugnare cum eo quod habetur in prima Epistola Ioannis: Tunc videbimus eum sicuti est, quia similes ei erimus. Eam disputationem cum agitet pluribus locis, nullibi melius videtur expedire, quam in Epistola ad Paulinam viduam. Quod tamen ad sensum praesentis loci attinet: simplex est responsio: nos videre Deum non posse in hac natura, quemadmodum alibidicitur: Caro et sanguis Regnum Dei non possidebunt. Renouari enim nos oportet, vt simus Deo similes, priusquam eum vi dere nobis detur. Et ne simus vltra modum curiosi, semper illud suceurrat, plus in hac inquisitione valere viuendi, quam loquendi modum. Simul et illud cauere meminerimus, quod prudenter Augustinus cauendum praecipit:ne dum altercando vestigamus, quomodo possit Deus videri. ipsam pacem, sanctificationemque perdamus, sine qua nemo Deum videbit vnquam.  \pend
\phantomsection
\addcontentsline{toc}{subsection}{\textit{His qui diuites sunt in hoc seculo praecipe, ne efferantur, neve sperent in diuitiarum incertitudine, sed in Deo viuo, qui abunde suppeditat omnia ad fruendum, vt benefaciant, vt diuites sint in operibus bonis, faciles ad largiendum, libenter communicantes, recondentes sibiipsis fundamentum bonum in posterum, vt vitam, quae vere aeterna est, apprehendant.O Timothee depositum custodi, deuitans profana vaniloquia, et oppositiones falso nominatae scientiae: quam quidam profitentes aberrarunt a fide. Gratia tecum, Amen. }}
\subsection*{\textit{His qui diuites sunt in hoc seculo praecipe, ne efferantur, neve sperent in diuitiarum incertitudine, sed in Deo viuo, qui abunde suppeditat omnia ad fruendum, vt benefaciant, vt diuites sint in operibus bonis, faciles ad largiendum, libenter communicantes, recondentes sibiipsis fundamentum bonum in posterum, vt vitam, quae vere aeterna est, apprehendant.O Timothee depositum custodi, deuitans profana vaniloquia, et oppositiones falso nominatae scientiae: quam quidam profitentes aberrarunt a fide. Gratia tecum, Amen. }}\pstart Verisimile est, cum multi inter Christianos tenues effent et abiecti, a diuitibus, vt fieri solet, fuisse contemptos. Idque Ephesipotissimùm, in vrbe opulenta, vsu venire potuit. Illic enim, vt plurimum magis regnat fastus. Atque hinc colligimus, quam periculosa res sit bonorum affluentia. Neque enim frustra Paulus tam seuera admonitione peculiariter diuites compellat: sed vitiis medetur, quae perpetuo fere sequuntur diuitias, non secus atque vmbra corpus: idque ingenii nostri prauitate, qui ex Dei donis peccandi materiam semper arripimus. Duo autem nominatim ponit, quae debeant cauere diuites, superbiam, et spem fallacem: quorum prius ex posteriore nascitur: et ideo videtur secundo loco additum, neve sperent, et c.vt  \pend
\section*{I. TIMOTH. }
\marginpar{[ p.84 ]}\pstart fontem omnis superbiae notaret Paulus. Vnde enim fit, vt insolescant diuites, et alios despiciendo sibi nimium placeant, nisi quia beatos se esse imaginantur? Praecedit igitur inanis fiducia, deinde elatio sequitur. Haec vitia dum vult Paulus corrigere, primùm contenptim de diuiriis loquitur. Nam particula, in hoc seculo, ad extenuationem tendit. Quidquid enim ex seculo est, eius naturam resipit: vt seilicet fluxum sit, ac cito praetereat. Spei quae in diuitüs locatur, fallaciam et vanitatem ex eo ostendit, quod earum possessio tam est euanida, vt similis sit rei incognitae. Nam dum eas tenere nos putamus, effluunt momento. Quam igitur stultum est, spem in illis figere?  \pend\pstart Sed in Deo viuo.) Qui hoc tenebit, facile spem a diuitiis suam abstrahet. Nam si solus Deus est, qui nobis ad necessarios vitae vsus omnia suppeditat: tranffertur eius officium ad diuitias, si spes in illis locatur. Obserua autem esse tacitam antithesin, cum praedicat Deum omnibus affatim dare. Sensus enim est: Etiam si plena rerum omnium copia affluamus, nos tamen nihil habere, nisi ex sola Dei benedictione: quoniam ea sola est, quae nobis subministrat quidquid opus est. Vnde sequitur valde falli, qui in diuitias recumbunt: et non toti pendent a Dei benedictione, in qua nobis sira est victus et rerum omnium sufficientia. Hinc etiam colligimus, non ideo tantùm a spe diuitiarum nos reuocari, quia tantùm ad vsum mortalis vitae pertineant sed etiam quia nihil sint nisi fumus. Neque enim pane alimur tantùm, sed Dei benedictione. Cum dicit abunde, aς ἀπόλαυσιν, significat quam liberalis sit erga nos Deus, imo erga vniuersos mortales, adeôque bruta animalia. Nam vltra necessitatem longe lateque se diffundit eius beneficentia.  \pend\pstart Vt benefaciant.) Superioribus remediis aliud etiam adiungit, ad prauos diuitum affectus corrigendos, dum praescribit quis legitimus sit bonorum vsus. Nam qud quisque plus opibus abundat. eo amplior illi beneficentiae materia suppetit. Et quia semper ad erogandum pigriores sumus quam deceats pluribus verbis eam virtutem commendat. Addit praeterea stimulum ex mercedis promissione. Largiendo, inquit, et communicando meliorem sibi thesaurum in coe lo comparent, quam habere possint in terra. Fundamenti nomine. stabilem et firmam durationem significat. Neque enim tineis aut furibus, aut incendiis obnoxiae sunt diuitiae spirituales, quas nobis in coelo recondimus: sed perpetuo manent, extra periculum constitutae. Nihil enim in terra vere fundatum est, sed omnia quodammodo fluctuant. Porro quod inde inferunt Papistae, nos igitur bonis operibus vitam aeternam promereri:nimis friuolum est. Acceptum certe fert Deus, quidquid pauperibus impensum fuerit.Sed quia vix centesima  \pend
\section*{CAPVT VI. }
\marginpar{[ p.85 ]}\pstart ex parte officio suo satiffaciunt etiam perfectissimi, indigna est nostra benignitas, quae in rationem coram Deo veniat. Imo si ad calculum nos vocaret Deus, nemo erit qui non conturbet. Adeo longe absumus a plena solutione. Sed postquam nos gratuito sibi reconciliauit Deus, officia nostra qualiacunque sint, sic accepta habet, vt indebita mercede dignetur. Proinde remuneratio non ex meriti ratione pendet, sed ex liberali Dei acceptione: adeôque fidei iustitiae non repugnat, vt cius sit quasi appendix.  \pend\pstart O Timothee.) Tametsi varie exponunt depositi nomen, ego simpliciter eo notari puto gratiam, qua instructus erat Timotheus ad munus suum exequendum. Eadem vero ratione depositum no minatur, qua et talentum. Nam ea lege nobis committit, quidquid donotum in nos confert Deus, vt reddenda sit olim ratio, si nostra ignauia perierit vtilitas, quae inde prouenire debuerat. Admonet igitur, vt diligenter conseruet quod sibi datum est, vel potius apud se depositum, ne patiatur corrumpi vel adulterari, vel ne seipsum priuet ac spoliet propria culpa. Saepe enim ingratitudo nostra, aut donorum Dei abusus facit, vt ipsa nobis auferantur. Proinde Timotheum admonet Paulus, vt bona conscientia rectóque vsu, quod sibi commis sum est, conseruare studeat.  \pend\pstart Deuitans profanas.) Huc spectat admonitio, vt intentus sit in solidam doctrinam: quod fieri non potest, quin ab ostentatione abhorreat. Nam vbi regnat placendi ambitio, nullum amplius viget aedificationis studium. Proinde cum mentionem faceret custodiendi depositi, non abs re hanc cautionem mox adhibuit, de fugienda profana garrulitate. Porro quod κινοβωνί αι vertunt inanitates voeum, non adeo displicet, nisi quod postea ambiguitate decepti, perperam exponunt. Voces enimaccipiunt pro vocabulis, vt est fatum, vel fortuna. Ego autem grandiloquentiam eorum et inflatam magnis ampullis garrulitatem notari existimo, qui Euangelii simplicitate non contenti, ipsum in profanam philosophiamconuertunt. Ergo κιν οευνίαι, sunt non in singulis vocibus, sed in perpetua resonantia, quam despumant ambitiosi homines, qui plausum magis captant, quam Ecclesiae profectum. Et aptissime expressit Paulus rem ipsam. Nam eum grande nescio quid resonent, nihil tamen subest. Est igitur inanis tinnitus: quem profanum etiam vocat, quoniam extinda est vis spiritus, simul atque Doctores itae inflant suas tibias, ad suam eloquentiam venditandam. Contra tam claram ac testatam Spiritussancti prohibitionem perrupit nihilominus haec pestis: ac statim quidem ab inirio coepit pullulare, ita vero demum grassando obtinuit in Papatu, vt adulterina theologiae larua, quae illic regnat, viuum sit speculum profanae istius et inanis resonantiae, cuius meminie Pau\pend
\section*{I. TIMOTH. }
\marginpar{[ p.86 ]}\pstart lus. Omitto, quod innumeris erroribus, deliriis, ac blasphemiis scaetent eorum libri ac clamosae illae disputationes. Verùm etiam si aliat nihil docerent impium: quoniam tamen tota eorum doctrina nihil continet, praeter sesquipedales ampullas: a maiestate Scripturae, ab energia spiritus, a prophetica grauitate, a sinceritate Apostolica prorsus aliena est: eo ipso mera est genuinae theologiae profanatio. Quid enim, obsecro, de fide, aut poenitentia, aut inuocatione Deit quid de hominum imbecillitate, Spiritussancti auxilio, gratuita peccatorum remissionet quid de officio Christi tractant, quod ad solidam pietatis aedificationem valeat? Sed de his adhuc dicendum erit in altera Epistola. Certe quisquis mediocri intelligentia et ingenuitate erit praeditus, fatebitur totam magniloquentiam Papalis theologiae, et omnes quae in scholis tonant magistrales determinationes, nihil esse aliud, quam profanas κενοφονίας: nec alio titulo posse aptius designari. Et certe haec iustissima est humanae arrogantiae poena, vt quicunque degenerant a Scripturae puritate, profanescant. Quare non possunt nimis sedulo intenti esse Ecclesiae Doctores ad cauendas eius modi corruptelas, et arcendam inde iuuentutem. Vetus interpres κανοφυνί ας legens eum diphrongo, transtulit nouitates. Et apparet ex veterum commentariis, lectionem hanc olim fuisse receptam inter multos, quam nunc etiam Graeci quidam codices retinent. Verùm prior illa, quam secutus sum, longe melius quadrat.  \pend\pstart Et oppositiones.) Hoc quoque admodum proprie et eleganter. Tantùm enim tumoris habent argutiae, quibus se iactant homines gloriae cupidi, vt germanam Euangelii doctrinam obruant. Pompam ergo illam, quae eminet, et quae cum plausu excipitur a mundo, vocat Apostolus oppositiones. Ambitio quidem semper contentiosa est, ac certaminum mater. Vnde fit, vt qui se iactare appetunt, quacunque de re digladiari sint parati. Huc tamen praecipue respexit Paulus, quod inanis sophistarum doctrina altius in aereas speculationes et argutias exurgens, sua specie non obscurat modo simplicitatem verae doctrinae, sed etiam opprimit ac reddit despicabilem, vt mundus fere externo splendore ducitur. Non vult Paulus Timotheum propterea aemulatione pungi, vt simile aliquid tentet. Cuni autem humanae curiositati magis adlubescant quae subtilitatis speciem praeseferunt, vel sunt ad ostentationem composita: Paulus ex aduerso pronunciat, falso nominari ac putari scientiam, quae se extollit aduersus simplicem et humilem pietatis doctrinam. Quod diligenter obseruandum est vt discamus secure ridere, ac contemnere to tam istam fucatam sapientiam, cuius admiratione mundus obstupescit, in qua nulla subest aedificatio. Ea enim demum, authore Paulo, vere et merito nominatur scientia, quae in fiducia timoreque Dei,  \pend
\marginpar{[ p.87 ]}CAPVT VI. \pstart toc est in pietate nos instituit.  \pend\pstart Quam quidam profitentes.) Ab exitu etiam demonstrat. quam noxiosa res sit, et quantopere fugienda. Deus enim ita eorum arrogantiam vleiscitur, qui comparandae famae studio, pietatis doctrinam corrumpunc ac deformant, quod eos excidere patitur ac sani tate mentis, vt multis absurdis erroribus se inuoluant. Quod etiam videmus accidisse in Papatu. Nam postquam coeperunt profano more de mysteriis religionis nostrae philosophari, secuta sunt innumera fal sarum opinionum portenta. Fides hic, vt superioribus aliquot locis, pro summa religionis et sana doctrina capitur. Eiusmodi exemplis moniti, si defectionem horremus a fide, maneamus in puro Dei verbo, sophisticem vero et omnes argutias, cum nefandae sint pietatis corruptelae. detestemur.  \pend
\endnumbering
\end{pages}
\end{document}
        
